% Chapter 4, Section 4 _Linear Algebra_ Jim Hefferon
%  http://joshua.smcvt.edu/linearalgebra
%  2001-Jun-12
\section{Bentuk Jordan}
\index{bentuk Jordan|(}
\noindent\textit{Bagian ini menggunakan materi dari tiga subbagian opsional:
Menggabungkan Subruang, Determinan Ada, dan Ekspansi Laplace.}

Kita mengawali bab ini
dengan mengingat bahwa setiap pemetaan linear $\map{h}{V}{W}$ dapat
direpresentasikan terhadap suatu
basis $B\subset V$ dan $D\subset W$ oleh matriks identitas parsial.
Dengan kata lain, bentuk identitas parsial merupakan bentuk kanonik untuk
ekuivalensi matriks.
Bab ini membahas kasus ketika
kodomain sama dengan domain, sehingga
wajar bila kita bertanya apa yang mungkin dilakukan ketika
kedua basisnya sama, yakni ketika kita mempunyai $\rep{t}{B,B}$.
Singkatnya, kita menginginkan bentuk kanonik untuk keserupaan matriks.

Kita telah mencatat bahwa dalam kasus $B,B$,
matriks identitas parsial tidak selalu dapat diperoleh.
Karena itu, kita memperluas bentuk matriks
yang diminati ke generalisasi alaminya,
yaitu matriks diagonal, dan
menunjukkan bahwa suatu transformasi atau matriks persegi dapat didiagonalkan
jika nilai eigennya berbeda-beda.
Namun, kita juga memberikan contoh
matriks persegi yang tidak dapat didiagonalkan
karena nilpoten,
sehingga bentuk diagonal tidak memadai sebagai bentuk kanonik
untuk keserupaan matriks.

Bagian sebelumnya mengembangkan contoh itu
untuk memperoleh bentuk kanonik bagi matriks nilpoten, yaitu bentuk subdiagonal.

Bagian ini menuntaskan program kita dengan
menunjukkan bahwa untuk setiap transformasi linear terdapat
basis~\( B \)
sedemikian sehingga representasi matriks $\rep{t}{B,B}$ merupakan jumlah
sebuah matriks diagonal dan sebuah matriks nilpoten.
Inilah bentuk kanonik Jordan.









\subsectionoptional{Polinomial Pemetaan dan Matriks}\index{polinomial!pemetaan dan matriks}
Ingatlah bahwa himpunan matriks persegi~\( \matspace_{\nbyn{n}} \)
merupakan ruang vektor terhadap penjumlahan entri demi entri dan perkalian skalar,
dan bahwa
%the unit matrices\Dash all entries are zero except
%for a single entry, which is one\Dash form a basis, so 
ruang ini berdimensi \( n^2 \).
Jadi, untuk setiap matriks \( \nbyn{n} \) $T$, himpunan beranggotakan
\( n^2+1 \), yaitu \( \set{I,T,T^2,\dots,T^{n^2} } \), bergantung
linear, sehingga terdapat skalar \( c_0,\dots,c_{n^2} \),
yang tidak semuanya nol, sedemikian sehingga
\begin{equation*}
  c_{n^2}T^{n^2}+\dots+c_1T+c_0I
\end{equation*}
merupakan matriks nol.
Dengan demikian, setiap transformasi memiliki semacam
nilpotensi tergeneralisasi\Dash pangkat-pangkat
sebuah matriks persegi tidak dapat terus meningkat tanpa suatu bentuk pengulangan.

\begin{definition}
Misalkan \( t \) adalah transformasi linear pada ruang vektor~\( V \).
Jika \( f(x)=c_nx^n+\dots+c_1x+c_0 \) adalah sebuah polinomial,
maka \( f(t) \) adalah
transformasi \( c_nt^n+\dots+c_1t+c_0(\identity) \) pada~\( V \).
Dengan cara yang sama, jika \( T \) adalah matriks persegi,
maka
\( f(T) \) adalah matriks \( c_nT^n+\dots+c_1T+c_0I \).
\end{definition}

\noindent Polinomial matriks merepresentasikan polinomial pemetaan:~jika
\( T=\rep{t}{B,B} \), maka \( f(T)=\rep{f(t)}{B,B} \).
Ini karena \( T^j=\rep{t^j}{B,B} \),
\( cT=\rep{ct}{B,B} \), dan \( T_1+T_2 =\rep{t_1+t_2}{B,B} \).

\begin{remark}
Kita akan menuliskan polinomial matriks sedikit berbeda dari
polinomial pemetaan.
Sebagai contoh, jika \( f(x)=x-3 \), maka
kita akan menuliskan matriks identitas, seperti pada~\( f(T)=T-3I \),
tetapi tidak menuliskan pemetaan identitas, seperti pada~\( f(t)=t-3 \).
\end{remark}

\begin{example}  \label{ex:PolySendRotMatToZ}
Rotasi vektor bidang sebesar \( \pi/3 \)~radian berlawanan arah jarum jam
direpresentasikan terhadap basis standar oleh
\begin{equation*}
  T=
  \begin{mat}[r]
     \cos(\pi/3)  &-\sin(\pi/3)  \\
     \sin(\pi/3)   &\cos(\pi/3)
  \end{mat}
  =
  \begin{mat}[r]
     1/2        &-\sqrt{3}/2  \\
     \sqrt{3}/2  &1/2
  \end{mat}
\end{equation*}
dan verifikasi bahwa \( T^2-T+I=Z \) bersifat rutin.
(Secara geometris, \( T^2 \) merotasi sebesar \( 2\pi/3\).
Pada vektor basis standar~\( \vec{e}_1 \), tindakan
$T^2-T$ menghasilkan selisih antara vektor satuan bersudut
\( 2\pi/3\) dan vektor satuan bersudut \( \pi/3\),
yaitu \( \binom{-1}{0} \).
Pada \( \vec{e}_2 \), hasilnya adalah \( \binom{0}{-1} \).
Jadi, \( T^2-T=-I \).)
\end{example}

Ruang $\matspace_{\nbyn{2}}$ berdimensi empat, sehingga kita mengetahui bahwa
untuk setiap matriks \( \nbyn{2} \)~\( T \), terdapat polinomial derajat empat
\( f \) sedemikian sehingga \( f(T)=Z \).
Namun, dalam contoh itu kita memperlihatkan polinomial derajat dua yang berhasil.
Jadi, meskipun derajat~$n^2$ selalu memadai, dalam beberapa kasus
polinomial berderajat lebih kecil sudah cukup.

\begin{definition}
\definend{polinomial minimal}\index{polinomial!minimal}%
\index{polinomial minimal}
\( m(x) \) dari transformasi \( t \)\index{transformasi!polinomial minimal}
atau matriks persegi \( T \)\index{matriks!polinomial minimal} adalah
polinomial tak nol berderajat terkecil dengan koefisien utama~\( 1 \)
sedemikian sehingga \( m(t) \) merupakan pemetaan nol atau \( m(T) \) merupakan matriks nol.
\end{definition}

\noindent Polinomial minimal
tidak dapat berupa polinomial konstan karena
pembatasan pada koefisien utamanya.
Jadi, polinomial minimal
harus berderajat sekurang-kurangnya satu.
Matriks nol memiliki polinomial minimal $p(x)=x$, sedangkan
matriks identitas memiliki polinomial minimal $\hat{p}(x)=x-1$.

\begin{lemma}
Setiap transformasi atau matriks persegi memiliki polinomial minimal yang unik.
\end{lemma}

\begin{proof}
Pertama-tama kita membuktikan keberadaannya.
Berdasarkan pengamatan sebelumnya bahwa
derajat~$n^2$ memadai, terdapat sekurang-kurangnya satu
polinomial tak nol $p(x)=c_kx^k+\cdots+c_0$ yang
mengirim pemetaan atau matriks tersebut ke nol.
Di antara semua polinomial semacam itu,
ambil salah satu yang berderajat terkecil.
Bagi polinomial ini dengan koefisien utamanya~$c_k$ agar koefisien utamanya menjadi~$1$.
Dengan demikian, setiap pemetaan atau matriks memiliki setidaknya satu polinomial minimal.

Sekarang untuk keunikannya.
Misalkan
\( m(x) \) dan \( \hat{m}(x) \) sama-sama mengirim pemetaan atau matriks tersebut ke nol,
keduanya berderajat
minimal sehingga berderajat sama,
dan keduanya memiliki koefisien utama~$1$.
Perhatikan selisih \( m(x)-\hat{m}(x) \).
Jika selisih itu bukan polinomial nol, maka ia memiliki koefisien utama tak nol.
Membaginya dengan
koefisien utama tersebut akan menghasilkan polinomial yang mengirim
pemetaan atau matriks itu ke nol, memiliki koefisien utama~\( 1\), dan
berderajat lebih kecil daripada $m$ dan $\hat{m}$
(karena dalam pengurangan itu, kedua koefisien utama \( 1\) saling meniadakan).
Hal ini bertentangan dengan minimalitas derajat
$m$ dan $\hat{m}$.
Jadi, \( m(x)-\hat{m}(x) \) adalah polinomial nol dan
keduanya sama.
\end{proof}

\begin{example}  \label{ex:MinPolyForRotMat}
Salah satu cara menghitung polinomial minimal matriks pada
\nearbyexample{ex:PolySendRotMatToZ} adalah mencari pangkat-pangkat $T$
hingga $n^2=4$.
\begin{equation*}
   T^2=
   \begin{mat}[r]
      -1/2         &-\sqrt{3}/2  \\
      \sqrt{3}/2   &-1/2
   \end{mat} 
   \quad
   T^3=
   \begin{mat}[r]
      -1    &0           \\
       0    &-1
   \end{mat}
   \quad
   T^4=
   \begin{mat}[r]
      -1/2         &\sqrt{3}/2  \\
      -\sqrt{3}/2  &-1/2
   \end{mat}
\end{equation*}
Samakan \( c_4T^4+c_3T^3+c_2T^2+c_1T+c_0I \) dengan matriks nol.
\begin{equation*}
  \begin{linsys}{5}
     -(1/2)c_4  
         &-  &c_3           
         &- &(1/2)c_2
         &+ &(1/2)c_1  
         &+  &c_0  &=  &0      \\
      (\sqrt{3}/2)c_4  
         &  &   
         &- &(\sqrt{3}/2)c_2
         &- &(\sqrt{3}/2)c_1  
         &   &     &=  &0        \\
     -(\sqrt{3}/2)c_4  
         &   &    
         &+ &(\sqrt{3}/2)c_2
         &+ &(\sqrt{3}/2)c_1  
        &   &            &=  &0      \\
     -(1/2)c_4  
         &-  &c_3             
         &- &(1/2)c_2
         &+ &(1/2)c_1  
         &+  &c_0  &=  &0
   \end{linsys}
\end{equation*}
Terapkan Metode Gauss.
\begin{equation*}
  \begin{linsys}{5}
     -(1/2)c_4  
         &- &c_3             
         &- &(1/2)c_2
         &+ &(1/2)c_1  
         &+  &c_0  &=  &0      \\
         &  &-(\sqrt{3})c_3 
         &- &\sqrt{3}c_2
         &  &    
         &+ &\sqrt{3}c_0 &=  &0
   \end{linsys} 
\end{equation*}
Dengan tujuan membuat derajat polinomial sekecil mungkin,
perhatikan bahwa menetapkan \( c_4 \), \( c_3 \), dan \( c_2 \) sama dengan nol
memaksa \( c_1 \) dan \( c_0 \) juga sama dengan nol,
sehingga persamaan-persamaan itu tidak memungkinkan polinomial minimal derajat satu.
Sebagai gantinya, tetapkan \( c_4 \) dan \( c_3 \) sama dengan nol.
Sistem
\begin{equation*}
  \begin{linsys}{3}           
         -(1/2)c_2
         &+ &(1/2)c_1  
         &+  &c_0  &=  &0      \\
         -\sqrt{3}c_2
         &  &    
         &+ &\sqrt{3}c_0 &=  &0
   \end{linsys} 
\end{equation*}
memiliki himpunan solusi \( c_1=-c_0 \) dan \( c_2=c_0 \).
Dengan mengambil koefisien utama~\( c_2=1 \), kita memperoleh
polinomial minimal $x^2-x+1$.
\end{example}

% Using the method of that example to find the minimal polynomial of a
% \( \nbyn{3} \) matrix 
% would mean doing Gaussian reduction on
% a system with nine equations in ten unknowns.
Perhitungan itu tidak praktis.
Kita akan mengembangkan cara lain.
% To begin, note that we can break a polynomial of a map or a matrix into
% its components.
% (For this lemma, recall that we are using complex numbers in this chapter
% so all polynomials break completely into linear factors.)

\begin{lemma} \label{le:PolyMapsFactor}
Misalkan polinomial \( f(x)=c_nx^n+\dots+c_1x+c_0 \) difaktorkan sebagai
\( k(x-\lambda_1)^{q_1}\cdots(x-\lambda_z)^{q_z} \).
Jika \( t \) adalah transformasi linear, maka kedua bentuk berikut merupakan pemetaan yang sama.
\begin{equation*}
  c_nt^n+\dots+c_1t+c_0
  =
  k\cdot\composed{\composed{(t-\lambda_1)^{q_1}}{\cdots}}{
      (t-\lambda_z)^{q_z}} 
\end{equation*}
Akibatnya, jika \( T \) adalah matriks persegi, maka \( f(T) \) dan
\( k\cdot(T-\lambda_1I)^{q_1}\cdots(T-\lambda_z I)^{q_z} \)
merupakan matriks yang sama.
\end{lemma}

\begin{proof}
Kita menggunakan induksi pada derajat polinomial.
Kasus ketika polinomial berderajat
nol dan berderajat satu sudah jelas.
Argumen induksi lengkap diberikan dalam \nearbyexercise{exer:PolyMapsFactor},
tetapi gagasannya akan kita tunjukkan melalui kasus derajat~dua.

Sebuah polinomial kuadrat dapat difaktorkan menjadi dua
faktor linear \( f(x)=k(x-\lambda_1)\cdot(x-\lambda_2)
                    =k(x^2+(-\lambda_1-\lambda_2)x+\lambda_1\lambda_2) \).
% (the roots $\lambda_1$ and $\lambda_2$ could be equal).
Menyubstitusikan \( t \)
untuk \( x \) dalam bentuk terfaktor dan
bentuk tak terfaktor menghasilkan pemetaan yang sama.
\begin{align*}
   \bigl(k\cdot\composed{(t-\lambda_1)}{(t-\lambda_2)}\bigr)\>(\vec{v})
   &=\bigl(k\cdot(t-\lambda_1)\bigr)\,(t(\vec{v})-\lambda_2\vec{v})    \\
   &=k\cdot\bigl(t(t(\vec{v}))-t(\lambda_2\vec{v})
      -\lambda_1 t(\vec{v})+\lambda_1\lambda_2\vec{v}\bigr)    \\
   &=k\cdot \bigl(\composed{t}{t}\,(\vec{v})-(\lambda_1+\lambda_2)t(\vec{v})
          +\lambda_1\lambda_2\vec{v}\bigr)                    \\
   &=k\cdot(t^2-(\lambda_1+\lambda_2)t+\lambda_1\lambda_2)\>(\vec{v})
\end{align*}
Kesamaan ketiga menggunakan linearitas \( t \) untuk mengeluarkan $\lambda_2$ dari
suku kedua.
\end{proof}

% In particular, if a minimal polynomial for a transformation $t$ 
% factors as
% $m(x)=(x-\lambda_1)^{q_1}\cdots (x-\lambda_z)^{q_z}$
% then 
% \( m(t)=\composed{\composed{(t-\lambda_1)^{q_1}}{\cdots}}{
%       (t-\lambda_z)^{q_z}} \) 
% is the zero map. 
% Since \( m(t) \) sends every vector to zero, at least
% one of the maps \( t-\lambda_i \)  sends some
% nonzero vectors to zero.
% Exactly the same holds in the matrix case\Dash if $m$ is minimal for $T$ then
% \( m(T)=(T-\lambda_1I)^{q_1}\cdots (T-\lambda_z I)^{q_z} \)
% is the zero matrix and at least one of the matrices $T-\lambda_iI$
% sends some nonzero vectors to zero. 
% That is, in both cases at least some of the \( \lambda_i \) are eigenvalues.
% (\nearbyexercise{exer:SomeRootsMinPolyAreEigs} expands on this.)

Hasil berikutnya menyatakan bahwa
setiap akar polinomial minimal merupakan nilai eigen dan
setiap nilai eigen merupakan akar polinomial minimal.
(Artinya, hasil tersebut menyatakan `$1\leq q_i$',
bukan sekadar `$0\leq q_i$'.)
Untuk itu, ingatlah bahwa untuk mencari nilai eigen,
kita menyelesaikan $\deter{T-xI}=0$, dan
determinan ini menghasilkan polinomial dalam $x$,
yaitu polinomial karakteristik,
yang akar-akarnya merupakan nilai eigen.

\begin{theorem}[Cayley-Hamilton]
\label{th:CayHam}
\index{Teorema Cayley--Hamilton}
\hspace*{0em plus2em}
Jika polinomial karakteristik suatu transformasi atau matriks persegi
difaktorkan menjadi
\begin{equation*}
  k\cdot (x-\lambda_1)^{p_1}(x-\lambda_2)^{p_2}\cdots(x-\lambda_z)^{p_z}
\end{equation*}
maka polinomial minimalnya difaktorkan menjadi
\begin{equation*}
  (x-\lambda_1)^{q_1}(x-\lambda_2)^{q_2}\cdots(x-\lambda_z)^{q_z}
\end{equation*}
dengan \( 1\leq q_i \leq p_i \) untuk setiap \( i \) antara \( 1 \) dan \( z \).
\end{theorem}

\noindent Pembuktiannya menggunakan tiga lema berikutnya.
Kita akan menyatakannya dalam istilah matriks
karena versi itu lebih praktis
untuk pembuktian pertama, tetapi semuanya berlaku sama baiknya
untuk pemetaan.

Hasil pertama merupakan kuncinya.
Untuk pembuktiannya, amati bahwa kita dapat memandang
matriks berentri polinomial sebagai polinomial dengan
koefisien berupa matriks.
\begin{equation*}
   \begin{mat}
     2x^2+3x-1  &x^2+2    \\
     3x^2+4x+1  &4x^2+x+1
   \end{mat}
 = \begin{mat}[r]
    2  &1  \\
    3  &4
  \end{mat}x^2
 + \begin{mat}[r]
    3  &0  \\
    4  &1
  \end{mat}x
 + \begin{mat}[r]
   -1  &2  \\
    1  &1
  \end{mat}
\end{equation*}

\begin{lemma}   \label{le:MatSatItsCharPoly}
Jika \( T \) adalah matriks persegi dengan polinomial karakteristik \( c(x) \),
maka \( c(T) \) adalah matriks nol.
\end{lemma}

\begin{proof}
Misalkan \( C \) adalah \( T-xI \),
yaitu matriks yang determinannya merupakan polinomial karakteristik
\( c(x)=c_nx^n+\dots+c_1x+c_0 \).
\begin{equation*}
  C=\begin{mat}
    t_{1,1}-x        &t_{1,2}   &\ldots        \\
    t_{2,1}          &t_{2,2}-x               \\
    \vdots           &          &\ddots       \\
                     &          &       &t_{n,n}-x
  \end{mat}
\end{equation*}
Ingat Teorema~Empat.III.\ref{th:MatTimesAdjEqDiagDets},
yang menyatakan bahwa hasil kali suatu matriks dengan adjoinnya sama dengan
determinan matriks tersebut dikalikan identitas.
\begin{equation*}
  c(x)\cdot I
  =\adj (C)C
  =\adj (C)(T-xI)
  =\adj (C)T- \adj(C)\cdot x
\tag*{($*$)}
\end{equation*}
Ruas kiri~($*$) adalah
$c_nIx^n+c_{n-1}Ix^{n-1}+\cdots+c_1Ix+c_0I$.
Untuk ruas kanan,
entri-entri \( \adj (C) \) adalah polinomial yang masing-masing berderajat
paling tinggi \( n-1 \), karena minor suatu matriks menghapus satu baris dan satu kolom.
Seperti disarankan sebelum pembuktian, tuliskan kembali sebagai polinomial dengan
koefisien matriks:
\( \adj (C)=C_{n-1}x^{n-1}+\cdots+C_1x+C_0 \)
dengan setiap \( C_i \) merupakan matriks skalar.
Sekarang, berikutlah ruas kanan~($*$).
\begin{equation*}
  [(C_{n-1}T)x^{n-1}+\cdots+(C_1T)x+C_0T]  
   -[C_{n-1}x^n+C_{n-2}x^{n-1}+\cdots+C_0x]
\end{equation*}
Samakan koefisien \( x^n \), $x^{n-1}$, dan seterusnya
pada ruas kiri dan kanan ($*$).
\begin{align*}
  c_nI
  &=-C_{n-1}    \\
  c_{n-1}I
  &=-C_{n-2}+C_{n-1}T    \\
  &\vdotswithin{=}             \\
  c_{1}I
  &=-C_{0}+C_{1}T    \\
  c_{0}I
  &=C_{0}T
\end{align*}
Kalikan dari kanan kedua ruas persamaan pertama dengan \( T^n \),
kedua ruas persamaan kedua dengan \( T^{n-1} \), dan seterusnya.
\begin{align*}
  c_nT^n
  &=-C_{n-1}T^n    \\
  c_{n-1}T^{n-1}
  &=-C_{n-2}T^{n-1}+C_{n-1}T^n    \\
  &\vdotswithin{=}             \\
  c_{1}T
  &=-C_{0}T+C_{1}T^2    \\
  c_{0}I
  &=C_{0}T
\end{align*}
Jumlahkan.
Ruas kirinya adalah
\( c_nT^n+c_{n-1}T^{n-1}+\cdots+c_0I \).
Ruas kanannya saling meniadakan secara berantai;
misalnya, $-C_{n-1}T^n$ dari baris pertama berpadu dengan
bagian $C_{n-1}T^n$ dari baris kedua.
Jumlah pada ruas kanan adalah matriks nol.
\end{proof}

Kita merujuk hasil tersebut dengan mengatakan bahwa suatu
matriks atau pemetaan
\definend{memenuhi}\index{karakteristik!polinomial!dipenuhi oleh}
polinomial karakteristiknya.

\begin{lemma} \label{le:tSatisImpMinPolyDivides}
Setiap polinomial yang dipenuhi oleh \( T \) habis dibagi oleh
polinomial minimal \( T \).
Artinya, untuk suatu polinomial \( f(x) \), jika \( f(T) \) merupakan matriks nol,
maka \( f(x) \) habis dibagi oleh polinomial minimal \( T \).
\end{lemma}

\begin{proof}
Misalkan \( m(x) \) adalah polinomial minimal untuk \( T \).
Teorema Pembagian Polinomial memberikan
\( f(x)=q(x)\cdot m(x)+r(x) \)
dengan derajat \( r \) lebih kecil secara ketat daripada derajat \( m \).
Karena $T$ memenuhi $f$ maupun $m$, penyubstitusian $T$ ke persamaan itu
menunjukkan bahwa \( r(T) \) adalah matriks nol.
Hal ini bertentangan dengan minimalitas \( m \), kecuali jika \( r \)
adalah polinomial nol.
\end{proof}

Menggabungkan dua lema sebelumnya menunjukkan bahwa polinomial minimal
membagi polinomial karakteristik.
Jadi,
setiap akar polinomial minimal juga merupakan akar polinomial
karakteristik.

Jadi, sejauh ini kita telah memperoleh bahwa jika polinomial minimal difaktorkan sebagai
\( m(x)=(x-\lambda_1)^{q_1}\cdots(x-\lambda_i)^{q_i} \), maka
polinomial karakteristik juga memiliki akar \( \lambda_1\), \ldots,
\( \lambda_i \).
Namun, berdasarkan apa yang telah kita buktikan sampai titik ini,
polinomial karakteristik mungkin memiliki akar tambahan:
\( c(x)=(x-\lambda_1)^{p_1}\cdots(x-\lambda_i)^{p_i}
     (x-\lambda_{i+1})^{p_{i+1}}\cdots(x-\lambda_z)^{p_z} \), dengan
\( 1\leq q_j\leq p_j \) untuk \( 1\leq j \leq i \).
Kita menyelesaikan pembuktian Teorema Cayley-Hamilton dengan menunjukkan bahwa
polinomial karakteristik tidak memiliki akar tambahan, sehingga
tidak ada $\lambda_{i+1}$, $\lambda_{i+2}$, dan seterusnya.

\begin{lemma}
Setiap faktor linear dari polinomial karakteristik suatu matriks persegi
juga merupakan faktor linear dari polinomial minimalnya.
\end{lemma}

\begin{proof}
Misalkan \( T \) adalah matriks persegi dengan polinomial minimal \( m(x) \)
berderajat~\( n \), dan
andaikan \( x-\lambda \) merupakan faktor polinomial karakteristik
\( T \), sehingga \( \lambda \) adalah nilai eigen \( T \).
Kita akan menunjukkan bahwa $m(\lambda)=0$, yakni bahwa $x-\lambda$ merupakan faktor $m$.

Andaikan $\lambda$ berkaitan dengan vektor eigen~$\vec{v}$, dan
perhatikan pangkat-pangkat $T^2(\vec{v})$, \ldots, \( T^n(\vec{v}) \).
Kita mempunyai
$T^2(\vec{v})
=T\cdot\lambda\vec{v}=\lambda\cdot T\vec{v}=\lambda^2\vec{v}$.
Hal yang sama berlaku untuk semua pangkat:
$T^i\vec{v}=\lambda^i\vec{v}\,$ untuk \( 1\leq i \leq n \).
Dengan demikian,
untuk setiap fungsi polinomial \( p(x) \), penerapan matriks \( p(T) \)
pada \( \vec{v} \) sama dengan hasil perkalian \( \vec{v} \) oleh skalar
\( p(\lambda) \).
\begin{multline*}
  p(T)\,\vec{v}
  =(c_kT^k+\dots+c_1T+c_0I)\,\vec{v}
  =c_kT^k\vec{v}+\dots+c_1T\vec{v}+c_0\vec{v}  \\
  =c_k\lambda^k\vec{v}+\dots+c_1\lambda\vec{v}+c_0\vec{v}
  =p(\lambda)\cdot\vec{v}
\end{multline*}
Karena \( m(T) \) adalah matriks nol,
\( \zero=m(T)(\vec{v})=m(\lambda)\cdot\vec{v} \), sehingga
\( m(\lambda)=0 \), sebab \( \vec{v}\neq\zero \)
karena ia merupakan vektor eigen.
\end{proof}

Ini menuntaskan pembuktian Teorema Cayley-Hamilton.

\begin{example} \label{ex:MinPolyUsingCH}
Kita dapat menggunakan Teorema Cayley-Hamilton untuk mencari polinomial minimal
matriks berikut.
\begin{equation*}
   T=
   \begin{mat}[r]
      2  &0  &0  &1  \\
      1  &2  &0  &2  \\
      0  &0  &2  &-1 \\
      0  &0  &0  &1
   \end{mat}
\end{equation*}
Pertama, kita mencari polinomial karakteristiknya \( c(x)=(x-1)(x-2)^3 \)
dengan determinan biasa $\deter{T-xI}$.
Dari sini, Teorema Cayley-Hamilton menyatakan bahwa
polinomial minimal \( T \) adalah salah satu dari
\( (x-1)(x-2) \), atau
\( (x-1)(x-2)^2 \), atau
\( (x-1)(x-2)^3 \).
Kita dapat memilih di antara kemungkinan itu cukup dengan menghitung
\begin{equation*}
   (T-1I)(T-2I)=\!
   \begin{mat}[r]
      1  &0  &0  &1  \\
      1  &1  &0  &2  \\
      0  &0  &1  &-1 \\
      0  &0  &0  &0
   \end{mat}
   \begin{mat}[r]
      0  &0  &0  &1  \\
      1  &0  &0  &2  \\
      0  &0  &0  &-1 \\
      0  &0  &0  &-1
   \end{mat}
   =
   \begin{mat}[r]
      0  &0  &0  &0  \\
      1  &0  &0  &1  \\
      0  &0  &0  &0  \\
      0  &0  &0  &0
   \end{mat}
\end{equation*}
dan
\begin{equation*}
   (T-1I)(T-2I)^2=
   \begin{mat}[r]
      0  &0  &0  &0  \\
      1  &0  &0  &1  \\
      0  &0  &0  &0  \\
      0  &0  &0  &0
   \end{mat}
   \begin{mat}[r]
      0  &0  &0  &1  \\
      1  &0  &0  &2  \\
      0  &0  &0  &-1 \\
      0  &0  &0  &-1
   \end{mat}
   =
   \begin{mat}[r]
      0  &0  &0  &0  \\
      0  &0  &0  &0  \\
      0  &0  &0  &0  \\
      0  &0  &0  &0
   \end{mat}
\end{equation*}
sehingga \( m(x)=(x-1)(x-2)^2 \).
\end{example}


% ==================================
\begin{exercises}
  \recommended \item 
    Apa saja polinomial minimal yang mungkin jika suatu matriks memiliki
    polinomial karakteristik yang diberikan?
    \begin{exparts*}
      \partsitem $(x-3)^4$
      \partsitem $(x+1)^3(x-4)$
      \partsitem $(x-2)^2(x-5)^2$
      \partsitem  \( (x+3)^2(x-1)(x-2)^2 \)
    \end{exparts*}
    Berapakah derajat setiap kemungkinan?
    \begin{answer} 
      Teorema Cayley-Hamilton \nearbytheorem{th:CayHam} menyatakan bahwa
      polinomial minimal harus memuat faktor-faktor linear yang sama
      dengan polinomial karakteristik, meskipun mungkin berpangkat lebih rendah,
      tetapi tidak berpangkat nol.
      \begin{exparts}
        \partsitem Kemungkinannya adalah
          $m_1(x)=x-3$, $m_2(x)=(x-3)^2$, $m_3(x)=(x-3)^3$,
          dan $m_4(x)=(x-3)^4$.
          Yang pertama adalah polinomial derajat satu, yang kedua berderajat dua,
          yang ketiga berderajat tiga, dan yang keempat berderajat empat.
        \partsitem Kemungkinannya adalah $m_1(x)=(x+1)(x-4)$,
          $m_2(x)=(x+1)^2(x-4)$, dan $m_3(x)=(x+1)^3(x-4)$.
          Yang pertama merupakan polinomial kuadrat, yakni berderajat dua.
          Yang kedua berderajat tiga, dan yang ketiga berderajat empat.
        \partsitem Kita mempunyai $m_1(x)=(x-2)(x-5)$, $m_2(x)=(x-2)^2(x-5)$,
          $m_3(x)=(x-2)(x-5)^2$, dan $m_4(x)=(x-2)^2(x-5)^2$.
          Polinomial-polinomial itu berturut-turut berderajat dua, tiga, tiga, dan empat.
        \partsitem Kemungkinannya adalah \( m_1(x)=(x+3)(x-1)(x-2) \),
          \( m_2(x)=(x+3)^2(x-1)(x-2) \),
          \( m_3(x)=(x+3)(x-1)(x-2)^2 \),
          dan \( m_4(x)=(x+3)^2(x-1)(x-2)^2 \).
          Derajat $m_1$ adalah tiga, derajat $m_2$ adalah empat,
          derajat $m_3$ adalah empat, dan derajat $m_4$ adalah lima.
      \end{exparts}
    \end{answer}
  \item Untuk matriks berikut,
    \begin{equation*}
      \begin{mat}
        0 &1 &0 &1 \\
        1 &0 &1 &0 \\
        0 &1 &0 &1 \\
        1 &0 &1 &0 \\
      \end{mat}
    \end{equation*}
    carilah polinomial karakteristik dan polinomial minimalnya.
    \begin{answer}
    Polinomial karakteristik
    \begin{equation*}
      \deter{T-xI}=
      \begin{vmat}
        -x &1  &0  &1 \\
         1 &-x &1  &0 \\
         0 &1  &-x &1 \\
         1 &0  &1  &-x \\
      \end{vmat}
    \end{equation*}
    adalah $x^2(x^2-4)$.

    Berdasarkan Teorema Cayley-Hamilton, semua akar berbeda dari
    polinomial karakteristik juga merupakan akar polinomial minimal, sehingga
    polinomial minimal berakar $0$, $2$, dan~$-2$.
    Jadi, polinomial minimalnya adalah $x(x^2-4)$ atau $x^2(x^2-4)$.
    
    Salah satu cara mencari polinomial minimal adalah menyubstitusikan matriks itu
    ke kedua bentuk hingga kita memperoleh matriks nol.
    \begin{equation*}
      T(T^2-4I)=
      \begin{mat}
        0 &1 &0 &1 \\
        1 &0 &1 &0 \\
        0 &1 &0 &1 \\
        1 &0 &1 &0 \\
      \end{mat}
      \begin{mat}
        -2 &0  &1  &0 \\
        0  &-2 &0  &2 \\
        2  &0  &-2 &0 \\
        0  &2  &0  &-2 \\
      \end{mat}
      =
      \begin{mat}
        0 &0 &0 &0 \\
        0 &0 &0 &0 \\
        0 &0 &0 &0 \\
        0 &0 &0 &0 \\
      \end{mat}
    \end{equation*}
    Jadi, polinomial minimalnya adalah $x(x^2-4)$.
    \end{answer}
   \item 
     Carilah polinomial minimal matriks berikut.
     \begin{equation*}
        \begin{mat}[r]
           0  &1  &0  \\
           0  &0  &1  \\
           1  &0  &0
        \end{mat}
     \end{equation*}
     \begin{answer}
       Polinomial karakteristiknya memiliki akar kompleks.
       \begin{equation*}
          \begin{vmat}
                   -x  &1  &0  \\
                    0  &-x &1  \\
                    1  &0  &-x
          \end{vmat}
          =(1-x)\cdot (x-(-\frac{1}{2}+\frac{\sqrt{3}}{2}i))
                \cdot (x-(-\frac{1}{2}-\frac{\sqrt{3}}{2}i))
       \end{equation*}
       Karena akar-akarnya berbeda, polinomial karakteristik sama dengan
       polinomial minimal.
     \end{answer}
  \recommended \item 
    Carilah polinomial minimal setiap matriks.
    \begin{exparts*}
       \partsitem \( \begin{mat}[r]
                   3  &0  &0  \\
                   1  &3  &0  \\
                   0  &0  &4
                \end{mat} \)
       \partsitem \( \begin{mat}[r]
                   3  &0  &0  \\
                   1  &3  &0  \\
                   0  &0  &3
                \end{mat} \)
       \partsitem \( \begin{mat}[r]
                   3  &0  &0  \\
                   1  &3  &0  \\
                   0  &1  &3
                \end{mat} \)
       \partsitem \( \begin{mat}[r]
                   2  &0  &1  \\
                   0  &6  &2  \\
                   0  &0  &2
                \end{mat} \)
       \partsitem \( \begin{mat}[r]
                   2  &2  &1  \\
                   0  &6  &2  \\
                   0  &0  &2
                \end{mat} \)
       \partsitem \( \begin{mat}[r]
                   -1 &4  &0  &0  &0  \\
                    0 &3  &0  &0  &0  \\
                    0 &-4 &-1 &0  &0  \\
                    3 &-9 &-4 &2  &-1 \\
                    1 &5  &4  &1  &4
                \end{mat} \)
    \end{exparts*}
    \begin{answer}
      Dalam setiap kasus, kita akan menggunakan metode pada \nearbyexample{ex:MinPolyUsingCH}.
      \begin{exparts}
       \partsitem Karena $T$ berbentuk segitiga, $T-xI$ juga berbentuk segitiga,
         \begin{equation*}
           T-xI=
           \begin{mat}
             3-x  &0    &0   \\
             1    &3-x  &0   \\
             0    &0    &4-x
           \end{mat}
         \end{equation*}
         sehingga polinomial karakteristiknya mudah dihitung,
         $c(x)=\deter{T-xI}=(3-x)^2(4-x)=-1\cdot (x-3)^2(x-4)$.
         Hanya ada dua kemungkinan untuk polinomial minimal,
         $m_1(x)=(x-3)(x-4)$ dan $m_2(x)=(x-3)^2(x-4)$.
         (Perhatikan bahwa polinomial karakteristik bertanda negatif,
         sedangkan polinomial minimal tidak, karena polinomial minimal harus
         memiliki koefisien utama satu.)
         Karena $m_1(T)$ bukan matriks nol,
         \begin{equation*}
           (T-3I)(T-4I)
           =
           \begin{mat}[r]
             0  &0  &0  \\
             1  &0  &0  \\
             0  &0  &1
           \end{mat}
           \begin{mat}[r]
             -1  &0  &0  \\
              1  &-1 &0  \\
              0  &0  &0
           \end{mat}
           =
           \begin{mat}[r]
             0  &0  &0  \\
            -1  &0  &0  \\
             0  &0  &0
           \end{mat}
         \end{equation*}
         polinomial minimalnya adalah $m(x)=m_2(x)$.
         \begin{multline*}
           (T-3I)^2(T-4I)
           =(T-3I)\cdot\bigl((T-3I)(T-4I)\bigr)              \\
           =
           \begin{mat}
             0  &0  &0  \\
             1  &0  &0  \\
             0  &0  &1
           \end{mat}
           \begin{mat}[r]
              0  &0  &0  \\
             -1  &0  &0  \\
              0  &0  &0
           \end{mat}
           =
           \begin{mat}[r]
             0  &0  &0  \\
             0  &0  &0  \\
             0  &0  &0
           \end{mat}
         \end{multline*}
       \partsitem Seperti pada butir sebelumnya, karena matriksnya berbentuk
        segitiga, polinomial karakteristik mudah dihitung.
        \begin{equation*}
          c(x)=\deter{T-xI}
              =
              \begin{vmat}
                3-x  &0   &0   \\
                1    &3-x &0   \\
                0    &0   &3-x
              \end{vmat}
              =(3-x)^3=-1\cdot (x-3)^3
        \end{equation*}
        Ada tiga kemungkinan untuk polinomial minimal,
        $m_1(x)=(x-3)$, $m_2(x)=(x-3)^2$, dan $m_3(x)=(x-3)^3$.
        Kita menentukannya dengan menghitung $m_1(T)$
        \begin{equation*}
          T-3I=
          \begin{mat}[r]
            0  &0  &0  \\
            1  &0  &0  \\
            0  &0  &0
          \end{mat}
        \end{equation*}
        dan $m_2(T)$.
        \begin{equation*}
          (T-3I)^2=
          \begin{mat}[r]
            0  &0  &0  \\
            1  &0  &0  \\
            0  &0  &0
          \end{mat}
          \begin{mat}[r]
            0  &0  &0  \\
            1  &0  &0  \\
            0  &0  &0
          \end{mat}
          =          
          \begin{mat}[r]
            0  &0  &0  \\
            0  &0  &0  \\
            0  &0  &0
          \end{mat}
        \end{equation*}
        Karena $m_2(T)$ adalah matriks nol, $m_2(x)$ merupakan polinomial
        minimal.
       \partsitem Sekali lagi, matriksnya berbentuk segitiga.
        \begin{equation*}
          c(x)=\deter{T-xI}
              =
              \begin{vmat}
                3-x  &0   &0   \\
                1    &3-x &0   \\
                0    &1   &3-x
              \end{vmat}
              =(3-x)^3=-1\cdot (x-3)^3
        \end{equation*}
        Sekali lagi, ada tiga kemungkinan untuk polinomial minimal,
        $m_1(x)=(x-3)$, $m_2(x)=(x-3)^2$, dan $m_3(x)=(x-3)^3$.
        Kita menghitung $m_1(T)$
        \begin{equation*}
          T-3I=
          \begin{mat}[r]
            0  &0  &0  \\
            1  &0  &0  \\
            0  &1  &0
          \end{mat}
        \end{equation*}
        dan $m_2(T)$
        \begin{equation*}
          (T-3I)^2=
          \begin{mat}[r]
            0  &0  &0  \\
            1  &0  &0  \\
            0  &1  &0
          \end{mat}
          \begin{mat}[r]
            0  &0  &0  \\
            1  &0  &0  \\
            0  &1  &0
          \end{mat}
          =          
          \begin{mat}[r]
            0  &0  &0  \\
            0  &0  &0  \\
            1  &0  &0
          \end{mat}
        \end{equation*}
        serta $m_3(T)$.
        \begin{equation*}
          (T-3I)^3
          =(T-3I)^2(T-3I)
          =
          \begin{mat}[r]
            0  &0  &0  \\
            0  &0  &0  \\
            1  &0  &0
          \end{mat}
          \begin{mat}[r]
            0  &0  &0  \\
            1  &0  &0  \\
            0  &1  &0
          \end{mat}
          =          
          \begin{mat}[r]
            0  &0  &0  \\
            0  &0  &0  \\
            0  &0  &0
          \end{mat}
        \end{equation*}
        Jadi, polinomial minimalnya adalah $m(x)=m_3(x)=(x-3)^3$.
       \partsitem Kasus ini juga berbentuk segitiga, kali ini segitiga atas.
         \begin{equation*}
           c(x)=\deter{T-xI}=
           \begin{vmat}
             2-x  &0   &1     \\
             0    &6-x &2     \\
             0    &0   &2-x
           \end{vmat}
           =(2-x)^2(6-x)=-(x-2)^2(x-6)
         \end{equation*}
         Ada dua kemungkinan untuk polinomial minimal,
         $m_1(x)=(x-2)(x-6)$ dan $m_2(x)=(x-2)^2(x-6)$.
         Perhitungan menunjukkan bahwa polinomial minimalnya bukan $m_1(x)$.
         \begin{equation*}
           (T-2I)(T-6I)=
           \begin{mat}[r]
             0  &0  &1  \\
             0  &4  &2  \\
             0  &0  &0  
           \end{mat}
           \begin{mat}[r]
             -4  &0  &1  \\
              0  &0  &2  \\
              0  &0  &-4
           \end{mat}
           =
           \begin{mat}[r]
             0  &0  &-4  \\
             0  &0  &0   \\
             0  &0  &0
           \end{mat}
         \end{equation*}
         Karena itu, haruslah $m(x)=m_2(x)=(x-2)^2(x-6)$.
         Berikut verifikasinya.
         \begin{multline*}
           (T-2I)^2(T-6I)=(T-2I)\cdot\bigl((T-2I)(T-6I)\bigr)             \\
           =
           \begin{mat}[r]
             0  &0  &1  \\
             0  &4  &2  \\
             0  &0  &0  
           \end{mat}
           \begin{mat}[r]
              0  &0  &-4   \\
              0  &0  &0   \\
              0  &0  &0
           \end{mat}
           =
           \begin{mat}[r]
             0  &0  &0  \\
             0  &0  &0   \\
             0  &0  &0
           \end{mat}
         \end{multline*}
       \partsitem Polinomial karakteristiknya adalah
         \begin{equation*}
           c(x)=\deter{T-xI}=
           \begin{vmat}
             2-x  &2   &1     \\
             0    &6-x &2     \\
             0    &0   &2-x
           \end{vmat}
           =(2-x)^2(6-x)=-(x-2)^2(x-6)
         \end{equation*}
         dan ada dua kemungkinan untuk polinomial minimal,
         $m_1(x)=(x-2)(x-6)$ dan $m_2(x)=(x-2)^2(x-6)$.
         Dengan memeriksa kemungkinan pertama,
         \begin{equation*}
           (T-2I)(T-6I)=
           \begin{mat}[r]
             0  &2  &1  \\
             0  &4  &2  \\
             0  &0  &0  
           \end{mat}
           \begin{mat}[r]
             -4  &2  &1  \\
              0  &0  &2  \\
              0  &0  &-4
           \end{mat}
           =
           \begin{mat}[r]
             0  &0  &0  \\
             0  &0  &0   \\
             0  &0  &0
           \end{mat}
         \end{equation*}
         terlihat bahwa polinomial minimalnya adalah
         $m(x)=m_1(x)=(x-2)(x-6)$.
       \partsitem Berikut polinomial karakteristiknya.
         \begin{equation*}
           c(x)=\deter{T-xI}=
           \begin{vmat} 
              -1-x &4    &0    &0    &0    \\
               0   &3-x  &0    &0    &0    \\
               0   &-4   &-1-x &0    &0    \\
               3   &-9   &-4   &2-x  &-1   \\
               1   &5    &4    &1    &4-x
           \end{vmat}     
           =-(x-3)^3(x+1)^2
         \end{equation*}
         Berikut kemungkinan untuk polinomial minimalnya,
         diurutkan menurut derajat yang meningkat:
         $m_1(x)=(x-3)(x+1)$, $m_2(x)=(x-3)^2(x+1)$, $m_3(x)=(x-3)(x+1)^2$,
         $m_4(x)=(x-3)^3(x+1)$, $m_5(x)=(x-3)^2(x+1)^2$,
         dan $m_6(x)=(x-3)^3(x+1)^2$.
         Kemungkinan pertama tidak berhasil,
         \begin{align*}
           (T-3I)(T+1I)
           &=
           \begin{mat}[r] 
              -4   &4    &0    &0    &0    \\
               0   &0    &0    &0    &0    \\
               0   &-4   &-4   &0    &0    \\
               3   &-9   &-4   &-1   &-1   \\
               1   &5    &4    &1    &1  
           \end{mat}     
           \begin{mat}[r] 
               0   &4    &0    &0    &0    \\
               0   &4    &0    &0    &0    \\
               0   &-4   &0    &0    &0    \\
               3   &-9   &-4   &3    &-1   \\
               1   &5    &4    &1    &5  
           \end{mat}                           \\     
           &=
           \begin{mat}[r] 
               0   &0    &0    &0    &0    \\
               0   &0    &0    &0    &0    \\
               0   &0    &0    &0    &0    \\
              -4   &-4   &0    &-4   &-4   \\
               4   &4    &0    &4    &4  
           \end{mat}           
         \end{align*}
         tetapi kemungkinan kedua berhasil.
         \begin{multline*}
           (T-3I)^2(T+1I)=(T-3I)\bigl((T-3I)(T+1I)\bigr) \\
           \begin{aligned}
           &=
           \begin{mat}[r] 
              -4   &4    &0    &0    &0    \\
               0   &0    &0    &0    &0    \\
               0   &-4   &-4   &0    &0    \\
               3   &-9   &-4   &-1   &-1   \\
               1   &5    &4    &1    &1  
           \end{mat}     
           \begin{mat}[r] 
               0   &0    &0    &0    &0    \\
               0   &0    &0    &0    &0    \\
               0   &0    &0    &0    &0    \\
              -4   &-4   &0    &-4   &-4   \\
               4   &4    &0    &4    &4  
           \end{mat}                          \\          
           &=
           \begin{mat}[r] 
               0   &0    &0    &0    &0    \\
               0   &0    &0    &0    &0    \\
               0   &0    &0    &0    &0    \\
               0   &0    &0    &0    &0    \\
               0   &0    &0    &0    &0  
           \end{mat}  
           \end{aligned}         
         \end{multline*}
         Polinomial minimalnya adalah \( m(x)=(x-3)^2(x+1) \).
      \end{exparts} 
    \end{answer}
  \recommended \item 
     Apa polinomial minimal operator diferensiasi
     $d/dx$ pada \( \polyspace_n \)?
     \begin{answer}
       Kita mengetahui bahwa $\polyspace_n$ merupakan ruang berdimensi $n+1$ dan
       operator diferensiasi bersifat
       nilpoten dengan indeks~$n+1$ (misalnya, dengan mengambil $n=3$,
       $\polyspace_3=\set{c_3x^3+c_2x^2+c_1x+c_0\suchthat c_3,\ldots,c_0\in\C}$
       dan turunan keempat suatu polinomial kubik adalah polinomial nol).
       Representasikan operator ini menggunakan bentuk kanonik
       untuk transformasi nilpoten.
       \begin{equation*}
         \begin{mat}
           0  &0  &0  &\ldots &  &0  \\
           1  &0  &0  &       &  &0  \\
           0  &1  &0  &       &  &   \\
              &   &\ddots            \\
           0  &0  &0  &       &1 &0 
         \end{mat}
       \end{equation*}
       Ini adalah matriks $\nbyn{(n+1)}$ dengan polinomial
       karakteristik yang mudah dihitung,
        $c(x)=(-x)^{n+1}$.
       (\textit{Catatan:} matriks ini adalah $\rep{d/dx}{B,B}$ dengan
        $B=\sequence{x^n,nx^{n-1},n(n-1)x^{n-2},\ldots,n!}$.)
       Untuk mencari polinomial minimal seperti dalam \nearbyexample{ex:MinPolyUsingCH},
       kita memperhatikan pangkat-pangkat $T-0I=T$.
       Namun, tentu saja, pangkat pertama $T$ yang menghasilkan matriks nol
       adalah pangkat $n+1$.
       Jadi, polinomial minimalnya adalah \( m(x)=x^{n+1} \).
     \end{answer}
  \recommended \item 
    Carilah polinomial minimal matriks berbentuk berikut,
    \begin{equation*}
      \begin{mat}
        \lambda  &0        &0          &\ldots  &        &0  \\
        1        &\lambda  &0          &        &        &0  \\
        0        &1        &\lambda                          \\
                 &         &           &\ddots                \\
                 &         &           &        &\lambda &0   \\
        0        &0        &\ldots     &        &1       &\lambda
      \end{mat}
    \end{equation*}
    dengan skalar $\lambda$ tetap (yakni bukan variabel).
    \begin{answer}
      Namakan matriks itu $T$ dan andaikan ukurannya \( \nbyn{n} \).
      Karena $T$ berbentuk segitiga, maka $T-xI$ juga berbentuk segitiga,
       dan polinomial karakteristiknya adalah
       $c(x)=(\lambda-x)^n=(-1)^n(x-\lambda)^n$.
       Untuk melihat bahwa polinomial minimalnya adalah bentuk monik yang
       bersesuaian, $m(x)=(x-\lambda)^n$, perhatikan
      $T-\lambda I$.
      \begin{equation*}
        \begin{mat}
          0        &0        &0          &\ldots  &0  \\
          1        &0        &0          &\ldots  &0  \\
          0        &1        &0                       \\
                   &         &\ddots                  \\
          0        &0        &\ldots     &1       &0      
        \end{mat}
      \end{equation*}
      Kenali matriks ini sebagai bentuk kanonik untuk transformasi yang
      nilpoten berderajat~$n$; pangkat $(T-\lambda I)^j$ pertama kali menjadi nol
      ketika $j$ sama dengan $n$.
    \end{answer}
  \item 
    Apa polinomial minimal transformasi pada
    \( \polyspace_n \) yang mengirim \( p(x) \) ke \( p(x+1) \)?
    \begin{answer}
      Kasus $n=3$ memberikan petunjuk.
      Basis alami untuk $\polyspace_3$ adalah
      $B=\sequence{1,x,x^2,x^3}$.
      Tindakan transformasinya adalah
      \begin{equation*}
        1\mapsto 1
        \quad
        x\mapsto x+1
        \quad
        x^2\mapsto x^2+2x+1
        \quad
        x^3\mapsto x^3+3x^2+3x+1
      \end{equation*}
      sehingga representasi $\rep{t}{B,B}$ adalah matriks segitiga atas berikut.
      \begin{equation*}
        \begin{mat}[r]
          1  &1  &1  &1  \\
          0  &1  &2  &3  \\
          0  &0  &1  &3  \\
          0  &0  &0  &1
        \end{mat}
      \end{equation*}
      Karena matriks itu berbentuk segitiga, jelas bahwa polinomial karakteristiknya
      adalah $c(x)=(x-1)^4$.
      Untuk polinomial minimal, kandidatnya adalah $m_1(x)=(x-1)$,
      \begin{equation*}
        T-1I=
        \begin{mat}[r]
          0  &1  &1  &1  \\
          0  &0  &2  &3  \\
          0  &0  &0  &3  \\
          0  &0  &0  &0
        \end{mat}
      \end{equation*}
      $m_2(x)=(x-1)^2$, 
      \begin{equation*}
        (T-1I)^2=
        \begin{mat}[r]
          0  &0  &2  &6  \\
          0  &0  &0  &6  \\
          0  &0  &0  &0  \\
          0  &0  &0  &0
        \end{mat}
      \end{equation*}
      $m_3(x)=(x-1)^3$,
      \begin{equation*}
        (T-1I)^3=
        \begin{mat}[r]
          0  &0  &0  &6  \\
          0  &0  &0  &0  \\
          0  &0  &0  &0  \\
          0  &0  &0  &0
        \end{mat}
      \end{equation*}
      dan $m_4(x)=(x-1)^4$.
      Karena $m_1$, $m_2$, dan $m_3$ tidak tepat, $m_4$ pasti tepat,
      sebagaimana mudah diverifikasi.
      
      Dalam kasus $n$ umum, representasinya adalah matriks segitiga
      atas dengan satu pada diagonalnya.
       Jadi, polinomial karakteristiknya adalah $c(x)=(1-x)^{n+1}$.
       Salah satu cara memverifikasi bahwa polinomial minimalnya adalah bentuk
       monik yang bersesuaian, $m(x)=(x-1)^{n+1}$, ialah menggunakan argumen berikut:
      sebutlah matriks segitiga atas sebagai segitiga-atas-$0$ jika
      terdapat entri tak nol pada diagonal, sebagai segitiga-atas-$1$ jika
      diagonalnya hanya memuat nol dan terdapat entri tak nol tepat di atas
      diagonal, dan seterusnya.
      Seperti diperlihatkan oleh contoh di atas, argumen induksi akan
      menunjukkan bahwa jika $T$ hanya memiliki entri tak negatif,
      maka $T^j$ berbentuk segitiga-atas-$j$.
      % We leave that argument to the reader.
    \end{answer}
   \item 
     Apa polinomial minimal dari
     pemetaan \( \map{\pi}{\C^3}{\C^3} \)
     yang memproyeksikan ke dua koordinat pertama?
      \begin{answer}
        Penerapan pemetaan dua kali sama dengan penerapannya sekali:~$\composed{\pi}{\pi}=\pi$,
        yakni $\pi^2=\pi$, sehingga polinomial minimal berderajat
        paling tinggi dua karena \( m(x)=x^2-x \) dapat digunakan.
        Fakta bahwa tidak ada polinomial linear yang dapat digunakan diperoleh dengan menerapkan
        pemetaan pada ruas kiri dan kanan
        $c_1\cdot \pi+c_0\cdot \identity=z$ (dengan $z$ merupakan pemetaan nol)
        pada kedua vektor berikut.
        \begin{equation*}
          \colvec[r]{0 \\ 0 \\ 1}
          \qquad
          \colvec[r]{1 \\ 0 \\ 0}
        \end{equation*}
        Jadi, polinomial minimalnya adalah $m(x)=x^2-x$.
      \end{answer}
   \item 
     Carilah matriks \( \nbyn{3} \) yang polinomial
     minimalnya adalah \( x^2 \).
     \begin{answer}
        Berikut salah satu jawabannya.
        \begin{equation*}
            \begin{mat}[r]
              0  &0  &0  \\
              1  &0  &0  \\
              0  &0  &0
            \end{mat}
        \end{equation*} 
      \end{answer}
  \item 
     Apa yang salah dalam pembuktian yang diklaim berikut untuk
     \nearbylemma{le:MatSatItsCharPoly}:
      ``jika \( c(x)=\deter{T-xI} \), maka \( c(T)=\deter{T-TI}=0 \,\)''\,?
     \cite{Cullen}
     \begin{answer}
       \( x \) harus berupa skalar, bukan matriks.
     \end{answer}
  \item 
    Verifikasikan \nearbylemma{le:MatSatItsCharPoly} untuk matriks
    \( \nbyn{2} \) melalui perhitungan langsung.
    \begin{answer}
      Polinomial karakteristik dari
      \begin{equation*}
         T=\begin{mat}
              a  &b  \\
              c  &d
           \end{mat}
      \end{equation*}
      adalah \( (a-x)(d-x)-bc=x^2-(a+d)x+(ad-bc) \).
      Substitusikan
      \begin{multline*}
         \begin{mat}
              a  &b  \\
              c  &d
         \end{mat}^2
         -
         (a+d)\begin{mat}
            a  &b  \\
            c  &d
         \end{mat}
         +
         (ad-bc)\begin{mat}[r]
            1  &0  \\
            0  &1
         \end{mat}                            \\
         =                     
         \begin{mat}
            a^2+bc  &ab+bd  \\
            ac+cd   &bc+d^2
         \end{mat}
         -
         \begin{mat}
            a^2+ad  &ab+bd   \\
            ac+cd   &ad+d^2
         \end{mat}
         +
         \begin{mat}
            ad-bc  &0      \\
            0      &ad-bc
         \end{mat}
      \end{multline*}
      lalu periksa jumlah pada setiap entri untuk melihat bahwa hasilnya adalah matriks nol.
    \end{answer}
  \recommended \item
    Buktikan bahwa polinomial minimal matriks \( \nbyn{n} \) berderajat
    paling tinggi \( n \) (bukan \( n^2 \), seperti yang mungkin diduga dari
    pembukaan subbagian ini).
    Verifikasikan bahwa nilai maksimum \( n \) ini dapat terjadi.
    \begin{answer}
      Berdasarkan Teorema Cayley-Hamilton, derajat polinomial minimal
      kurang dari atau sama dengan derajat polinomial karakteristik, yaitu
      \( n \).
      \nearbyexample{ex:MinPolyForRotMat} menunjukkan bahwa \( n \) dapat terjadi.
    \end{answer}
   \recommended \item 
     Tunjukkan bahwa pada ruang vektor nontrivial, suatu transformasi linear
     nilpoten jika dan hanya jika satu-satunya nilai eigennya adalah nol.
     \begin{answer}
       Misalkan transformasi linearnya adalah $\map{t}{V}{V}$.
       Jika $t$ nilpoten, maka terdapat $n$ sedemikian sehingga $t^n$ merupakan pemetaan nol,
       sehingga $t$ memenuhi polinomial $p(x)=x^n=(x-0)^n$.
       Berdasarkan \nearbylemma{le:tSatisImpMinPolyDivides}, polinomial minimal
       $t$ membagi $p$, sehingga satu-satunya akar polinomial minimal adalah nol.
       Berdasarkan Cayley-Hamilton, \nearbytheorem{th:CayHam},
       satu-satunya akar polinomial karakteristik adalah nol.
       Jadi, satu-satunya nilai eigen $t$ adalah nol.

       Sebaliknya, jika transformasi \( t \) pada ruang berdimensi
       \( n \) hanya memiliki nilai eigen nol,
       maka polinomial karakteristiknya adalah \( x^n \).
       Teorema Cayley-Hamilton menyatakan bahwa suatu pemetaan memenuhi
       polinomial karakteristiknya, sehingga \( t^n \) merupakan pemetaan nol.
       Jadi, $t$ nilpoten.
     \end{answer}
   \item 
       Apa polinomial minimal suatu pemetaan atau matriks nol?
       Bagaimana dengan pemetaan atau matriks identitas?
       \begin{answer}
         Polinomial minimal harus memiliki koefisien utama $1$,
         sehingga jika polinomial minimal suatu pemetaan atau matriks
         berderajat nol, maka polinomial itu adalah $m(x)=1$.
         Namun, pemetaan atau matriks identitas sama dengan pemetaan atau matriks nol
         hanya pada ruang vektor trivial.

         Jadi, dalam kasus nontrivial, polinomial minimal harus berderajat
         setidaknya satu.
         Pemetaan atau matriks nol memiliki polinomial minimal \( m(x)=x \), sedangkan
         pemetaan atau matriks identitas memiliki polinomial minimal \( m(x)=x-1 \).
       \end{answer}
   \recommended \item 
     Apa polinomial minimal suatu matriks diagonal?
     \begin{answer}
       Untuk matriks diagonal
       \begin{equation*}
          T=
          \begin{mat}
             t_{1,1}   &0        \\
             0         &t_{2,2}  \\
                       &        &\ddots  \\
                       &        &      &t_{n,n}
          \end{mat}
       \end{equation*}
       polinomial karakteristiknya adalah
       $(t_{1,1}-x)(t_{2,2}-x)\cdots (t_{n,n}-x)$.     
       Tentu saja, sebagian faktor itu mungkin berulang; misalnya, matriksnya mungkin
       memenuhi $t_{1,1}=t_{2,2}$.
       Sebagai contoh, polinomial karakteristik dari
       \begin{equation*}
          D=
          \begin{mat}[r]
             3 &0 &0  \\
             0 &3 &0  \\
             0 &0 &1
          \end{mat}
       \end{equation*}
       adalah \( (3-x)^2(1-x)=-1\cdot (x-3)^2(x-1) \).

       Untuk membentuk polinomial minimal,
       ambil faktor-faktor \( x-t_{i,i} \), buang pengulangannya,
       lalu kalikan semuanya.
       Sebagai contoh, polinomial minimal $D$
       adalah \( (x-3)(x-1) \).
       Untuk memeriksanya, pertama-tama perhatikan bahwa \nearbytheorem{th:CayHam},
       Teorema Cayley-Hamilton, mengharuskan setiap faktor linear dalam
       polinomial karakteristik muncul setidaknya sekali dalam polinomial
       minimal.
       Salah satu cara memeriksa arah sebaliknya\Dash bahwa dalam kasus
       matriks diagonal,
       setiap faktor linear hanya perlu muncul paling banyak sekali\Dash adalah
       menggunakan argumen matriks.
       Ketika mengalikan dari kiri, matriks diagonal menskalakan baris-baris
       menurut entri pada diagonal.
       Namun, dalam hasil kali $(T-t_{1,1}I)\cdots\hbox{}$, bahkan tanpa faktor
       berulang, setiap baris bernilai nol pada setidaknya salah satu faktor.

       Sebagai contoh, dalam hasil kali
       \begin{equation*}
         (D-3I)(D-1I)=(D-3I)(D-1I)I=
         \begin{mat}[r]
           0  &0  &0  \\
           0  &0  &0  \\
           0  &0  &-2        
         \end{mat}
         \begin{mat}[r]
           2  &0  &0  \\
           0  &2  &0  \\
           0  &0  &0
         \end{mat}
         \begin{mat}[r]
           1  &0  &0  \\
           0  &1  &0  \\
           0  &0  &1
         \end{mat}
       \end{equation*}
       karena baris pertama dan kedua matriks pertama $D-3I$
       bernilai nol, seluruh hasil kali akan memiliki baris pertama dan kedua
       yang bernilai nol.
       Dan karena baris ketiga matriks tengah $D-1I$ bernilai nol,
       seluruh hasil kali memiliki baris ketiga yang bernilai nol.
    \end{answer}
  \recommended \item 
    Suatu \definend{proyeksi}\index{proyeksi}%
    \index{transformasi!proyeksi} 
    adalah sebarang transformasi \( t \) sedemikian sehingga \( t^2=t \).
    (Sebagai contoh, perhatikan transformasi pada bidang $\Re^2$ yang memproyeksikan
    setiap vektor ke koordinat pertamanya.
    Jika kita memproyeksikan dua kali, hasilnya sama seperti jika kita memproyeksikan sekali.)
    Apa polinomial minimal suatu proyeksi?
    \begin{answer}
      Subbagian ini diawali dengan pengamatan bahwa pangkat-pangkat
      suatu transformasi linear tidak dapat terus meningkat tanpa ``pengulangan'',
      yakni bahwa untuk suatu pangkat~$n$ terdapat hubungan linear
      $c_n\cdot t^n+\dots+c_1\cdot t+c_0\cdot \identity=z$ dengan $z$ sebagai
      transformasi nol.
      Definisi proyeksi menyatakan bahwa untuk pemetaan semacam itu,
      salah satu hubungan linearnya bersifat kuadratik, $t^2-t=z$.
      Untuk menyelesaikannya, kita hanya perlu mempertimbangkan apakah hubungan ini
      mungkin tidak minimal; yakni, adakah proyeksi yang polinomial minimalnya
      konstan atau linear?

      Agar polinomial minimalnya konstan, pemetaan harus memenuhi
      $c_0\cdot\identity=z$, dengan $c_0=1$ karena koefisien utama
      polinomial minimal adalah $1$.
      Persamaan ini hanya dipenuhi oleh transformasi nol pada ruang trivial.
      Ini memang suatu proyeksi, tetapi bukan proyeksi yang menarik.

      Jika polinomial minimal suatu transformasi linear, kita memperoleh
      $c_1\cdot t+c_0\cdot\identity=z$ dengan $c_1=1$.
      Persamaan ini memberikan $t=-c_0\cdot \identity$.
      Menggabungkannya dengan syarat $t^2=t$ memberikan
      $t^2=(-c_0)^2\cdot\identity=-c_0\cdot\identity$, yang menghasilkan
      $c_0=0$ dan $t$ merupakan transformasi nol, atau $c_0=1$ dan
      $t$ merupakan identitas.

      Jadi, kecuali ketika proyeksinya merupakan pemetaan nol atau
      pemetaan identitas, polinomial minimalnya adalah $m(x)=x^2-x$.
    \end{answer}
  \item \label{exer:SomeRootsMinPolyAreEigs}
    \textit{Dua butir pertama soal ini merupakan tinjauan ulang.}
    \begin{exparts}
      \partsitem Buktikan bahwa komposisi pemetaan satu-satu juga
        satu-satu.
      \partsitem Buktikan bahwa jika suatu pemetaan linear tidak satu-satu,
        maka setidaknya satu vektor tak nol dari domain dipetakan ke
        vektor nol dalam kodomain.
      \partsitem Verifikasikan pernyataan berikut, yang dikutip dari bagian
         sebelum \nearbytheorem{th:CayHam}.
         \begin{quotation}
           \noindent \ldots{} 
           jika polinomial minimal $m(x)$ untuk transformasi $t$
           difaktorkan sebagai
           $m(x)=(x-\lambda_1)^{q_1}\cdots (x-\lambda_z)^{q_z}$
           maka
           \( m(t)=\composed{\composed{(t-\lambda_1)^{q_1}}{\cdots}}{
                                                        (t-\lambda_z)^{q_z}} \) 
          merupakan pemetaan nol.
          Karena \( m(t) \) mengirim setiap vektor ke nol, setidaknya
          salah satu pemetaan \( t-\lambda_i \) mengirim beberapa
          vektor tak nol ke nol. \ldots{}
          Artinya, \ldots{}
          setidaknya sebagian dari \( \lambda_i \) merupakan nilai eigen.
        \end{quotation}
    \end{exparts}
    \begin{answer}
      \begin{exparts}
       \partsitem \textit{Ini merupakan sifat fungsi secara umum,
          bukan hanya fungsi linear.}
          Misalkan $f$ dan $g$ adalah fungsi satu-satu sedemikian sehingga
          $\composed{f}{g}$ terdefinisi.
          Misalkan $\composed{f}{g}(x_1)=\composed{f}{g}(x_2)$, sehingga
          $f(g(x_1))=f(g(x_2))$.
          Karena $f$ satu-satu, ini mengakibatkan $g(x_1)=g(x_2)$.
          Karena $g$ juga satu-satu, hal ini selanjutnya mengakibatkan
          $x_1=x_2$.
          Jadi, singkatnya, $\composed{f}{g}(x_1)=\composed{f}{g}(x_2)$
          mengakibatkan $x_1=x_2$, sehingga $\composed{f}{g}$ satu-satu.
        \partsitem Jika pemetaan linear $h$
          tidak satu-satu, maka terdapat vektor berbeda
          $\vec{v}_1$, $\vec{v}_2$ yang dipetakan ke nilai yang sama,
          $h(\vec{v}_1)=h(\vec{v}_2)$.
          Karena $h$ linear, kita mempunyai
          $\zero=h(\vec{v}_1)-h(\vec{v}_2)=h(\vec{v}_1-\vec{v}_2)$
          sehingga $\vec{v}_1-\vec{v}_2$ merupakan vektor tak nol dari domain
          yang dipetakan $h$ ke vektor nol dalam kodomain
          ($\vec{v}_1-\vec{v}_2$
          tidak sama dengan vektor nol domain karena $\vec{v}_1$
          tidak sama dengan $\vec{v}_2$).
        \partsitem Polinomial minimal
          $m(t)$ mengirim setiap vektor dalam domain ke
          nol, sehingga tidak satu-satu (kecuali dalam ruang trivial, yang
          kita abaikan).
          Berdasarkan butir pertama soal ini,
          karena komposisi $m(t)$ tidak satu-satu,
          setidaknya salah satu komponennya $t-\lambda_i$ tidak satu-satu.
          Berdasarkan butir kedua, $t-\lambda_i$ memiliki ruang nol nontrivial.
          Karena $(t-\lambda_i)(\vec{v})=\zero$ berlaku jika dan hanya jika
          $t(\vec{v})=\lambda_i\cdot\vec{v}$, kalimat sebelumnya menunjukkan bahwa
          $\lambda_i$ merupakan nilai eigen (ingat bahwa definisi
          nilai eigen mensyaratkan hubungan tersebut berlaku untuk setidaknya satu
          $\vec{v}$ tak nol).
      \end{exparts}
    \end{answer}
  \item 
    Benar atau salah:~untuk transformasi pada ruang berdimensi
    \( n \), jika polinomial minimal berderajat \( n \),
    maka pemetaannya dapat didiagonalkan.
    \begin{answer}
      Pernyataan ini salah.
      Contoh alami transformasi yang tidak dapat didiagonalkan dapat digunakan di sini.
      Perhatikan transformasi pada $\C^2$ yang direpresentasikan terhadap
      basis standar oleh matriks berikut.
      \begin{equation*}
        N=
        \begin{mat}[r]
          0  &1  \\
          0  &0
        \end{mat}
      \end{equation*}
      Polinomial karakteristiknya adalah $c(x)=x^2$.
      Jadi, polinomial minimalnya adalah $m_1(x)=x$ atau $m_2(x)=x^2$.
      Kemungkinan pertama tidak tepat karena $N-0\cdot I$ bukan matriks nol;
      dengan demikian, dalam contoh ini polinomial minimal berderajat sama dengan
      dimensi ruang dasarnya, dan seperti telah disebutkan,
      kita mengetahui bahwa matriks ini tidak dapat didiagonalkan karena nilpoten.
    \end{answer}
   \item 
     Misalkan $f(x)$ adalah sebuah polinomial.
     Buktikan bahwa jika $A$ dan $B$ merupakan matriks serupa, maka $f(A)$
     serupa dengan $f(B)$.
     \begin{exparts}
       \partsitem Selanjutnya tunjukkan bahwa matriks-matriks serupa memiliki
         polinomial karakteristik yang sama.
       \partsitem Tunjukkan bahwa matriks-matriks serupa memiliki polinomial minimal yang sama.
       \partsitem Tentukan apakah kedua matriks berikut serupa.
          \begin{equation*}
            \begin{mat}[r]
              1  &3  \\
              2  &3
            \end{mat}
            \qquad
            \begin{mat}[r]
              4  &-1 \\
              1  &1
            \end{mat}
          \end{equation*}
     \end{exparts}
     \begin{answer}
       Misalkan \( A \) dan \( B \) serupa, \( A=PBP^{-1} \).
       Dari fakta bahwa
       \begin{multline*}
           A^n=(PBP^{-1})^n=(PBP^{-1})(PBP^{-1})\cdots(PBP^{-1})   \\
                           =PB(P^{-1}P)B(P^{-1}P)\cdots (P^{-1}P)BP^{-1}
                           =PB^nP^{-1}
       \end{multline*}
       dan $c\cdot A=c\cdot(PBP^{-1})=P(c\cdot B)P^{-1}$, diperoleh
       fakta yang diperlukan bahwa untuk setiap fungsi polinomial $f$, kita mempunyai
       \( f(A)=P\,f(B)\,P^{-1} \).
       Sebagai contoh, jika $f(x)=x^2+2x+3$, maka
       \begin{multline*}
         A^2+2A+3I=(PBP^{-1})^2+2\cdot PBP^{-1}+3\cdot I           \\
                  =(PBP^{-1})(PBP^{-1})+P(2B)P^{-1}+3\cdot PP^{-1}
                  =P(B^2+2B+3I)P^{-1}
       \end{multline*}
       menunjukkan bahwa $f(A)$ serupa dengan $f(B)$.
       \begin{exparts}
         \partsitem Dengan mengambil $f$ sebagai polinomial linear, kita memperoleh bahwa
           $A-xI$ serupa dengan $B-xI$.
           Matriks-matriks serupa memiliki determinan yang sama (karena
           $\deter{A}=\deter{PBP^{-1}}
               =\deter{P}\cdot\deter{B}\cdot\deter{P^{-1}}
               =1\cdot\deter{B}\cdot 1=\deter{B}$).
           Jadi, polinomial karakteristiknya sama.
         \partsitem 
           Karena \( P \) dan \( P^{-1} \) invertibel, \( f(A) \) merupakan
           matriks nol jika dan hanya jika \( f(B) \) merupakan matriks nol.
         \partsitem Keduanya tidak mungkin serupa karena tidak memiliki
           polinomial karakteristik yang sama.
           Polinomial karakteristik matriks pertama adalah
           $x^2-4x-3$, sedangkan polinomial karakteristik matriks
           kedua adalah $x^2-5x+5$.
       \end{exparts}  
     \end{answer}
   \item 
    \begin{exparts}
     \partsitem Tunjukkan bahwa suatu matriks invertibel jika dan hanya jika
       suku konstan
       dalam polinomial minimalnya bukan $0$.
     \partsitem Tunjukkan bahwa jika matriks persegi \( T \) tidak invertibel,
       maka terdapat
       matriks tak nol \( S \) sedemikian sehingga \( ST \) dan \( TS \) keduanya
       sama dengan matriks nol.
     \end{exparts}
     \begin{answer}
      Misalkan \( m(x)=x^n+m_{n-1}x^{n-1}+\dots+m_1x+m_0 \)
      merupakan polinomial minimal untuk \( T \).
      \begin{exparts}
       \partsitem 
         Untuk argumen `jika',
         karena \( T^n+\dots+m_1T+m_0I \) adalah matriks nol, kita
         mempunyai \( I=(T^n+\dots+m_1T)/(-m_0)=
         T\cdot (T^{n-1}+\dots+m_1I)/(-m_0) \), sehingga
         matriks $(-1/m_0)\cdot (T^{n-1}+\dots+m_1I)$ merupakan invers $T$.
         Untuk argumen `hanya jika', andaikan \( m_0=0 \)
         (kita sisihkan kasus \( n=1 \), yang mudah), sehingga
         \( T^n+\dots+m_1T=(T^{n-1}+\dots+m_1I)T \) adalah matriks nol.
         Perhatikan bahwa \( T^{n-1}+\dots+m_1I \) bukan matriks nol karena
         derajat polinomial minimal adalah \( n \).
         Jika \( T^{-1} \) ada, maka mengalikan dari kanan baik
         \( (T^{n-1}+\dots+m_1I)T \) maupun matriks nol
         dengan $T^{-1}$ menghasilkan kontradiksi.
       \partsitem Jika \( T \) tidak invertibel, maka suku konstan dalam
         polinomial minimalnya adalah nol.
         Jadi,
         \begin{equation*}
            T^n+\dots+m_1T=(T^{n-1}+\dots+m_1I)T=T(T^{n-1}+\dots+m_1I)
         \end{equation*}
         adalah matriks nol.
       \end{exparts} 
     \end{answer}
   \recommended \item \label{exer:PolyMapsFactor}
      \begin{exparts}
        \partsitem Selesaikan pembuktian \nearbylemma{le:PolyMapsFactor}.
        \partsitem Berikan contoh yang menunjukkan bahwa hasil tersebut tidak berlaku
          jika $t$ tidak linear.
      \end{exparts}
      \begin{answer}
        \begin{exparts}
          \partsitem
            Untuk langkah induksi, andaikan \nearbylemma{le:PolyMapsFactor}
            benar untuk polinomial
            berderajat \( i,\ldots,k-1 \), dan perhatikan polinomial \( f(x) \)
            berderajat \( k \).
            Faktorkan $f(x)=k(x-\lambda_1)^{q_1}\cdots(x-\lambda_z)^{q_z}$,
            dan misalkan
            \( k(x-\lambda_1)^{q_1-1}\cdots(x-\lambda_z)^{q_z} \)
            adalah \( c_{n-1}x^{n-1}+\cdots+c_1x+c_0 \).
            Substitusikan:
            \begin{align*}
               k\composed{\composed{(t-\lambda_1)^{q_1}}{\cdots}}{
                                            (t-\lambda_z)^{q_z}}(\vec{v})
               &=
               \composed{(t-\lambda_1)}{
                  \composed{\composed{(t-\lambda_1)^{q_1}}{\cdots}}{
                  (t-\lambda_z)^{q_z}} }
                          (\vec{v})    \\
               &=
               (t-\lambda_1)\,(c_{n-1}t^{n-1}(\vec{v})+\cdots+c_0\vec{v}) \\
               &=
               f(t)(\vec{v})
            \end{align*}
           (kesamaan kedua diperoleh dari hipotesis induksi, dan
           kesamaan ketiga dari linearitas \( t \)).
         \partsitem Salah satu contohnya adalah pemetaan kuadrat
            $\map{s}{\Re}{\Re}$ yang diberikan oleh $s(x)=x^2$.
            Pemetaan ini tidak linear.
            Tindakan yang didefinisikan oleh polinomial $f(t)=t^2-1$
            mengubah $s$ menjadi $f(s)=s^2-1$, yaitu pemetaan berikut.
            \begin{equation*} 
              x\mapsunder{s^2-1} \composed{s}{s}(x)-1=x^4-1
            \end{equation*}
            Amati bahwa pemetaan ini berbeda dari pemetaan
            $\composed{(s-1)}{(s+1)}$; misalnya, pemetaan pertama mengirim
            $x=5$ ke $624$, sedangkan pemetaan kedua mengirim $x=5$ ke $675$.
       \end{exparts} 
     \end{answer}
%  \item
%    Give an example of two \( \nbyn{4} \) matrices that have the 
%    same characteristic polynomial and the same minimal polynomial but
%    nonetheless are not similar.
%    \begin{answer}
%      These two have the same characteristic polynomial,
%      $c_S(x)=c_T(x)=(x-2)^4$.
%      \begin{equation*}
%        S=
%        \begin{mat}[r]
%          2  &0  &0  &0  \\
%          1  &2  &0  &0  \\
%          0  &0  &2  &0  \\
%          0  &0  &0  &2
%        \end{mat}
%        \qquad
%        T=
%        \begin{mat}[r]
%          2  &0  &0  &0  \\
%          1  &2  &0  &0  \\
%          0  &0  &2  &0  \\
%          0  &0  &1  &2
%        \end{mat}
%      \end{equation*}
%      For each of the two
%      the candidates for the minimal polynomial are 
%      $m_1(x)=(x-2)$, $m_2(x)=(x-2)^2$, $m_3(x)=(x-2)^3$, and 
%      $m_4(x)=(x-2)^4$.
%      We have 
%      \begin{equation*}
%        S-2I=
%        \begin{mat}[r]
%          0  &0  &0  &0  \\
%          1  &0  &0  &0  \\
%          0  &0  &0  &0  \\
%          0  &0  &0  &0
%        \end{mat}
%        \qquad
%        T-2I=
%        \begin{mat}[r]
%          0  &0  &0  &0  \\
%          1  &0  &0  &0  \\
%          0  &0  &0  &0  \\
%          0  &0  &1  &0
%        \end{mat}
%      \end{equation*}
%      and
%      \begin{equation*}
%        (S-2I)^2=
%        \begin{mat}[r]
%          0  &0  &0  &0  \\
%          0  &0  &0  &0  \\
%          0  &0  &0  &0  \\
%          0  &0  &0  &0
%        \end{mat}
%        =
%        (T-2I)^2
%      \end{equation*}
%      and so the minimal polynomial for each is $m_2$.
%    \end{answer}
  \item
    Setiap transformasi atau matriks persegi memiliki polinomial minimal.
    Apakah konversnya berlaku?
    \begin{answer}
      Ya.
      Lakukan ekspansi sepanjang kolom terakhir untuk memeriksa bahwa
      \( x^n+m_{n-1}x^{n-1}+\dots+m_1x+m_0 \) sama dengan, mungkin hingga tanda,
      determinan matriks berikut.
      \begin{equation*}
         \begin{mat}
            -x  &0  &0  &      &    &m_0    \\
             0  &1-x&0  &      &    &m_1    \\
             0  &0  &1-x&      &    &m_2    \\
                &   &   &\ddots             \\
                &   &   &      &1-x &m_{n-1}
         \end{mat}
      \end{equation*} 
     \end{answer}
\end{exercises}












\subsectionoptional{Bentuk Kanonik Jordan}
Kita sedang mencari bentuk kanonik untuk keserupaan matriks.
Subbagian ini menuntaskan program tersebut dengan beralih
dari bentuk kanonik untuk kelas-kelas
matriks nilpoten ke
bentuk kanonik untuk semua kelas.

\begin{lemma} \label{le:NilIffOnlyEigenZero}
Suatu transformasi linear pada ruang vektor nontrivial
nilpoten jika dan hanya jika satu-satunya nilai eigennya adalah nol.
\end{lemma}

\begin{proof}
Jika transformasi linear $t$ pada ruang vektor nontrivial
nilpoten, maka terdapat
$n$ sedemikian sehingga $t^n$ merupakan pemetaan nol,
sehingga $t$ memenuhi polinomial $p(x)=x^n=(x-0)^n$.
Berdasarkan \nearbylemma{le:tSatisImpMinPolyDivides}, polinomial minimal
$t$ membagi $p$, sehingga satu-satunya akar polinomial minimal adalah nol.
Berdasarkan Cayley-Hamilton, \nearbytheorem{th:CayHam},
satu-satunya akar polinomial karakteristik adalah nol.
Jadi, satu-satunya nilai eigen $t$ adalah nol.

Sebaliknya, jika transformasi \( t \) pada ruang
berdimensi \( n \) hanya memiliki nilai eigen nol,
maka polinomial karakteristiknya adalah \( x^n \).
\nearbylemma{le:MatSatItsCharPoly} menyatakan bahwa suatu pemetaan memenuhi
polinomial karakteristiknya, sehingga \( t^n \) merupakan pemetaan nol.
Jadi, $t$ nilpoten.
\end{proof}

\noindent Frasa `ruang vektor nontrivial' terdapat dalam pernyataan
lema itu karena pada ruang trivial
$\set{\zero}$, satu-satunya transformasi adalah pemetaan nol, yang tidak memiliki
nilai eigen karena tidak ada vektor eigen tak nol yang terkait.

\begin{corollary} \label{cor:tMinLambdaNilpotent}
Transformasi $t-\lambda$ nilpoten jika dan hanya jika
satu-satunya nilai eigen $t$ adalah $\lambda$.
\end{corollary}

\begin{proof}
Transformasi \( t-\lambda \) nilpoten jika dan hanya jika
satu-satunya nilai eigen $t-\lambda$ adalah \( 0 \).
Hal ini berlaku jika dan hanya jika satu-satunya nilai eigen $t$ adalah $\lambda$, karena
\( t(\vec{v})=\lambda\vec{v} \,\) jika dan
hanya jika \( (t-\lambda)\,(\vec{v})=0\cdot\vec{v} \).
\end{proof}

\begin{lemma}   \label{le:SimRespAddScalar}
Jika matriks \( T-\lambda I \) dan \( N \) serupa,
maka \( T \) dan \( N+\lambda I \) juga serupa,
melalui matriks perubahan basis yang sama.
\end{lemma}

\begin{proof}
Dari \( N=P(T-\lambda I)P^{-1}=PTP^{-1}-P(\lambda I)P^{-1} \),
kita memperoleh $N=PTP^{-1}-PP^{-1}(\lambda I)$
karena matriks diagonal \( \lambda I \) berkomutasi dengan sebarang matriks.
Jadi, \( N=PTP^{-1}-\lambda I \),
dan dengan demikian \( N+\lambda I=PTP^{-1} \).
\end{proof}

Kita telah memiliki bentuk kanonik untuk
kasus matriks nilpoten,
yakni untuk kasus matriks yang satu-satunya nilai eigennya adalah nol.
Berdasarkan Korolari~III.\ref{cor:NilpotentMatCanonForm}, setiap
matriks semacam itu serupa dengan matriks yang semua entrinya
nol kecuali blok-blok satu pada subdiagonal.
% (To make this representation unique we can fix some arrangement of
% the blocks, perhaps from longest to shortest.)
Hasil sebelumnya memungkinkan kita memperluas ini ke semua
matriks yang memiliki satu nilai eigen~$\lambda$.
Dalam bentuk baru,
selain blok-blok satu pada subdiagonal, semua entri pada diagonal
bernilai~$\lambda$.
(Lihat \nearbydefinition{def:JordanBlock} di bawah.)

\begin{example}   \label{ex:SingJordBlock}
Polinomial karakteristik dari
\begin{equation*}
  T=\begin{mat}[r]
      2  &-1  \\
      1  &4
    \end{mat}
\end{equation*}
adalah \( (x-3)^2 \), sehingga \( T \) memiliki satu nilai eigen, yaitu \( 3 \).
% We can find a matrix similar to $T$ that is simple in 
% the way that our canonical form for nilpotent matrices is simple.
% The form is described in \nearbydefinition{def:JordanBlock} below but
% as a preview here the new matrix.
% \begin{equation*}
%   \begin{mat}[r]
%       3  &0  \\
%       1  &3
%     \end{mat}
% \end{equation*}

Karena $3$ adalah satu-satunya nilai eigen,
$T-3I$ hanya memiliki nilai eigen \( 0 \), sehingga nilpoten.
Untuk memperoleh bentuk kanonik~$T$, mulailah dengan
mencari basis rantai untuk $T-3I$
seperti pada bagian sebelumnya, dengan menghitung pangkat-pangkat matriks ini
dan ruang nol pangkat-pangkat tersebut.
\begin{center}
  \begin{tabular}{c|cc}
    \multicolumn{1}{c}{\textit{pangkat} \( p \)}  &\( (T-3I)^p \)  &\( \nullspace{(T-3I)^p}  \)   \\  
    \hline
    \( 1 \)
    &\(  
       \matrixvenlarge{\begin{mat}[r]
         -1  &-1 \\
         1   &1   \\
       \end{mat}}  \)
    &\( \set{\matrixvenlarge{\colvec{-y \\ y}}\suchthat
                               y\in\C}  \)   \\
    \( 2 \)
    &\(  
       \matrixvenlarge{\begin{mat}[r]
         0  &0  \\
         0  &0   \\
       \end{mat}}  \)
    &\( \set{\matrixvenlarge{\colvec{x \\ y}}\suchthat
                               x,y\in\C}  \)   \\
  \end{tabular}
\end{center}
Ruang nol $t-3$ berdimensi satu,
dan ruang nol $(t-3)^2$ berdimensi dua.
Jadi, berikut adalah matriks representasi kanonik dan bentuk
basis rantai~$B$ yang terkait untuk $t-3$.
\begin{equation*}
  \rep{t-3}{B,B}
  =N=
  \begin{mat}[r]
      0  &0   \\
      1  &0
  \end{mat}
  \qquad
  \begin{array}{r}
    \vec{\beta}_1\xmapsto{t-3}\vec{\beta}_2\xmapsto{t-3}\zero  \\
  \end{array}  
\end{equation*}
Carilah basis semacam itu
dengan terlebih dahulu memilih $\vec{\beta}_2$ yang dipetakan ke~$\zero$
oleh $t-3$ (yakni, pilih sebuah anggota
ruang nol $\nullspace{t-3}$ yang diberikan dalam tabel).
Kemudian pilih $\vec{\beta}_1$ yang dipetakan ke~$\zero$ oleh $(t-3)^2$.
Satu-satunya batasan pada pilihan tersebut, selain keanggotaan dalam ruang nol,
adalah bahwa $\vec{\beta}_1$ dan $\vec{\beta}_2$
harus membentuk himpunan bebas linear.
Berikut salah satu pilihan yang mungkin.
\begin{equation*}
  \vec{\beta}_1=\colvec[r]{1 \\ 1}
  \quad
  \vec{\beta}_2=\colvec[r]{-2 \\ 2}
\end{equation*}
Dari sini, \nearbylemma{le:SimRespAddScalar} menyatakan bahwa \( T \) serupa dengan
matriks berikut.
\begin{equation*}
  \rep{t}{B,B}=
  N+3I=
  \begin{mat}[r]
     3  &0  \\
     1  &3
  \end{mat}
\end{equation*}

Kita dapat memperlihatkan perhitungan keserupaannya.
Ambil
$T$ sebagai representasi transformasi $\map{t}{\C^2}{\C^2}$ terhadap
basis standar (seperti yang akan kita lakukan sepanjang sisa bab ini).
% Recall how to find the change of
% basis matrices $P$ and $P^{-1}$ to express \( N \) as \( P(T-3I)P^{-1} \).
Diagram keserupaan
\begin{equation*}
  \begin{CD}
    \C^2_{\wrt{\stdbasis_2}}      @>t-3>T-3I>      \C^2_{\wrt{\stdbasis_2}}     \\
    @V\scriptstyle\identity V\scriptstyle PV  
                                 @V\scriptstyle\identity V\scriptstyle PV \\
    \C^2_{\wrt{B}}                 @>t-3>N>         \C^2_{\wrt{B}}
  \end{CD}
\end{equation*}
menjelaskan bahwa untuk berpindah dari kiri bawah ke kiri atas, kita mengalikan dengan
\begin{equation*}
  P^{-1}=\bigl(\rep{\identity}{\stdbasis_2,B}\bigr)^{-1}
    =\rep{\identity}{B,\stdbasis_2}
    =\begin{mat}[r]
        1  &-2  \\
        1  &2
     \end{mat}
\end{equation*}
dan untuk berpindah dari kanan atas ke kanan bawah, kita mengalikan dengan
matriks berikut.
\begin{equation*}
  P=\begin{mat}[r]
      1  &-2  \\
      1  &2
     \end{mat}^{-1}\!\!
   =\begin{mat}[r]
      1/2  &1/2  \\
      -1/4 &1/4
   \end{mat}
\end{equation*}
Jadi, persamaan berikut memperlihatkan keserupaannya.
\begin{equation*}
  \begin{mat}[r]
      1/2  &1/2  \\
      -1/4 &1/4
   \end{mat}
   \begin{mat}[r]
      2  &-1  \\
      1  &4
    \end{mat}
    \begin{mat}[r]
        1  &-2  \\
        1  &2
     \end{mat}
  =
  \begin{mat}[r]
     3  &0  \\
     1  &3
  \end{mat}
\end{equation*}
\end{example}

\begin{example}  \label{ex:FullNilpotentCase}
Matriks berikut
\begin{equation*}
  T=
  \begin{mat}[r]
    4  &1  &0  &-1  \\
    0  &3  &0  &1   \\
    0  &0  &4  &0   \\
    1  &0  &0  &5
   \end{mat}
\end{equation*}
memiliki polinomial karakteristik \( (x-4)^4 \),
sehingga hanya memiliki nilai eigen~$4$.
\begin{center}
  \begin{tabular}{c|cc}
    \multicolumn{1}{c}{\textit{pangkat} \( p \)}  &\( (T-4I)^p \)  &\( \nullspace{(T-4I)^p}  \)   \\  
    \hline
    \( 1 \)
    &\(  
       \matrixvenlarge{\begin{mat}[r]
         0  &1  &0  &-1  \\
         0  &-1 &0  &1   \\
         0  &0  &0  &0   \\
         1  &0  &0  &1
       \end{mat}}  \)
    &\( \set{\matrixvenlarge{\colvec{-w \\ w \\ z \\ w}}\suchthat
                               z,w\in\C}  \)   \\
    \( 2 \)
    &\(  
       \matrixvenlarge{\begin{mat}[r]
         -1  &-1 &0  &0  \\
          1  &1  &0  &0   \\
          0  &0  &0  &0   \\
          1  &1  &0  &0
       \end{mat}}  \)
    &\( \set{\matrixvenlarge{\colvec{-y \\ y \\ z \\ w}}\suchthat
                               y,z,w\in\C}  \)   \\
    \( 3 \)
    &\(  
       \matrixvenlarge{\begin{mat}[r]
          0  &0   &0  &0  \\
          0  &0  &0  &0   \\
          0  &0  &0  &0   \\
          0  &0  &0  &0   \\
       \end{mat}}  \)
    &\( \set{\matrixvenlarge{\colvec{x \\ y \\ z \\ w}}\suchthat
                               x,y,z,w\in\C}  \)   \\
  \end{tabular}
% sage: T=matrix(QQ, [[4,1,0,-1], [0,3,0,1], [0,0,4,0], [1,0,0,5]])
% sage: S=T-4*identity_matrix(4)
% sage: S^2
% [-1 -1  0  0]
% [ 1  1  0  0]
% [ 0  0  0  0]
% [ 1  1  0  0]
% sage: S^3
% [0 0 0 0]
% [0 0 0 0]
% [0 0 0 0]
% [0 0 0 0]
% sage: S
% [ 0  1  0 -1]
% [ 0 -1  0  1]
% [ 0  0  0  0]
% [ 1  0  0  1]
% sage: S*vector(QQ, [-1,1, 0, 1])
% (0, 0, 0, 0)
\end{center}
Ruang nol $t-4$ berdimensi dua, ruang nol $(t-4)^2$
berdimensi tiga, dan ruang nol $(t-4)^3$ berdimensi empat.
Ini memberikan bentuk kanonik untuk $t-4$.
\begin{equation*}
  N=\rep{t-4}{B,B}=
  \begin{mat}[r]
    0  &0  &0  &0   \\
    1  &0  &0  &0   \\
    0  &1  &0  &0   \\
    0  &0  &0  &0
   \end{mat}
   \qquad
  \begin{array}{r}
    \vec{\beta}_{1}\xmapsto{t-4}\vec{\beta}_{2}
       \xmapsto{t-4}\vec{\beta}_{3}\xmapsto{t-4}\zero \\
   \vec{\beta}_{4}\xmapsto{t-4}\zero \\
  \end{array}  
\end{equation*}
Untuk menghasilkan basis semacam itu, pilih
$\vec{\beta}_3$ dan~$\vec{\beta}_4$ yang dipetakan ke~$\zero$
oleh~$t-4$.
Pilih pula~$\vec{\beta}_2$ yang dipetakan ke~$\zero$ oleh~$(t-4)^2$
dan $\vec{\beta}_1$ yang dipetakan ke~$\zero$ oleh~$(t-4)^3$.
\begin{equation*}
  \vec{\beta}_1=\colvec{1 \\ 0 \\ 0 \\ 0}
  \quad
  \vec{\beta}_2=\colvec{-1 \\ 1 \\ 0 \\ 0}
  \quad
  \vec{\beta}_3=\colvec{-1 \\ 1 \\ 0 \\ 1}
  \quad
  \vec{\beta}_4=\colvec{0 \\ 0 \\ 1 \\ 0}
\end{equation*}
Perhatikan bahwa keempat vektor ini dipilih agar tidak memiliki kebergantungan linear.
Berikut matriks bentuk kanonik yang serupa dengan~$T$.
\begin{equation*}
  \rep{t}{B,B}
  =N+4I=
  \begin{mat}[r]
    4  &0  &0  &0   \\
    1  &4  &0  &0   \\
    0  &1  &4  &0   \\
    0  &0  &0  &4
   \end{mat}
\end{equation*}
\end{example}

\begin{definition}  \label{def:JordanBlock}
Suatu \definend{blok Jordan}\index{blok Jordan}
adalah matriks persegi, atau blok persegi entri-entri di dalam suatu matriks,
yang seluruh entrinya nol kecuali bahwa
terdapat bilangan $\lambda\in\C$ sedemikian sehingga
setiap entri diagonal adalah $\lambda$,
dan setiap entri subdiagonal adalah~$1$.
(Jika bloknya berukuran $\nbyn{1}$, maka blok itu tidak memiliki entri subdiagonal~$1$.)
\end{definition}

Rantai-rantai dalam basis terkait disebut
\definend{untaian Jordan}\index{untaian Jordan}\index{basis!untaian Jordan}
atau \definend{rantai Jordan}.\index{rantai Jordan}\index{basis!rantai Jordan}

Contoh-contoh di atas memperlihatkan bahwa untuk matriks
dengan satu nilai eigen, matriks-matriks blok Jordan merupakan
wakil kanonik kelas-kelas keserupaan.
Kita dapat membuat bentuk matriks ini unik dengan
menyusun elemen basis sehingga blok-blok satu pada subdiagonal berurutan dari
yang terpanjang ke yang terpendek, dibaca dari kiri ke kanan.

\begin{example}
Matriks-matriks \( \nbyn{3} \) yang satu-satunya nilai eigennya adalah \( 1/2 \)
terbagi menjadi tiga kelas keserupaan.
Ketiga kelas itu memiliki wakil kanonik berikut.
\begin{equation*}
  \begin{mat}[r]
     1/2  &0    &0  \\
     0    &1/2  &0  \\
     0    &0    &1/2
   \end{mat}
   \qquad 
   \begin{mat}[r]
     1/2  &0    &0  \\
     1    &1/2  &0  \\
     0    &0    &1/2
   \end{mat}
   \qquad 
   \begin{mat}[r]
     1/2  &0    &0  \\
     1    &1/2  &0  \\
     0    &1    &1/2
   \end{mat}
\end{equation*}
Secara khusus, matriks berikut
\begin{equation*}
   \begin{mat}[r]
     1/2  &0    &0    \\
     0    &1/2  &0    \\
     0    &1    &1/2
   \end{mat}
\end{equation*}
termasuk dalam kelas keserupaan yang diwakili oleh matriks tengah, karena
konvensi pengurutan blok-blok satu pada subdiagonal.
\end{example}

Sekarang kita siap menuntaskan program bab ini.
Pertama, kita meninjau apa yang sejauh ini telah kita ketahui tentang
transformasi $\map{t}{\C^n}{\C^n}$.

Kasus yang dapat didiagonalkan terjadi ketika
$t$ memiliki~$n$ nilai eigen berbeda
$\lambda_1$, \ldots, $\lambda_n$, yakni ketika banyaknya nilai eigen
sama dengan dimensi ruang.
Dalam kasus ini terdapat basis
$\sequence{\vec{\beta}_1,\ldots,\vec{\beta}_n}$ dengan
$\vec{\beta}_i$ merupakan vektor eigen yang terkait dengan nilai eigen~$\lambda_i$.
Contoh di bawah memiliki
$n=3$, dan ketiga nilai eigennya adalah $\lambda_1=1$, $\lambda_2=3$,
dan $\lambda_3=-1$.
Contoh tersebut memperlihatkan matriks wakil kanonik dan basis terkait.
\begin{equation*}
  T_1=\begin{mat}
    1  &0  &0  \\
    0  &3  &0  \\
    0  &0  &-1
  \end{mat}
  \qquad
    \begin{array}{r}
    \vec{\beta}_1\mapsunder{t-1}\zero  \\
    \vec{\beta}_2\mapsunder{t-3}\zero  \\
    \vec{\beta}_3\mapsunder{t+1}\zero
    \end{array}  
\end{equation*}
Salah satu contoh diagonalisasi adalah
Lima.II.\ref{ex:SummaryOfDiagonalizable}.

Kasus ketika $t$ memiliki satu nilai eigen memanfaatkan
hasil-hasil tentang nilpotensi untuk memperoleh basis
vektor eigen terkait yang membentuk rantai-rantai saling lepas.
Contoh ini memiliki $n=10$, satu nilai eigen~$\lambda=2$,
dan lima blok Jordan.
\begin{equation*}
  T_2=\begin{mat}
    2  &0  &0  &0 &0 &0 &0 &0 &0 &0 \\
    1  &2  &0  &0 &0 &0 &0 &0 &0 &0 \\
    0  &1  &2  &0 &0 &0 &0 &0 &0 &0 \\
    0  &0  &0  &2 &0 &0 &0 &0 &0 &0 \\
    0  &0  &0  &1 &2 &0 &0 &0 &0 &0 \\
    0  &0  &0  &0 &1 &2 &0 &0 &0 &0 \\
    0  &0  &0  &0 &0 &0 &2 &0 &0 &0 \\
    0  &0  &0  &0 &0 &0 &1 &2 &0 &0 \\
    0  &0  &0  &0 &0 &0 &0 &0 &2 &0 \\
    0  &0  &0  &0 &0 &0 &0 &0 &0 &2 \\
  \end{mat}
  \qquad
  {\renewcommand{\arraystretch}{1.75}
    \begin{array}{r}
    \vec{\beta}_1\mapsunder{t-2}\vec{\beta}_2\mapsunder{t-2}\vec{\beta}_3\mapsunder{t-2}\zero  \\
    \vec{\beta}_4\mapsunder{t-2}\vec{\beta}_5\mapsunder{t-2}\vec{\beta}_6\mapsunder{t-2}\zero  \\
    \vec{\beta}_7\mapsunder{t-2}\vec{\beta}_8\mapsunder{t-2}\zero  \\
    \vec{\beta}_9\mapsunder{t-2}\zero  \\
    \vec{\beta}_{10}\mapsunder{t-2}\zero
    \end{array}}  
\end{equation*}
Kita telah melihat contoh lengkap dalam
\nearbyexample{ex:FullNilpotentCase}.

\nearbytheorem{th:JordanForm} di bawah memperluas kedua kasus ini
ke sebarang transformasi
$\map{t}{\C^n}{\C^n}$.
Bentuk kanoniknya terdiri atas
blok-blok Jordan yang memuat nilai eigen.
Berikut ilustrasi matriks semacam itu untuk $n=6$ dengan
nilai eigen $3$, $2$, dan~$-1$
(contoh-contoh dengan perhitungan lengkap diberikan setelah teorema).
\begin{equation*}
  T_3=\begin{mat}
    3  &0  &0  &0 &0 &0 \\
    1  &3  &0  &0 &0 &0  \\
    0  &0  &3  &0 &0 &0  \\
    0  &0  &0  &2 &0 &0  \\
    0  &0  &0  &1 &2 &0  \\
    0  &0  &0  &0 &0 &-1  \\
  \end{mat}
  \qquad
  {\renewcommand{\arraystretch}{1.5}
    \begin{array}{r}
    \vec{\beta}_1\mapsunder{t-3}\vec{\beta}_2\mapsunder{t-3}\zero  \\
    \vec{\beta}_3\mapsunder{t-3}\zero  \\
    \vec{\beta}_4\mapsunder{t-2}\vec{\beta}_5\mapsunder{t-2}\zero  \\
    \vec{\beta}_6\mapsunder{t+1}\zero  \\
    \end{array}}  
\end{equation*}
Matriks tersebut memiliki empat blok, dua
terkait dengan nilai eigen~$3$, serta masing-masing satu dengan \( 2 \) dan~\( -1 \).

 
\begin{theorem}  \label{th:JordanForm}
\index{serupa!bentuk kanonik}%
\index{bentuk kanonik!untuk keserupaan}\index{transformasi!bentuk Jordan untuk}
Setiap transformasi $\map{t}{\C^n}{\C^n}$ dapat direpresentasikan dalam
\definend{bentuk Jordan},\index{bentuk Jordan!merepresentasikan kelas keserupaan}\index{bentuk Jordan!definisi}
dengan setiap $J_{\lambda}$ merupakan blok Jordan.\index{blok Jordan}
\begin{equation*}
  \begin{mat}
    J_{\lambda_1}  &            &\text{\textit{--nol--}}                 \\
               &J_{\lambda_2}                                              \\
               &     &\ddots                                     \\
             %    &     &                           &J_{\lambda_{k-1}} &     \\
               &     &\text{\textit{--nol--}} &            &J_{\lambda_{k}}
  \end{mat}
\end{equation*}
Dengan kata lain, setiap matriks $\nbyn{n}$ serupa dengan matriks berbentuk ini.
\end{theorem} 

Dalam contoh sebelumnya, jika kita menerapkan
$t-3$ pada basis, maka ia mengirim dua elemen,
$\vec{\beta}_2$ dan~$\vec{\beta}_3$, ke~$\zero$.
Ini menyarankan pembuktian dengan induksi
karena dimensi daerah hasil~$\rangespace{t-3}$
lebih kecil daripada dimensi ruang semula.

\begin{proof}
% This proof is closely related to that of
% Five.III.\ref{th:NilMapHasStrBas}, that any nilpotent transformation
% has a basis of disjoint $t$-strings.
%
Kita akan menggunakan induksi pada~$n$.
Kasus dasar $n=1$ bersifat trivial karena setiap matriks $\nbyn{1}$ sudah berbentuk Jordan.
Untuk langkah induksi, andaikan bahwa
setiap transformasi pada $\C^1$, \ldots, $\C^{n-1}$ memiliki representasi
bentuk Jordan,
dan tetapkan $\map{t}{\C^n}{\C^n}$.

Setiap transformasi memiliki setidaknya satu nilai eigen~$\lambda_1$ karena
ia memiliki polinomial minimal, yang
memiliki setidaknya satu akar atas bilangan kompleks.
Nilai eigen tersebut memiliki vektor eigen
tak nol~$\vec{v}$, sehingga $(t-\lambda_1)(\vec{v})=\zero$.
% and so the null space of $t-\lambda_1$ is nontrivial.
Jadi, dimensi $\nullspace{t-\lambda_1}$, yaitu nulitas pemetaan,
lebih besar dari nol.
Karena jumlah peringkat dan nulitas sama dengan dimensi domain,
peringkat $t-\lambda_1$
lebih kecil secara ketat daripada~$n$.
Tuliskan $W$ untuk daerah hasil $\rangespace{t-\lambda_1}$ dan
$r$ untuk peringkat, yaitu dimensinya.

Jika
$\vec{w}\in W=\rangespace{t-\lambda_1}$, maka $\vec{w}\in\C^n$, sehingga
$(t-\lambda_1)(\vec{w})\in\rangespace{t-\lambda_1}$.
Jadi, restriksi $t-\lambda_1$ ke~$W$ menginduksi suatu transformasi pada~$W$,
yang juga akan kita sebut $t-\lambda_1$.
Dimensi~$W$ lebih kecil daripada~$n$,
sehingga hipotesis induksi berlaku dan
kita memperoleh basis untuk~$W$ yang terdiri atas rantai-rantai saling lepas,
yang terkait dengan nilai-nilai eigen $t-\lambda_1$
pada~$W$.
Untuk setiap $\lambda$ yang merupakan nilai eigen~$t-\lambda_1$,
transformasi pada elemen-elemen basis adalah
$(t-\lambda_1)-\lambda=t-(\lambda_1+\lambda)$,
yang akan kita tulis sebagai $t-\lambda_s$.
Diagram berikut memperlihatkan beberapa rantai untuk $t-\lambda_1$, yakni
ketika $\lambda=0$;
argumen ini tetap berlaku jika tidak ada rantai semacam itu.
\begin{equation*}
\begin{strings}{*{10}{l}}
  \vec{w}_{1,n_1} &\xmapsto{t-\lambda_1} &\cdots 
     &\xmapsto{t-\lambda_1} &\vec{w}_{1,1} &\xmapsto{t-\lambda_1} &\zero \\
     \vdotswithin{\xmapsto{t-\lambda_1}}                   \\
  \vec{w}_{q,n_q} &\xmapsto{t-\lambda_1} &\cdots 
     &\xmapsto{t-\lambda_1} &\vec{w}_{q,1} &\xmapsto{t-\lambda_1} &\zero \\
  \vec{w}_{q+1,n_{q+1}} &\xmapsto{t-\lambda_2} &\cdots 
     &\xmapsto{t-\lambda_2} &\vec{w}_{q+1,1} &\xmapsto{t-\lambda_2} &\zero \\
     \vdotswithin{\xmapsto{t-\lambda_1}}                              \\
  \vec{w}_{k,n_k} &\xmapsto{t-\lambda_i} &\cdots 
     &\xmapsto{t-\lambda_i} &\vec{w}_{k,1} &\xmapsto{t-\lambda_i} &\zero \\
\end{strings}
\end{equation*}
Rantai-rantai ini mungkin memiliki panjang berbeda;
misalnya, mungkin $n_1\neq n_k$.

Vektor-vektor tak-$\zero$ paling kanan,
$\vec{w}_{1,1}$, \ldots, 
$\vec{w}_{k,1}$
berada dalam ruang nol pemetaan yang terkait,
$(t-\lambda_1)(\vec{w}_{1,1})=\zero$, \ldots,
$(t-\lambda_i)(\vec{w}_{k,1})=\zero$.
Jadi, banyaknya rantai yang terkait dengan $t-\lambda_1$,
yang dinotasikan $q$ dalam diagram di atas,
sama dengan dimensi $W\cap\nullspace{t-\lambda_1}$.

Basis rantai di atas merupakan basis ruang $W=\rangespace{t-\lambda_1}$.
Sekarang kita memperluasnya menjadi basis seluruh~$\C^n$.
Pertama, karena
setiap vektor $\vec{w}_{1,n_1}$, \ldots,~$\vec{w}_{q,n_q}$
merupakan anggota~$\rangespace{t-\lambda_1}$, masing-masing adalah citra
suatu vektor dari~$\C^n$.
Tambahkan satu vektor semacam itu di depan setiap rantai,
\( \vec{x}_1 \), \ldots, \( \vec{x}_q \),
seperti diperlihatkan di bawah.
Kedua, dimensi $W\cap\nullspace{t-\lambda_1}$ adalah~$q$,
sedangkan nulitasnya adalah~$n-r$,
sehingga pencacahannya mengharuskan adanya subruang
$Y\subseteq\nullspace{t-\lambda_1}$ berdimensi $n-r-q$
yang irisannya dengan $W$ adalah~$\set{\zero}$.
Karena itu,
pilih $n-r-q$ vektor bebas linear
$\vec{y}_1$, \ldots, $\vec{y}_{n-r-q}\in\nullspace{t-\lambda_1}$,
dan masukkan semuanya ke daftar rantai,
sekali lagi seperti diperlihatkan di bawah.
\begin{equation*}
\begin{strings}{*{10}{c}}
  \vec{x}_1  &\xmapsto{t-\lambda_1} 
     &\vec{w}_{1,n_1} &\xmapsto{t-\lambda_1} &\cdots 
     &\xmapsto{t-\lambda_1} &\vec{w}_{1,1} &\xmapsto{t-\lambda_1} &\zero \\
     \vdotswithin{\xmapsto{t-\lambda_1}}                   \\
  \vec{x}_q  &\xmapsto{t-\lambda_1}  
     &\vec{w}_{q,n_q} &\xmapsto{t-\lambda_1} &\cdots 
     &\xmapsto{t-\lambda_1} &\vec{w}_{q,1} &\xmapsto{t-\lambda_1} &\zero \\
  &&\vec{w}_{q+1,n_{q+1}} &\xmapsto{t-\lambda_2} &\cdots 
     &\xmapsto{t-\lambda_2} &\vec{w}_{q+1,1} &\xmapsto{t-\lambda_2} &\zero \\
     &&\vdotswithin{\xmapsto{t-\lambda_1}}                              \\
  &&\vec{w}_{k,n_k} &\xmapsto{t-\lambda_i} &\cdots 
     &\xmapsto{t-\lambda_i} &\vec{w}_{k,1} &\xmapsto{t-\lambda_i} &\zero \\
     &&&&&&\vec{y}_1  &\xmapsto{t-\lambda_1} &\zero \\
     &&&&&&\vdotswithin{\xmapsto{t-\lambda_1}}                              \\
     &&&&&&\vec{y}_{n-r-q}  &\xmapsto{t-\lambda_1} &\zero \\
\end{strings}
\end{equation*}
Kita akan menunjukkan bahwa inilah basis yang diinginkan untuk~$\C^n$.

Karena pencacahan tersebut, banyaknya vektor dalam himpunan
sama dengan dimensi ruang, sehingga
cukup kita verifikasikan bahwa elemen-elemennya bebas linear.
Andaikan bentuk berikut sama dengan~$\zero$.
\begin{equation*}
  a_1\vec{x}_1+\cdots+a_q\vec{x}_q   
    +b_{1,n_1}\vec{w}_{1,n_1}+\cdots+b_{k,1}\vec{w}_{k,1} 
    +c_1\vec{y}_1+\cdots+c_{n-r-q}\vec{y}_{n-r-q}
  \tag{$*$}
\end{equation*}
Pertama-tama kita menunjukkan bahwa sebagai akibatnya,
semua $a$ bernilai nol.
Terapkan $t-\lambda_1$ pada~($*$).
Vektor-vektor $\vec{y}_1$, \ldots, $\vec{y}_{n-r-q}$
dipetakan ke~$\zero$.
Setiap vektor $\vec{x}_i$ dikirim ke
$\vec{w}_{i,n_i}$.
Adapun vektor-vektor $\vec{w}$, semuanya memiliki sifat
$(t-\lambda_k)\,\vec{w}_{i,j}=\vec{w}_{i,j-1}$ 
untuk suatu~$\lambda_k$.
Tuliskan $\hat{\lambda}$ untuk $\lambda_k-\lambda_1$ agar diperoleh
\begin{equation*}
  (t-\lambda_1)(\vec{w}_{i,j})
    =(t-(\lambda_1+\hat{\lambda})+\hat{\lambda})(\vec{w}_{i,j})
    =\vec{w}_{i,j-1}+\hat{\lambda}\cdot\vec{w}_{i,j} 
  \tag{$**$}
\end{equation*}
Jadi, penerapan $t-\lambda_1$ pada~($*$) menghasilkan
kombinasi linear vektor-vektor $\vec{w}$.
Vektor-vektor ini bebas linear karena
membentuk basis~$W$, sehingga dalam kombinasi linear ini
semua koefisiennya bernilai~$0$.
Untuk menunjukkan bahwa $a_i=0$ bagi $i=1,\ldots,q$, cukup
kita verifikasikan bahwa setelah menerapkan $t-\lambda_1$ pada~($*$),
koefisien $\vec{w}_{i,n_i}$ adalah $a_i$.
Hal itu benar karena untuk $i=1,\ldots,q$,
persamaan~($**$) memiliki $\hat{\lambda}=0$,
sehingga $\vec{w}_{i,n_i}$ merupakan citra $x_{i}$ saja.
% Thus $a_1$, \ldots, $a_q$ are zero. 

Dengan semua $a$ sama dengan nol, yang tersisa dari~($*$) adalah
\begin{equation*}
  \zero
  =
    b_{1,n_1}\vec{w}_{1,n_1}+\cdots+b_{k,1}\vec{w}_{k,1}
    +c_1\vec{y}_1+\cdots+c_{n-r-q}\vec{y}_{n-r-q}
\end{equation*}
Tuliskan kembali sebagai $-(b_{1,n_1}\vec{w}_{1,n_1}+\cdots+b_{k,1}\vec{w}_{k,1})=
    c_1\vec{y}_1+\cdots+c_{n-r-q}\vec{y}_{n-r-q}$.
Vektor-vektor $\vec{w}$ berasal dari~$W$, sedangkan vektor-vektor $\vec{y}$ berasal dari~$Y$,
dan kedua ruang itu hanya memiliki irisan trivial.
Karena itu, $b_{1,n_1}\vec{w}_{1,n_1}+\cdots+b_{k,1}\vec{w}_{k,1}=\zero$
dan $c_1\vec{y}_1+\cdots+c_{n-r-q}\vec{y}_{n-r-q}=\zero$,
serta semua $b$ dan $c$ bernilai nol.
\end{proof}




% Recall that Lemma~III.\ref{le:RangeAndNullChains} shows that for any
% transformation~$t$ on an $n$~dimensional space
% the range spaces of iterates are stable 
% \begin{equation*}
%   \rangespace{t^n}=\rangespace{t^{n+1}}=\cdots=\genrangespace{t}
% \end{equation*} 
% as are the null spaces. 
% \begin{equation*}
%   \nullspace{t^n}=\nullspace{t^{n+1}}=\cdots=\gennullspace{t}
% \end{equation*} 
% Thus,
% the generalized null space $\gennullspace{t}$ and the generalized range space
% $\genrangespace{t}$ are $t$~invariant.
% % For the generalized null space, if $\vec{v}\in\gennullspace{t}$ then
% % $t^n(\vec{v})=\zero$ where $n$ is the dimension of the underlying space
% % and so $t(\vec{v})\in\gennullspace{t}$ because
% % $t^n(\,t(\vec{v})\,)$ is zero also.
% % For the generalized rangespace, if $\vec{v}\in\genrangespace{t}$ then
% % $\vec{v}=t^n(\vec{w})$ for some $\vec{w}$ and then 
% % $t(\vec{v})=t^{n+1}(\vec{w})=t^n(\,t(\vec{w})\,)$ 
% % shows that $t(\vec{v})$ is also a member of $\genrangespace{t}$.
% Also, 
% $\gennullspace{t-\lambda_i}$ and $\genrangespace{t-\lambda_i}$
% are $t-\lambda_i$~invariant.

% The action of the transformation $t-\lambda_i$ on $\gennullspace{t-\lambda_i}$
% is especially easy to understand.
% Observe that any transformation~$t$ is nilpotent on $\gennullspace{t}$, 
% because if $\vec{v}\in\gennullspace{t}$ then 
% by definition $t^n(\vec{v})=\zero$. 
% Thus $t-\lambda_i$ is nilpotent on 
% $\gennullspace{t-\lambda_i}$.

% % Thus, the generalized null space $\gennullspace{t-\lambda_i}$ is a part of
% % the space on which the action of $t-\lambda_i$ is easy to understand.

% We shall take three steps to prove this section's major result. 
% The next result is the first.
% % It establishes that if $j=neq i$ then \( t-\lambda_j \) leaves
% % \( t-\lambda_i \)'s part unchanged.

% \begin{lemma} \label{le:tInvIfftMinLambdaInv}
% A subspace is \( t \) invariant if and only if 
% it is \( t-\lambda \) invariant for all scalars \( \lambda \).
% In particular, 
% if \( \lambda_i \) is an eigenvalue of  a linear transformation
% \( t \) then for any other eigenvalue $\lambda_j$
% the spaces \( \gennullspace{t-\lambda_i} \) 
% and \( \genrangespace{t-\lambda_i} \)
% are \( t-\lambda_j \) invariant.
% \end{lemma}

% \begin{proof}
% For the first sentence we check the two implications separately.
% The `if' half is easy: if the subspace is $t-\lambda$ invariant for 
% all scalars $\lambda$ then using $\lambda=0$ shows that it is $t$ invariant.
% For `only if' suppose that the subspace is $t$ invariant,
% so that if $\vec{m}\in M$ then $t(\vec{m})\in M$, and let $\lambda$ be 
% a scalar.
% The subspace $M$ is closed under linear combinations and so if 
% $t(\vec{m})\in M$ then $t(\vec{m})-\lambda\vec{m}\in M$.
% Thus if $\vec{m}\in M$ then $(t-\lambda)\,(\vec{m})\in M$.

% The lemma's second sentence follows from its first.
% The
% two spaces are $t-\lambda_i$~invariant so they are \( t \)~invariant.
% Apply the first sentence again to
% conclude that they are also \( t-\lambda_j \) invariant.
% \end{proof}

% The second step of the three that we will take to prove this
% section's major result makes use of an additional property of 
% \( \gennullspace{t-\lambda_i} \) and
% \( \genrangespace{t-\lambda_i} \), that they are complementary.
% Recall that if a space is the direct sum of two others 
% \( V=\mathscr{N}\directsum \mathscr{R} \) 
% then any vector \( \vec{v} \) in the space breaks into
% two parts \( \vec{v}=\vec{n}+\vec{r} \) where \( \vec{n}\in \mathscr{N} \) and
% \( \vec{r}\in \mathscr{R} \), and recall also 
% that if \( B_{\mathscr{N}} \) and \( B_{\mathscr{R}} \) are bases for
% \( \mathscr{N} \) and \( \mathscr{R} \) then the concatenation
% \( \cat{B_{\mathscr{N}}}{B_{\mathscr{R}}} \) is linearly independent.
% The next result says that for any subspaces
% \( \mathscr{N} \) and \( \mathscr{R} \) that are complementary 
% as well as \( t \)~invariant,
% the action
% of \( t \) on \( \vec{v} \) breaks into the actions of
% \( t \) on \( \vec{n} \) and on \( \vec{r} \).

% \begin{lemma} \label{le:InvCompSubspSplitTrans}
% Let \( \map{t}{V}{V} \) be a transformation and let \( \mathscr{N} \) and 
% \( \mathscr{R} \) be
% \( t \) invariant complementary subspaces of \( V \).
% Then we can represent \( t \) by a matrix with
% blocks of square submatrices $T_1$ and $T_2$
% \begin{equation*} \renewcommand{\arraystretch}{1.2}
%   \begin{pmat}{c|c}
%       T_1   &Z_2  \\  \cline{1-2}
%       Z_1 &T_2
%    \end{pmat}
%    \begin{array}{@{}l}
%      \} \text{\ $\dim(\mathscr{N})$-many rows}  \\
%      \} \text{\ $\dim(\mathscr{R})$-many rows}
%    \end{array}
% \end{equation*}
% where \( Z_1 \) and \( Z_2 \) are blocks of zeroes.
% \end{lemma}

% \begin{proof}
% Since the two subspaces are complementary, the concatenation of a basis
% for \( \mathscr{N} \) with a basis for \( \mathscr{R} \) makes a basis
% \( B=\sequence{\vec{\nu}_1,\dots,\vec{\nu}_p,
%         \vec{\mu}_1,\ldots,\vec{\mu}_q}  \)
% for \( V \).
% We shall show that the matrix
% \begin{equation*}
%   \rep{t}{B,B}=
%   \begin{pmat}{c@{\hspace*{1em}}c@{\hspace*{1em}}c}
%      \vdots                   &        &\vdots     \\
%      \rep{t(\vec{\nu}_1)}{B}  &\cdots  &\rep{t(\vec{\mu}_q)}{B}  \\
%      \vdots                   &        &\vdots     \\
%   \end{pmat}
% \end{equation*}
% has the desired form.

% Any vector \( \vec{v}\in V \) is a member of \( \mathscr{N} \) 
% if and only if when it is represented with respect to \( B \)
% the final \( q \)
% coefficients are zero.
% As \( \mathscr{N} \) is \( t \)~invariant, each of the vectors
% \( \rep{t(\vec{\nu}_1)}{B} \),
% \ldots, \( \rep{t(\vec{\nu}_p)}{B} \) has this form.
% Hence the lower left of \( \rep{t}{B,B} \) is all zeroes.
% The argument for the upper right is similar.
% \end{proof}

% To see that we have decomposed \( t \) into its action on the parts, 
% let $B_{\mathscr{N}}=\sequence{\vec{\nu}_1,\dots,\vec{\nu}_p}$ and 
% $B_{\mathscr{R}}=\sequence{\vec{\mu}_1,\ldots,\vec{\mu}_q}$.
% The restrictions of \( t \) to the subspaces \( \mathscr{N} \) 
% and~\( \mathscr{R} \) 
% are represented
% with respect to the bases $B_{\mathscr{N}},B_{\mathscr{N}}$
% and $B_{\mathscr{R}},B_{\mathscr{R}}$
% by the matrices \( T_1 \) and \( T_2 \).
% So with subspaces that are invariant and complementary 
% we can split the problem of examining
% a linear transformation into two lower-dimensional subproblems.
% The next result illustrates this decomposition into blocks.

% \begin{lemma} \label{le:DetIsProdOfSubDets}
% If $T$ is a matrix with square submatrices $T_1$ and $T_2$
% \begin{equation*} \renewcommand{\arraystretch}{1.2}
%   T=
%   \begin{pmat}{c|c}
%       T_1   &Z_2  \\  \cline{1-2}
%       Z_1   &T_2
%    \end{pmat}
% \end{equation*}
% where the \( Z \)'s are blocks of zeroes,
% then \( \deter{T}=\deter{T_1}\cdot\deter{T_2} \).
% \end{lemma}

% \begin{proof}
% Suppose that \( T \) is \( \nbyn{n} \),
% that \( T_1 \) is \( \nbyn{p} \),
% and that \( T_2 \) is \( \nbyn{q} \).
% In the permutation formula for the determinant
% \begin{equation*}
%   \deter{T}=
%   \sum_{\text{\scriptsize permutations\ }\phi}
%           t_{1,\phi(1)}t_{2,\phi(2)}\cdots t_{n,\phi(n)}\sgn(\phi)
% \end{equation*}
% each term comes from a rearrangement of the column numbers
% \( 1,\dots,n \) into a new order \( \phi(1),\dots,\phi(n) \).
% The upper right block $Z_2$ is all zeroes, so if a
% \( \phi \) has at least one of \( p+1,\dots,n \) among its first
% \( p \) column numbers \( \phi(1),\dots,\phi(p) \) then the term arising
% from \( \phi \) does not contribute to the sum because it is zero,
% e.g., if \( \phi(1)=n \) then
% \( t_{1,\phi(1)}t_{2,\phi(2)}\dots t_{n,\phi(n)}
%    =0\cdot t_{2,\phi(2)}\dots t_{n,\phi(n)}=0 \).

% So the above formula reduces to a sum over all permutations with 
% two halves:~any contributing $\phi$ is the composition of a $\phi_1$ that
% rearranges only \( 1,\dots,p \) 
% and a $\phi_2$ that rearranges only \( p+1,\dots,p+q \).
% Now, the distributive law 
% and the fact that the signum of a composition is the product
% of the signums gives that this
% \begin{multline*}
%    \deter{T_1}\cdot\deter{T_2}=
%    \bigg(\sum_{\begin{subarray}{c}
%                 \text{\scriptsize perms\ }\phi_1 \\
%                 \text{\scriptsize of\ } 1,\dots,p
%                \end{subarray}}
%        \!\!\! t_{1,\phi_1(1)}\cdots t_{p,\phi_1(p)}\sgn(\phi_1) \bigg)  \\
%    \cdot
%    \bigg(\sum_{\begin{subarray}{c}
%                 \text{\scriptsize perms\ }\phi_2 \\
%                 \text{\scriptsize of\ } p+1,\dots,p+q
%                \end{subarray}}
%        \!\!\! t_{p+1,\phi_2(p+1)}\cdots t_{p+q,\phi_2(p+q)}\sgn(\phi_2) 
%         \bigg)
% \end{multline*}
% equals
%   $\deter{T}=
%   \sum_{\text{\scriptsize contributing\ }\phi}
%           t_{1,\phi(1)}t_{2,\phi(2)}\cdots t_{n,\phi(n)}\sgn(\phi)$.
% \end{proof}

% \begin{example}
% \begin{equation*}
%     \begin{vmat}[r]
%        2  &0  &0  &0  \\
%        1  &2  &0  &0  \\
%        0  &0  &3  &0  \\
%        0  &0  &0  &3
%     \end{vmat}
%    =\begin{vmat}[r]
%        2  &0  \\
%        1  &2
%     \end{vmat}
%     \cdot
%     \begin{vmat}[r]
%        3  &0  \\
%        0  &3
%     \end{vmat}
%    =36
% \end{equation*}
% \end{example}

% From \nearbylemma{le:DetIsProdOfSubDets} we conclude that
% if two subspaces 
% are complementary and \( t \)~invariant then
% \( t \) is one-to-one if and only if its 
% restriction %\appendrefs{restrictions}
% to each subspace is nonsingular.

% Now for the promised third, and final, step to the main result.

% \begin{lemma}
% If a linear transformation \( \map{t}{V}{V} \) has the 
% characteristic polynomial
% \( (x-\lambda_1)^{p_1}\dots(x-\lambda_k)^{p_k} \) then
% (1)~\( V=\gennullspace{t-\lambda_1}\directsum\cdots
%              \directsum\gennullspace{t-\lambda_k} \) 
% and
% (2)~\( \dim(\gennullspace{t-\lambda_i})=p_i  \).
% \end{lemma}

% \begin{proof}
% This argument consists of proving
% two preliminary claims, followed by proofs of clauses~(1) and~(2).

% The first claim is that
% \( \gennullspace{t-\lambda_i}\intersection\gennullspace{t-\lambda_j}=\set{\zero} \)
% when \( i\neq j \).
% By \nearbylemma{le:tInvIfftMinLambdaInv}
% both \( \gennullspace{t-\lambda_i} \) and 
% \( \gennullspace{t-\lambda_j} \) are \( t \)~invariant.
% The intersection of \( t \) invariant subspaces is \( t \)
% invariant and so the restriction of \( t \) to
% \( \gennullspace{t-\lambda_i}\intersection\gennullspace{t-\lambda_j} \)
% is a linear transformation.
% Now,
% $t-\lambda_i$ is nilpotent on \( \gennullspace{t-\lambda_i} \)
% and  
% $t-\lambda_j$ is nilpotent on \( \gennullspace{t-\lambda_j} \),
% so both \( t-\lambda_i \) and \( t-\lambda_j \) are 
% nilpotent on the intersection.
% Therefore by \nearbylemma{le:NilIffOnlyEigenZero} and the observation
% following it, if \( t \) has any eigenvalues on the intersection 
% then the ``only'' eigenvalue is both
% \( \lambda_i \) and \( \lambda_j \).
% This cannot be, so the restriction has no eigenvalues:
% \( \gennullspace{t-\lambda_i}\intersection\gennullspace{t-\lambda_j} \)
% is the trivial space
% (\nearbylemma{le:MapNonTrivSpHasEigen} shows that 
% the only transformation that is without any eigenvalues is 
% the transformation on the trivial space).

% The second claim is that  when \( i\neq j \) then
% \( \gennullspace{t-\lambda_i} \subseteq \genrangespace{t-\lambda_j} \).
% To verify it we will show that \( t-\lambda_j \) is one-to-one on
% \( \gennullspace{t-\lambda_i} \).
% This will imply, 
% since \( \gennullspace{t-\lambda_i} \) is \( t-\lambda_j \) invariant
% by \nearbylemma{le:tInvIfftMinLambdaInv},
% that the map \( t-\lambda_j \) is an automorphism of the subspace
% \( \gennullspace{t-\lambda_i} \),
% and therefore that 
% \( \gennullspace{t-\lambda_i} \)  is a subset of each
% \( \rangespace{t-\lambda_j} \), \( \rangespace{(t-\lambda_j)^2} \), etc.
% To verify that the map is one-to-one,
% suppose that \( \vec{v}\in\gennullspace{t-\lambda_i} \) is in the 
% null space of \( t-\lambda_j \), aiming to show that \( \vec{v}=\zero \).
% Consider the map 
% \( [(t-\lambda_i)-(t-\lambda_j)]^n \).
% On the one hand, the only
% vector that \( (t-\lambda_i)-(t-\lambda_j)=\lambda_i-\lambda_j \) maps to 
% zero is the zero vector itself.
% On the other hand, as in the proof of 
% \nearbylemma{le:PolyMapsFactor}
% the binomial expansion gives this.
% \begin{equation*}
%   (t-\lambda_i)^n(\vec{v})
%     +\binom{n}{1}(t-\lambda_i)^{n-1}(t-\lambda_j)^1(\vec{v})
%     +\binom{n}{2}(t-\lambda_i)^{n-2}(t-\lambda_j)^2(\vec{v})
%     +\cdots 
% \end{equation*}
% The first term is the zero vector 
% because \( \vec{v}\in\gennullspace{t-\lambda_i} \).
% The remaining terms are the zero vector because
% \( \vec{v} \) is in the null space of \( t-\lambda_j \).
% Therefore \( \vec{v}=\zero \). 

% With those two preliminary claims done 
% we can prove clause~(1), that the space is the direct sum of the
% generalized null spaces.
% By Corollary~III.\ref{GenRngNullDirSumToSp} the space is the direct sum
% \( V=\gennullspace{t-\lambda_1}\directsum\genrangespace{t-\lambda_1} \).
% By the second claim
% \( \gennullspace{t-\lambda_2}\subseteq\genrangespace{t-\lambda_1} \)
% and so we can get a basis for \( \genrangespace{t-\lambda_1} \) by
% starting with a basis for \( \gennullspace{t-\lambda_2} \) and adding
% extra basis elements taken from
% \( \genrangespace{t-\lambda_1}\intersection\genrangespace{t-\lambda_2} \). 
% Thus \( V=\gennullspace{t-\lambda_1}\directsum\gennullspace{t-\lambda_2}
%           \directsum 
%        (\genrangespace{t-\lambda_1}\intersection\genrangespace{t-\lambda_2}) \).
% Continuing in this way we get this.
% \begin{equation*}
%   V=\gennullspace{t-\lambda_1}\directsum\cdots\directsum\gennullspace{t-\lambda_k}
%           \directsum 
%        (\genrangespace{t-\lambda_1}\intersection\cdots\intersection\genrangespace{t-\lambda_k})
% \end{equation*}
% The first claim above shows that the final space is trivial. 

% We finish by verifying clause~(2).
% Decompose \( V \) as 
% \( \gennullspace{t-\lambda_i}\directsum\genrangespace{t-\lambda_i} \)
% %with basis $B=\cat{B_{\mathscr{N}}}{B_{\mathscr{R}}}$
% and apply \nearbylemma{le:InvCompSubspSplitTrans}.
% \begin{equation*}   \renewcommand{\arraystretch}{1.2}
%   T=
% %  \rep{t}{B,B}=
%   \begin{pmat}{c|c}
%       T_1   &Z_2  \\  \cline{1-2}
%       Z_1   &T_2
%    \end{pmat}
%    \begin{array}{@{}l}
%      \} \text{\ $\dim(\,\gennullspace{t-\lambda_i}\,)$-many rows}  \\
%      \} \text{\ $\dim(\,\genrangespace{t-\lambda_i}\,)$-many rows}
%    \end{array}
% \end{equation*}
% \nearbylemma{le:DetIsProdOfSubDets} says that
% \( \deter{T-xI}=\deter{T_1-xI}\cdot\deter{T_2-xI} \).
% By the uniqueness clause of the Fundamental Theorem of Algebra, 
% Theorem~I.\ref{th:FundThmAlg},
% the determinants of the blocks have the same factors as the
% characteristic polynomial
% \( \deter{T_1-xI}=(x-\lambda_1)^{q_1}\cdots(x-\lambda_z)^{q_k} \)
% and
% \( \deter{T_2-xI}=(x-\lambda_1)^{r_1}\cdots(x-\lambda_z)^{r_k} \),
% where
% \( q_1+r_1=p_1 \), \dots, \( q_k+r_k=p_k \).
% We will finish by establishing that (i)~$q_j=0$ for all $j\neq i$,
% and (ii)~$q_i=p_i$.
% Together these prove clause~(2) because they show that the degree of 
% the polynomial $\deter{T_1-xI}$ is $q_i$ and
% the degree of that polynomial equals 
% the dimension of the
% generalized null space $\gennullspace{t-\lambda_i}$. 

% For (i),
% because the restriction of \( t-\lambda_i \) to \( \gennullspace{t-\lambda_i} \)
% is nilpotent on that space,
% $t$'s only eigenvalue on that space is \( \lambda_i \),
% by \nearbylemma{cor:tMinLambdaNilpotent}.
% So $q_j=0$ for $j\neq i$.

% For (ii),
% consider the restriction of \( t \) to \( \genrangespace{t-\lambda_i} \).
% By Lemma~III.\ref{lem:RestONeToOne}, the map
% \( t-\lambda_i \) is one-to-one on
% \( \genrangespace{t-\lambda_i} \) and so \( \lambda_i \) is not an
% eigenvalue of \( t \) on that subspace.
% Therefore \( x-\lambda_i \) is not a factor of \( \deter{T_2-xI} \),
% so $r_i=0$, and so \( q_i=p_i \).
% \end{proof}

% Recall the goal of this chapter, to give a canonical form for matrix similarity.
% That result is next.
% It 
% translates the above steps into matrix terms.

% \begin{theorem}
% \index{Jordan form!represents similarity classes}\index{similar!canonical form}
% \index{canonical form!for similarity}\index{transformation!Jordan form for}
% Any square matrix is similar to one in \definend{Jordan form}
% \begin{equation*}
%   \begin{mat}
%     J_{\lambda_1}  &            &\text{\textit{--zeroes--}}                 \\
%                &J_{\lambda_2}                                              \\
%                &     &\ddots                                     \\
%                 &     &                           &J_{\lambda_{k-1}} &     \\
%                &     &\text{\textit{--zeroes--}} &            &J_{\lambda_{k}}
%   \end{mat}
% \end{equation*}
% where each \( J_{\lambda} \) is the Jordan block associated with an
% eigenvalue $\lambda$ of the original matrix (that is, each \( J_{\lambda} \)
% is all zeroes except for
% \( \lambda \)'s down the diagonal and some subdiagonal ones).
% \end{theorem}

% \begin{proof}
% Given an \( \nbyn{n} \) matrix \( T \), consider the linear map
% \( \map{t}{\C^n}{\C^n} \) that it represents
% with respect to the standard bases.
% Use the prior lemma to write
% \( \C^n=\gennullspace{t-\lambda_1}\directsum\cdots
%         \directsum\gennullspace{t-\lambda_k} \)
% where \( \lambda_1 \), \ldots, \(\lambda_k \) are the eigenvalues of \( t \).
% Because each \( \gennullspace{t-\lambda_i} \)  is \( t \) invariant,
% \nearbylemma{le:InvCompSubspSplitTrans} and the prior lemma show
% that \( t \) is represented by a matrix that is all zeroes except for square
% blocks along the diagonal.
% To make those blocks into Jordan blocks, pick each \( B_{\lambda_i} \)
% to be a string basis for the action of \( t-\lambda_i \) on
% \( \gennullspace{t-\lambda_i} \). 
% \end{proof}

% \begin{corollary}
% Every square matrix is similar to the sum of a diagonal matrix and a nilpotent
% matrix.
% \end{corollary}

\begin{remark}
Secara teknis, agar menjadi bentuk kanonik untuk
keserupaan matriks,\index{wakil!kelas keserupaan}
bentuk Jordan harus unik.
Artinya, untuk setiap matriks persegi harus ada tepat satu
matriks berbentuk Jordan yang serupa dengannya.
Kita dapat membuat bentuk dalam teorema itu unik
dengan menyusun blok-blok Jordan sehingga nilai eigennya mengikuti
suatu urutan yang ditetapkan, lalu menyusun
blok-blok satu pada subdiagonal dari
yang terpanjang ke yang terpendek.
Dalam contoh-contoh di bawah, kita tidak akan membahas hal itu.
\end{remark}

% As in the nilpotent cases earlier in this section,
% to find the structure of the basis 
% associated with~$T_3$ we apply
% the maps $t-3$ and~$(t-3)^2$, the maps $t-2$ and~$(t-2)^2$, and 
% the map~$t+1$.
% For instance,
% the basis has a total of three elements that are mapped to~$\zero$
% when we apply $t-3$ twice.
% It has two elements
% that are sent to~$\zero$ on two applications of~$t-2$.
% And it has one element sent to~$\zero$ by~$t+1$.
% That is, to find the elements of this basis we want to
% iterate the maps, so we look in  
% the generalized null spaces $\gennullspace{t-3}$, $\gennullspace{t-2}$,
% and~$\gennullspace{t+1}$.

Contoh-contoh berikut dimulai dengan sebuah matriks dan menghitung
matriks berbentuk Jordan yang serupa dengannya.
Sebelumnya, kita akan mengemukakan satu hal lagi.
Sebagai contoh, pada contoh pertama di bawah kita memperoleh nilai eigen $2$
dan~$6$.
Untuk memilih elemen basis yang berkaitan dengan~\( 2\), kita mencari ruang nol
$t-2$, \( (t-2)^2 \), dan~\( (t-2)^3\).
Artinya, untuk mencari elemen-elemen basis dalam contoh ini, kita akan
melihat ruang nol tergeneralisasi $\gennullspace{t-2}$
dan~$\gennullspace{t-6}$.
Namun, kita perlu memastikan bahwa perhitungan untuk setiap pemetaan tidak
memengaruhi perhitungan untuk pemetaan lainnya.

% So before those examples, 
% suppose that $\map{t}{\C^n}{\C^n}$ is a transformation and 
% that $\lambda\in\C$ is an eigenvalue.
% Note that the generalized null space~$\gennullspace{t-\lambda}$ 
% and range space~$\genrangespace{t-\lambda}$ are $t-\lambda$
% invariant spaces.
% For the generalized null space, 
% if $\vec{v}\in\gennullspace{t-\lambda}$ then 
% $(t-\lambda)^n(\vec{v})=\zero$ and it follows that
% $(t-\lambda)^n(\,(t-\lambda)(\vec{v})\,)=(t-\lambda)^{n+1}(\vec{v})=\zero$ 
% also. 
% As to the generalized range space,
% Lemma Five.III.\ref{le:RangeAndNullChains} and its proof show that
% $\genrangespace{t-\lambda}
%   =\rangespace{(t-\lambda)^n}
%   =\rangespace{(t-\lambda)^{n+1}}$ 
% and so $\vec{v}\in\genrangespace{t-\lambda}$ implies that
% $(t-\lambda)(\vec{v})\in\genrangespace{t-\lambda}$ also.

% Before the theorem, we first will show that applying $t-\lambda$ 
% won't do something awkward to the basis elements not in 
% $\gennullspace{t-\lambda}$.
% For instance, in the prior example we want that applying
% $t-3$ won't do something awkward to~$\vec{\beta}_4$, 
% $\vec{\beta}_5$, and~$\vec{\beta}_6$. 
% Here, awkward means that
% where \( \map{t}{V}{V} \) is a linear transformation,
% the restriction\appendrefs{restrictions of functions}\spacefactor=1000  %
% of \( t \) to a subspace \( M\subseteq V \) 
% need not be a linear transformation on \( M \):~there
% may be an \( \vec{m}\in M \)
% with \( t(\vec{m})\not\in M \). 
% An example is where $t$ is the transformation 
% that rotates the plane by
% a quarter turn and $M$ is the $x=y$ line; the vector $\binom{1}{1}$
% is rotated away from the subspace.
% We next show this does not happen in the special case of 
% eigenvalues;~if $j\neq i$ then 
% application of \( t-\lambda_j \) leaves 
% \( t-\lambda_i \)'s part unchanged.

\begin{definition} \label{def:invariant}
Misalkan \( \map{t}{V}{V} \) adalah suatu transformasi.
Subruang \( M\subseteq V \) disebut \definend{invarian terhadap $t$}%
\index{subruang invarian}\index{subruang!invarian}
jika setiap kali \( \vec{m}\in M \), berlaku \( t(\vec{m})\in M \)
(secara ringkas: \( t(M)\subseteq M \)).
\end{definition}

\begin{lemma} \label{le:tInvIfftMinLambdaInv}
Suatu subruang invarian terhadap \( t \) jika dan hanya jika
subruang itu invarian terhadap \( t-\lambda \) untuk semua skalar \( \lambda \).
Secara khusus,
jika \( \lambda_i,\lambda_j\in\C \) merupakan nilai eigen berbeda dari
\( t \), maka
ruang \( \gennullspace{t-\lambda_i} \)
dan \( \genrangespace{t-\lambda_i} \)
invarian terhadap \( t-\lambda_j \).
\end{lemma}

\begin{proof}
Implikasi dari kanan ke kiri pada kalimat pertama bersifat trivial:~jika
subruang invarian terhadap $t-\lambda$ untuk
semua skalar $\lambda$, maka mengambil $\lambda=0$
menunjukkan bahwa subruang itu invarian terhadap $t$.
Untuk implikasi lainnya, andaikan subruang tersebut invarian terhadap $t$,
sehingga jika $\vec{m}\in M$, maka $t(\vec{m})\in M$, dan misalkan $\lambda$
suatu skalar.
Subruang $M$ tertutup terhadap kombinasi linear, sehingga jika
$t(\vec{m})\in M$, maka $t(\vec{m})-\lambda\vec{m}\in M$.
Jadi, jika $\vec{m}\in M$, maka $(t-\lambda)\,(\vec{m})\in M$.

Kalimat kedua mengikuti kalimat pertama.
Kedua ruang tersebut invarian terhadap $t-\lambda_i$, sehingga invarian terhadap \( t \).
Menerapkan kembali kalimat pertama menunjukkan bahwa
keduanya juga invarian terhadap \( t-\lambda_j \).
\end{proof}

\begin{example} \label{ex:FirstJordForm}
Matriks berikut
\begin{equation*}
   T=
   \begin{mat}[r]
     2  &0  &1  \\
     0  &6  &2  \\
     0  &0  &2
   \end{mat}
\end{equation*}
memiliki polinomial karakteristik \( -(x-2)^2(x-6) \).
\begin{equation*}
   |T-xI|=
   \begin{vmat}
     2-x  &0    &1  \\
     0    &6-x  &2  \\
     0    &0    &2-x
   \end{vmat}
   =-(x-2)^2(x-6)
\end{equation*}
Pertama-tama kita menangani nilai eigen~$2$.
Perhitungan pangkat-pangkat $T-2I$, serta ruang nol dan nulitasnya,
bersifat rutin.
(Ingat konvensi kita untuk
mengambil $T$ sebagai representasi transformasi $\map{t}{\C^3}{\C^3}$
terhadap basis standar.)
\begin{center}
  \renewcommand{\arraystretch}{1.25}
  \begin{tabular}{r|ccc} 
    \multicolumn{1}{c}{\( p \)}  
         &\( (T-2I)^p \) &\( \nullspace{(t-2)^p}  \) 
         &\textit{nulitas}                                            \\  
    \hline
    \( 1 \)
    &\matrixvenlarge{\begin{mat}[r]
          0  &0  &1  \\
          0  &4  &2  \\
          0  &0  &0
        \end{mat}}
    &\( \set{\matrixvenlarge{\colvec{x \\ 0 \\ 0}}
         \suchthat x\in\C}  \)  
    &$1$                                                   \\
    \( 2 \)
    &\matrixvenlarge{\begin{mat}[r]
          0  &0  &0  \\
          0  &16 &8  \\
          0  &0  &0
        \end{mat}}
    &\( \set{\matrixvenlarge{\colvec{x \\ -z/2 \\  z}}
               \suchthat x,z\in\C}  \) 
    &$2$                                                   \\
    \( 3 \)
    &\matrixvenlarge{\begin{mat}[r]
          0  &0  &0  \\
          0  &64 &32 \\
          0  &0  &0
        \end{mat}}
    &\textit{--sama--}
    &\textit{--sama--}
  \end{tabular}
\end{center}
Jadi, ruang nol tergeneralisasi $\gennullspace{t-2}$ berdimensi dua.
Kita mengetahui bahwa restriksi $t-2$ nilpoten pada subruang ini.
Dari pola pertumbuhan nulitas, kita mengetahui bahwa tindakan
$t-2$ pada suatu basis rantai adalah
$\vec{\beta}_1\mapsto\vec{\beta}_2\mapsto\zero$.  
Jadi, kita dapat merepresentasikan restriksi itu dalam bentuk kanonik
\begin{equation*}
  N_2=
  \begin{mat}[r]
    0  &0  \\
    1  &0  
  \end{mat}
  =\rep{t-2}{B,B}
  \qquad
   B_2=\sequence{\colvec[r]{1 \\ 1 \\ -2},
                 \colvec[r]{-2 \\ 0 \\ 0}}  
\end{equation*}
(pilihan basis lain juga mungkin).
Akibatnya, tindakan restriksi $t$ pada
$\gennullspace{t-2}$ direpresentasikan oleh blok Jordan berikut.
\begin{equation*}
  J_2=N_2+2I=\rep{t}{B_2,B_2}=
  \begin{mat}[r]
    2  &0  \\
    1  &2
  \end{mat}
\end{equation*}

Nilai eigen kedua adalah $6$.
Perhitungannya lebih mudah.
Karena pangkat $x-6$ dalam polinomial karakteristik adalah satu,
restriksi $t-6$ pada $\gennullspace{t-6}$
harus nilpoten berindeks satu
(indeksnya tidak mungkin kurang dari satu, dan karena $x-6$ merupakan
faktor polinomial karakteristik berpangkat satu, indeksnya juga tidak mungkin
lebih dari satu).
Tindakannya pada basis rantai harus berupa $\vec{\beta}_3\mapsto\zero$, dan
karena merupakan pemetaan nol, bentuk kanoniknya $N_6$
adalah matriks nol $\nbyn{1}$.
Akibatnya, bentuk kanonik $J_6$ untuk tindakan $t$ pada
$\gennullspace{t-6}$ adalah matriks blok Jordan $\nbyn{1}$
dengan satu-satunya entri $6$.
Untuk basisnya, kita dapat menggunakan sebarang vektor tak nol dari ruang nol tergeneralisasi.
\begin{equation*}
   B_6=\sequence{\colvec[r]{0 \\ 1 \\ 0}}
\end{equation*}

Jika digabungkan, keduanya memberikan
bentuk Jordan \( T \), yaitu
\begin{equation*}
   \rep{t}{B,B}=
   \begin{mat}[r]
      2  &0  &0  \\
      1  &2  &0  \\
      0  &0  &6
   \end{mat}
\end{equation*}
dengan \( B \) merupakan konkatenasi $B_2$ dan $B_6$.
(Jika kita ingin cermat memperoleh bentuk Jordan yang unik,
kita dapat, misalnya, menyusun ulang elemen basis agar nilai eigennya
berurutan menurun.
Namun, urutan meningkat juga boleh.)
\end{example}

\begin{example}  \label{SecJordanForm}
Matriks berikut
memiliki polinomial karakteristik yang sama
dengan contoh sebelumnya, \( -(x-2)^2(x-6) \).
\begin{equation*}
   T=
   \begin{mat}[r]
     2  &2  &1  \\
     0  &6  &2  \\
     0  &0  &2
   \end{mat}
\end{equation*}
Namun, di sini
tindakan $t-2$ stabil setelah satu kali penerapan saja\Dash restriksi
$t-2$ pada $\gennullspace{t-2}$ nilpoten berindeks satu.
\begin{center}
  % \renewcommand{\arraystretch}{1.25}
  \begin{tabular}{r|ccc}
    \multicolumn{1}{c}{\( p \)}  &\( (T-2I)^p \)  &\( \nullspace{(t-2)^p}  \)  
       &\textit{nulitas} \\                                     \hline
   \( 1 \)
   &\matrixvenlarge{\begin{mat}[r]
         0  &2  &1  \\
         0  &4  &2  \\
         0  &0  &0
       \end{mat}}
   &\( \set{\matrixvenlarge{\colvec{x \\ (-1/2)z \\ z}}\suchthat x,z\in\C}  \) 
   &$2$                                                     \\
   \( 2 \)
   &\matrixvenlarge{\begin{mat}[r]
         0  &8  &4 \\
         0  &16 &8 \\
         0  &0  &0
       \end{mat}}
   &\textit{--sama--}
   &\textit{--sama--}
 \end{tabular}
\end{center}
Jadi, restriksi $t-2$ pada ruang nol tergeneralisasi bertindak pada basis
rantai melalui dua rantai $\vec{\beta}_1\mapsto\zero$
dan $\vec{\beta}_2\mapsto\zero$.
Kita memperoleh blok Jordan berikut yang terkait dengan nilai eigen~$2$.
\begin{equation*}
  J_2=
  \begin{mat}[r]
    2  &0  \\
    0  &2  
  \end{mat}
\end{equation*}

Jadi, perbedaannya dengan contoh sebelumnya adalah bahwa meskipun
polinomial karakteristik memberi tahu kita untuk melihat
tindakan $t-2$ pada ruang nol tergeneralisasinya, polinomial karakteristik
tidak sepenuhnya menggambarkan tindakan $t-2$.
Kita harus melakukan beberapa perhitungan untuk memperoleh bahwa
polinomial minimalnya adalah \( (x-2)(x-6) \).

Untuk nilai eigen $6$, argumen bagi nilai eigen kedua pada
contoh sebelumnya kembali berlaku:~karena pangkat $x-6$ dalam
polinomial karakteristik adalah satu,
restriksi $t-6$ pada $\gennullspace{t-6}$
harus nilpoten berindeks satu.
Sebagai alternatif, kita dapat langsung menghitungnya.
\begin{center}
  \begin{tabular}{r|ccc}
    \multicolumn{1}{c}{\( p \)}  &\( (T-6I)^p \)  &\( \nullspace{(t-6)^p}  \)  
       &\textit{nulitas} \\                                     \hline
   \( 1 \)
   &\matrixvenlarge{\begin{mat}[r]
        -4  &2  &1  \\
         0  &0  &2  \\
         0  &0  &-4
       \end{mat}}
   &\( \set{\matrixvenlarge{\colvec{x \\ 2x \\ 0}}\suchthat x\in\C}  \) 
   &$1$                                                     \\
   \( 2 \)
   &\matrixvenlarge{\begin{mat}[r]
        16  &-8 &-4 \\
         0  &0  &-8 \\
         0  &0  &16
       \end{mat}}
   &\textit{--sama--}
   &\textit{--sama--}
 \end{tabular}
\end{center}
Bagaimanapun caranya, pada $\gennullspace{t-6}$, bentuk kanonik $N_6$
dari restriksi $t-6$ adalah matriks nol $\nbyn{1}$.
Blok Jordan $J_6$ adalah matriks $\nbyn{1}$ dengan entri $6$.
 
Dengan demikian, bentuk Jordan $T$ adalah matriks diagonal, sehingga~\( t \)
merupakan transformasi yang dapat didiagonalkan.
\begin{equation*}
  \rep{t}{B,B}=
  \begin{mat}[r]
    2  &0  &0  \\
    0  &2  &0  \\
    0  &0  &6
  \end{mat}
  \qquad
  B=\cat{B_2}{B_6}
   =\sequence{\colvec[r]{1 \\ 0 \\ 0},
              \colvec[r]{0 \\ 1 \\ -2},
              \colvec[r]{1 \\ 2 \\ 0}}
\end{equation*}
Dari ketiga vektor basis, dua yang pertama berasal dari ruang nol
$t-2$, dan yang ketiga berasal dari ruang nol~$t-6$.
\end{example}

\begin{example} \label{ThirdJordanForm}
Sedikit perhitungan dengan
\begin{equation*}
   T=
   \begin{mat}[r]
     -1  &4  &0  &0  &0  \\
      0  &3  &0  &0  &0  \\
      0  &-4 &-1 &0  &0  \\
      3  &-9 &-4 &2  &-1 \\
      1  &5  &4  &1  &4
   \end{mat}
\end{equation*}
menunjukkan bahwa polinomial karakteristiknya
adalah \( -(x-3)^3(x+1)^2 \).
Tabel berikut
\begin{center}
  % \renewcommand{\arraystretch}{1.25}
  \begin{tabular}{@{}r|c@{}c@{}c@{}}
    \multicolumn{1}{c}{\( p \)}  
            &\( (T-3I)^p \)  &\( \nullspace{(t-3)^p}  \)
             &\textit{nulitas}
    \\  \hline
    \( 1 \)
    &\matrixvenlarge{\begin{mat}[r]
         -4  &4  &0  &0  &0  \\
          0  &0  &0  &0  &0  \\
          0  &-4 &-4 &0  &0  \\
          3  &-9 &-4 &-1 &-1 \\
          1  &5  &4  &1  &1
       \end{mat}}
    &\( \set{\matrixvenlarge{\colvec{-(u+v)/2 \\
                     -(u+v)/2 \\
                      (u+v)/2 \\
                       u      \\
                       v}}\suchthat u,v\in\C}  \) 
    &$2$                                            \\
    \( 2 \)
    &\matrixvenlarge{\begin{mat}[r]
         16  &-16&0  &0  &0  \\
          0  &0  &0  &0  &0  \\
          0  &16 &16 &0  &0  \\
        -16  &32 &16 &0  &0  \\
          0  &-16&-16&0  &0
       \end{mat}}
    &\( \set{\matrixvenlarge{\colvec{ -z      \\
                      -z      \\
                       z      \\
                       u      \\
                       v}}\suchthat z,u,v\in\C}  \) 
    &$3$                                              \\
    \( 3 \)
    &\matrixvenlarge{\begin{mat}[r]
        -64  &64   &0   &0   &0  \\
          0  &0    &0   &0   &0  \\
          0  &-64  &-64 &0   &0  \\
         64  &-128 &-64 &0   &0  \\
          0  &64   &64  &0   &0
       \end{mat}}
    &\textit{--sama--}
    &\textit{--sama--}
  \end{tabular}
\end{center}
menunjukkan bahwa restriksi $t-3$ pada $\gennullspace{t-3}$ bertindak pada
basis rantai melalui dua rantai
$\vec{\beta}_1\mapsto\vec{\beta}_2\mapsto\zero$
dan
$\vec{\beta}_3\mapsto\zero$.

Perhitungan serupa untuk nilai eigen lainnya,
\begin{center}
  \renewcommand{\arraystretch}{1.25}
  \begin{tabular}{r|ccc}
    \multicolumn{1}{c}{\( p \)}  
         &\( (T+1I)^p \)  &\( \nullspace{(t+1)^p}  \) 
         &\textit{nulitas}  \\  
    \hline
    \( 1 \)
    &\matrixvenlarge{\begin{mat}[r]
          0  &4  &0  &0  &0  \\
          0  &4  &0  &0  &0  \\
          0  &-4 &0  &0  &0  \\
          3  &-9 &-4 &3  &-1 \\
          1  &5  &4  &1  &5
       \end{mat}}
    &\( \set{\matrixvenlarge{\colvec{-(u+v)   \\
                       0      \\
                      -v      \\
                       u      \\
                       v}}\suchthat u,v\in\C}  \)  
    &$2$                                              \\
    \( 2 \)
    &\matrixvenlarge{\begin{mat}[r]
          0  &16 &0  &0  &0  \\
          0  &16 &0  &0  &0  \\
          0  &-16&0  &0  &0  \\
          8  &-40&-16&8  &-8 \\
          8  &24 &16 &8  &24
       \end{mat}}
    &\textit{--sama--}
    &\textit{--sama--}
  \end{tabular}
\end{center}
memberikan bahwa restriksi $t+1$ pada ruang nol tergeneralisasinya
bertindak pada basis rantai melalui dua rantai terpisah,
$\vec{\beta}_4\mapsto\zero$ dan $\vec{\beta}_5\mapsto\zero$.

Dengan demikian,
\( T \) serupa dengan matriks berbentuk Jordan berikut.
\begin{equation*}
   \begin{mat}[r]
     -1  &0  &0  &0  &0  \\
      0  &-1 &0  &0  &0  \\    
      0  &0  &3  &0  &0  \\
      0  &0  &1  &3  &0  \\
      0  &0  &0  &0  &3  
   \end{mat}
\end{equation*}
\end{example}





\begin{exercises}
  \item 
    Lakukan pemeriksaan untuk \nearbyexample{ex:SingJordBlock}.
    \begin{answer}
      Kita harus memeriksa bahwa
      \begin{equation*}
         \begin{mat}[r]
           3  &0  \\
           1  &3
         \end{mat}
         =
         N+3I=PTP^{-1}
         =
         \begin{mat}[r]
           1/2  &1/2  \\
          -1/4 &1/4
         \end{mat}
         \begin{mat}[r]
           2  &-1  \\
           1  &4
         \end{mat}
         \begin{mat}[r]
           1  &-2  \\
           1  &2
         \end{mat}
      \end{equation*}
      Perhitungan itu mudah.
    \end{answer}
  \item 
    Setiap matriks berikut berbentuk Jordan.
    Nyatakan polinomial karakteristik dan polinomial minimalnya.
    \begin{exparts*}
      \partsitem 
        $\begin{mat}[r]
           3  &0  \\
           1  &3        
         \end{mat}$
      \partsitem 
        $\begin{mat}[r]
           -1  &0  \\
            0  &-1 
         \end{mat}$
      \partsitem 
        $\begin{mat}[r]
           2  &0  &0  \\
           1  &2  &0  \\
           0  &0  &-1/2
         \end{mat}$
      \partsitem 
        $\begin{mat}[r]
           3  &0  &0  \\
           1  &3  &0  \\
           0  &1  &3  \\ 
         \end{mat}$
      \partsitem 
        $\begin{mat}[r]
           3  &0  &0  &0  \\
           1  &3  &0  &0  \\
           0  &0  &3  &0  \\
           0  &0  &1  &3 
         \end{mat}$
      \partsitem 
        $\begin{mat}[r]
           4  &0  &0  &0  \\
           1  &4  &0  &0  \\
           0  &0  &-4 &0  \\
           0  &0  &1  &-4
         \end{mat}$
      \partsitem 
        $\begin{mat}[r]
           5  &0  &0  \\
           0  &2  &0  \\
           0  &0  &3  
         \end{mat}$
      \partsitem 
        $\begin{mat}[r]
           5  &0  &0  &0  \\
           0  &2  &0  &0  \\
           0  &0  &2  &0  \\
           0  &0  &0  &3 
         \end{mat}$
      \partsitem 
        $\begin{mat}[r]
           5  &0  &0  &0  \\
           0  &2  &0  &0  \\
           0  &1  &2  &0  \\
           0  &0  &0  &3 
         \end{mat}$
    \end{exparts*}
    \begin{answer}
      \begin{exparts}
        \partsitem Polinomial karakteristiknya adalah $c(x)=(x-3)^2$, dan
          polinomial minimalnya sama.
        \partsitem Polinomial karakteristiknya adalah $c(x)=(x+1)^2$.
          Polinomial minimalnya adalah $m(x)=x+1$.
        \partsitem Polinomial karakteristiknya adalah
          $c(x)=-(x+(1/2))(x-2)^2$, dan
          polinomial minimalnya sama.
        \partsitem Polinomial karakteristiknya adalah $c(x)=-(x-3)^3$.
          Polinomial minimalnya sama.
        \partsitem Polinomial karakteristiknya adalah $c(x)=(x-3)^4$.
          Polinomial minimalnya adalah $m(x)=(x-3)^2$.
        \partsitem Polinomial karakteristiknya adalah $c(x)=(x+4)^2(x-4)^2$, dan
          polinomial minimalnya sama.
        \partsitem Polinomial karakteristiknya adalah
          $c(x)=-(x-5)(x-2)(x-3)$, dan polinomial minimalnya sama.
        \partsitem Polinomial karakteristiknya adalah
          $c(x)=(x-2)^2(x-3)(x-5)$, dan
          polinomial minimalnya adalah $m(x)=(x-2)(x-3)(x-5)$.
        \partsitem Polinomial karakteristiknya adalah
          $c(x)=(x-2)^2(x-3)(x-5)$, dan
          polinomial minimalnya sama.
      \end{exparts}
    \end{answer}
  \recommended \item
    Carilah bentuk Jordan dari data yang diberikan.
    \begin{exparts}
       \partsitem Matriks
         \( T \) berukuran \( \nbyn{5} \) dengan satu nilai eigen $3$.
         Nulitas pangkat-pangkatnya adalah:
         \( T-3I \) bernulitas dua, \( (T-3I)^2 \) bernulitas tiga,
         \( (T-3I)^3 \) bernulitas empat, dan \( (T-3I)^4 \) bernulitas
         lima.
       \partsitem Matriks \( S \) berukuran \( \nbyn{5} \) dengan dua nilai eigen.
         Untuk nilai eigen $2$, nulitasnya adalah:
         \( S-2I \) bernulitas dua, dan \( (S-2I)^2 \) bernulitas empat.
         Untuk nilai eigen $-1$, nulitasnya adalah:
         \( S+1I \) bernulitas satu.
    \end{exparts}
    \begin{answer}
      \begin{exparts}
         \partsitem Transformasi $t-3$ nilpoten
          (yakni, $\gennullspace{t-3}$ adalah seluruh ruang)
          dan bertindak pada basis rantai melalui dua rantai,
          $\vec{\beta}_1\mapsto\vec{\beta}_2\mapsto\vec{\beta}_3
            \mapsto\vec{\beta}_4\mapsto\zero$
          dan $\vec{\beta}_5\mapsto\zero$.
          Akibatnya, $t-3$ dapat direpresentasikan dalam bentuk kanonik berikut.
          \begin{equation*}
            N_3=
            \begin{mat}[r]
               0  &0  &0  &0  &0  \\
               1  &0  &0  &0  &0  \\
               0  &1  &0  &0  &0  \\
               0  &0  &1  &0  &0  \\
               0  &0  &0  &0  &0
             \end{mat}
          \end{equation*}
          dan karena itu $T$ serupa dengan matriks bentuk kanonik berikut.
          \begin{equation*}
           J_3=N_3+3I=
           \begin{mat}[r]
               3  &0  &0  &0  &0  \\
               1  &3  &0  &0  &0  \\
               0  &1  &3  &0  &0  \\
               0  &0  &1  &3  &0  \\
               0  &0  &0  &0  &3
             \end{mat}
         \end{equation*} 
         \partsitem Restriksi transformasi $s+1$ nilpoten
          pada subruang $\gennullspace{s+1}$, dan tindakannya pada
          basis rantai adalah $\vec{\beta}_1\mapsto\zero$.
          Restriksi transformasi $s-2$ nilpoten
          pada subruang $\gennullspace{s-2}$, dengan tindakan pada
          basis rantai berupa $\vec{\beta}_2\mapsto\vec{\beta}_3\mapsto\zero$
          dan $\vec{\beta}_4\mapsto\vec{\beta}_5\mapsto\zero$.
          Akibatnya, bentuk Jordannya adalah berikut.
          \begin{equation*}
            \begin{mat}[r]
              -1  &0  &0  &0  &0  \\
               0  &2  &0  &0  &0  \\
               0  &1  &2  &0  &0  \\
               0  &0  &0  &2  &0  \\
               0  &0  &0  &1  &2
             \end{mat} 
          \end{equation*}
      \end{exparts}  
    \end{answer}
  \item 
    Carilah matriks perubahan basis untuk setiap contoh.
    \begin{exparts*}
      \partsitem \nearbyexample{ex:FirstJordForm}
      \partsitem \nearbyexample{SecJordanForm}
      \partsitem \nearbyexample{ThirdJordanForm}
    \end{exparts*}
    \begin{answer}
      Untuk masing-masing contoh, karena banyak pilihan basis dimungkinkan,
      banyak jawaban lain juga dimungkinkan.
      Tentu saja, perhitungan untuk memeriksa apakah suatu jawaban membuat
      $PTP^{-1}$ berbentuk Jordan merupakan penentu kebenarannya.
      \begin{exparts}
        \partsitem Berikut diagram panahnya.
          \begin{equation*}
            \begin{CD}
              \C^3_{\wrt{\stdbasis_3}}      @>t>T>    \C^3_{\wrt{\stdbasis_3}}   \\
                @V\scriptstyle\identity V\scriptstyle PV  
                                    @V\scriptstyle\identity V\scriptstyle PV \\
              \C^3_{\wrt{B}}                 @>t>J>         \C^3_{\wrt{B}}
            \end{CD}
          \end{equation*}
          Matriks untuk berpindah dari kiri bawah ke kiri atas adalah berikut.
          \begin{equation*}
            P^{-1}=\bigl(\rep{\identity}{\stdbasis_3,B}\bigr)^{-1}
                  =\rep{\identity}{B,\stdbasis_3}   
                  =\begin{mat}[r]
                     1  &-2   &0 \\
                     1  &0   &1 \\
                    -2  &0   &0
                   \end{mat}
          \end{equation*}
          Matriks $P$ untuk berpindah dari kanan atas ke kanan
          bawah adalah invers $P^{-1}$.
        \partsitem Kita menginginkan matriks berikut beserta inversnya.
          \begin{equation*}
            P^{-1}=
            \begin{mat}[r]
              1  &0  &3  \\
              0  &1  &4  \\
              0  &-2 &0
            \end{mat}
          \end{equation*}
        \partsitem Konkatenasi basis-basis untuk
          ruang nol tergeneralisasi berikut dapat digunakan sebagai basis
          seluruh ruang.
          \begin{equation*}
             B_{-1}=\sequence{\colvec[r]{-1\\ 0 \\  0 \\ 1 \\ 0},
                 \colvec[r]{-1\\ 0 \\ -1 \\ 0 \\ 1}}
            \qquad
             B_3=\sequence{\colvec[r]{1 \\ 1 \\ -1 \\ 0 \\ 0},
                 \colvec[r]{0 \\ 0 \\  0 \\-2 \\ 2},
                 \colvec[r]{-1\\-1 \\  1 \\ 2 \\ 0}}
          \end{equation*}
          Matriks perubahan basisnya adalah matriks berikut beserta inversnya.
          \begin{equation*}
            P^{-1}=
            \begin{mat}[r]
              -1  &-1  &1  &0  &-1  \\
              0   &0   &1  &0  &-1  \\
              0   &-1  &-1 &0  &1   \\
              1   &0   &0  &-2 &2   \\
              0   &1   &0  &2  &0   \\
            \end{mat}
          \end{equation*}
      \end{exparts}
    \end{answer}
  \recommended \item 
    Carilah bentuk Jordan dan basis Jordan untuk setiap matriks.
    \begin{exparts*}
      \partsitem 
        \(
        \begin{mat}[r]
          -10  &4  \\
          -25  &10
        \end{mat} \)
      \partsitem 
        \(
        \begin{mat}[r]
           5   &-4 \\
           9   &-7
        \end{mat} \)
      \partsitem 
        \(
        \begin{mat}[r]
           4   &0    &0  \\
           2   &1    &3  \\
           5   &0    &4
        \end{mat} \)
      \partsitem 
        \(
        \begin{mat}[r]
           5   &4    &3  \\
          -1   &0    &-3 \\
           1   &-2   &1
        \end{mat} \)
      \partsitem
        \(
        \begin{mat}[r]
           9   &7    &3  \\
          -9   &-7   &-4 \\
           4   &4    &4
        \end{mat} \)
      \partsitem 
        \(
        \begin{mat}[r]
           2   &2    &-1 \\
          -1   &-1   &1  \\
          -1   &-2   &2
        \end{mat} \)
      \partsitem 
        \(
        \begin{mat}[r]
           7   &1    &2   &2 \\
           1   &4    &-1  &-1\\
          -2   &1    &5   &-1\\
           1   &1    &2   &8
        \end{mat} \)
    \end{exparts*}
    \begin{answer}
      Prosedur umumnya adalah memfaktorkan polinomial karakteristik
      $c(x)=(x-\lambda_1)^{p_1}(x-\lambda_2)^{p_2}\cdots$
      untuk memperoleh nilai eigen $\lambda_1$, $\lambda_2$, dan seterusnya.
      Kemudian, untuk setiap $\lambda_i$, kita mencari
      basis rantai bagi tindakan transformasi $t-\lambda_i$
      ketika direstriksi pada $\gennullspace{t-\lambda_i}$,
      dengan menghitung pangkat-pangkat matriks $T-\lambda_iI$ dan mencari
      ruang nol terkait hingga ruang-ruang nol tersebut stabil
      (tidak berubah); pada saat itu, kita telah memperoleh ruang nol tergeneralisasi.
      Dimensi ruang-ruang nol tersebut (nulitasnya) memberi tahu kita
      tindakan $t-\lambda_i$ pada basis rantai untuk ruang nol
      tergeneralisasi, sehingga kita dapat menuliskan pola satu pada subdiagonal
      untuk memperoleh $N_{\lambda_i}$.
      Dari matriks ini, blok Jordan $J_{\lambda_i}$ yang terkait
      dengan $\lambda_i$ langsung diperoleh: $J_{\lambda_i}=N_{\lambda_i}+\lambda_iI$.
      Terakhir, setelah melakukannya untuk setiap nilai eigen, kita menggabungkannya
      menjadi bentuk kanonik.
      \begin{exparts}
        \partsitem Polinomial karakteristik matriks ini
           adalah $c(x)=(-10-x)(10-x)+100=x^2$,
           sehingga matriks itu hanya memiliki nilai eigen $\lambda=0$.
           \begin{center}
             \renewcommand{\arraystretch}{1.25}
             \begin{tabular}{r|ccc}
                \multicolumn{1}{c}{\textit{pangkat}~$p$}
                    &$(T+0\cdot I)^p$ &$\nullspace{(t-0)^p}$
                    &\textit{nulitas}                   \\ 
                \hline
                $1$  
                &\matrixvenlarge{\begin{mat}[r]
                  -10  &4  \\ 
                  -25  &10
                \end{mat}}
                &$\set{\matrixvenlarge{\colvec{2y/5 \\ y}}
                                     \suchthat
                                     y\in\C}$
                &$1$                                       \\
                $2$  
                &\matrixvenlarge{\begin{mat}[r]
                    0  &0  \\ 
                    0  &0
                \end{mat}}
                &$\C^2$
                &$2$
             \end{tabular}
           \end{center}
           (Jadi, transformasi ini nilpoten:
           $\gennullspace{t-0}$ adalah seluruh ruang.)
           Dari nulitasnya, kita mengetahui bahwa
           tindakan $t$ pada basis rantai adalah
           $\vec{\beta}_1\mapsto\vec{\beta}_2\mapsto\zero$.
           Berikut matriks bentuk kanonik untuk tindakan $t-0$ pada
           $\gennullspace{t-0}=\C^2$,
           \begin{equation*}
             N_0=
             \begin{mat}[r]
               0  &0  \\
               1  &0
            \end{mat}
           \end{equation*}
           dan berikut bentuk Jordan matriks tersebut.
           \begin{equation*}
             J_0=N_0+0\cdot I=
             \begin{mat}[r]
               0  &0  \\
               1  &0
            \end{mat}
           \end{equation*}
           Perhatikan bahwa jika suatu matriks nilpoten, maka bentuk kanoniknya
           sama dengan bentuk Jordannya.

           Kita dapat mencari basis rantai semacam itu menggunakan teknik-teknik pada
           bagian sebelumnya.
           \begin{equation*}
                 B=\sequence{\colvec[r]{1 \\ 0},
                             \colvec[r]{-10 \\ -25}}
           \end{equation*}
           Kita memilih vektor basis pertama agar berada dalam
           ruang nol $t^2$, tetapi tidak dalam ruang nol $t$.
           Vektor basis kedua merupakan citra vektor pertama oleh $t$.
        \partsitem Polinomial karakteristik matriks ini
           adalah \( c(x)=(x+1)^2 \), sehingga matriks ini hanya memiliki satu nilai eigen.
           (Artinya, ruang nol tergeneralisasi $t+1$ adalah seluruh
           ruang.)
           Kita mempunyai
           \begin{equation*}
             \nullspace{t+1}=\set{\colvec{2y/3 \\ y}\suchthat
                                     y\in\C} 
             \qquad
             \nullspace{(t+1)^2}=\C^2 
           \end{equation*}
           sehingga tindakan $t+1$ pada
           basis rantai terkait adalah
           $\vec{\beta}_1\mapsto\vec{\beta}_2\mapsto\zero$.
           Jadi,
           \begin{equation*}
             N_{-1}
             =
             \begin{mat}[r]
                0  &0  \\
                1  &0 
             \end{mat}
           \end{equation*}
           bentuk Jordan $T$ adalah
           \begin{equation*}
             J_{-1}=N_{-1}+(-1)\cdot I
             =
             \begin{mat}[r]
               -1  &0  \\
                1  &-1
             \end{mat}
           \end{equation*}
           dan memilih vektor-vektor dari ruang nol di atas menghasilkan
           basis rantai berikut (pilihan lain juga mungkin).
           \begin{equation*}
             B=\sequence{\colvec[r]{1 \\ 0},
                         \colvec[r]{6 \\ 9}}
           \end{equation*}
        \partsitem Polinomial karakteristik
            \( c(x)=(1-x)(4-x)^2=-1\cdot (x-1)(x-4)^2 \) memiliki dua akar,
            yang merupakan nilai eigen $\lambda_1=1$ dan $\lambda_2=4$.

            Kita menangani kedua nilai eigen secara terpisah.
            Untuk $\lambda_1$, perhitungan pangkat-pangkat $T-1I$
            menghasilkan
            \begin{equation*}
              \nullspace{t-1}=\set{\colvec{0 \\ y \\ 0}
                                      \suchthat y\in\C}
            \end{equation*}
            dan ruang nol $(t-1)^2$ sama.
            Jadi, himpunan ini merupakan ruang nol tergeneralisasi
            $\gennullspace{t-1}$.
            Nulitasnya menunjukkan bahwa tindakan restriksi $t-1$
            pada ruang nol tergeneralisasi terhadap suatu basis rantai
            adalah $\vec{\beta}_1\mapsto\zero$.

            Perhitungan serupa untuk $\lambda_2=4$ memberikan ruang nol berikut.
            \begin{equation*}
              \nullspace{t-4}=\set{\colvec{0 \\ z \\ z}
                                      \suchthat z\in\C}
              \qquad
              \nullspace{(t-4)^2}=\set{\colvec{z-y \\ y \\ z}
                                          \suchthat y,z\in\C}
            \end{equation*}
            (Ruang nol $(t-4)^3$ sama, sebagaimana seharusnya karena
            pangkat faktor yang terkait dengan $\lambda_2=4$ dalam
            polinomial karakteristik adalah dua, sehingga restriksi
            $t-4$ pada ruang nol tergeneralisasi $\gennullspace{t-4}$
            nilpoten berindeks paling tinggi dua\Dash diperlukan paling banyak
            dua penerapan $t-4$ agar ruang nol menjadi stabil.)
            Pola pertumbuhan nulitas memberi tahu kita bahwa
            tindakan $t-4$ pada basis rantai terkait
            untuk $\gennullspace{t-4}$ adalah
            $\vec{\beta}_2\mapsto\vec{\beta}_3\mapsto\zero$.

            Menggabungkan informasi untuk kedua nilai eigen
            memberikan bentuk Jordan transformasi $t$.
            \begin{equation*}
              \begin{mat}[r]
                1  &0  &0  \\
                0  &4  &0  \\
                0  &1  &4
              \end{mat}
            \end{equation*}
            Kita dapat mengambil elemen-elemen ruang nol untuk memperoleh
            basis yang sesuai.
            \begin{equation*}
              B=\cat{B_{1}}{B_4}=
               \sequence{\colvec[r]{0 \\ 1 \\ 0},
                          \colvec[r]{1 \\ 0 \\ 1},
                          \colvec[r]{0 \\ 5 \\ 5}}
            \end{equation*}
        \partsitem Polinomial karakteristiknya adalah
            \( c(x)=(-2-x)(4-x)^2=-1\cdot (x+2)(x-4)^2 \).

            Untuk nilai eigen $\lambda_{-2}$
            (subskripnya adalah \( -2\)), perhitungan
            pangkat-pangkat $T+2I$ menghasilkan
            \begin{equation*}
              \nullspace{t+2}=\set{\colvec{z \\ z \\ z}
                                      \suchthat z\in\C}
            \end{equation*}
            Ruang nol $(t+2)^2$ sama, sehingga
            inilah ruang nol tergeneralisasi $\gennullspace{t+2}$.
            Jadi, tindakan restriksi $t+2$ pada
            $\gennullspace{t+2}$ terhadap basis
            rantai terkait adalah $\vec{\beta}_1\mapsto\zero$.

            Untuk $\lambda_2=4$,
            menghitung pangkat-pangkat $T-4I$ menghasilkan
            \begin{equation*}
              \nullspace{t-4}=\set{\colvec{z \\ -z \\ z}
                                      \suchthat z\in\C} 
              \qquad
              \nullspace{(t-4)^2}=\set{\colvec{x \\ -z \\ z}
                                           \suchthat x,z\in\C}
            \end{equation*}
            sehingga tindakan $t-4$ pada basis rantai untuk
            $\gennullspace{t-4}$ adalah
            $\vec{\beta}_2\mapsto\vec{\beta}_3\mapsto\zero$.

            Dengan demikian, bentuk Jordannya adalah
            \begin{equation*}
              \begin{mat}[r]
                -2  &0  &0  \\
                 0  &4  &0  \\
                 0  &1  &4
              \end{mat}
            \end{equation*}
            dan berikut basis yang sesuai.
            \begin{equation*}
              B=\cat{B_{-2}}{B_4}=
                \sequence{\colvec[r]{1 \\ 1 \\ 1},
                          \colvec[r]{0 \\ -1 \\ 1},
                          \colvec[r]{-1 \\ 1 \\ -1}}
            \end{equation*}
        \partsitem Polinomial karakteristik
            matriks ini adalah \( c(x)=(2-x)^3=-1\cdot (x-2)^3 \).
            Matriks ini hanya memiliki satu nilai eigen, $\lambda=2$.
            Dengan mencari pangkat-pangkat $T-2I$, kita memperoleh
            \begin{equation*}
              \nullspace{t-2}=\set{\colvec{-y \\ y \\ 0}
                                      \suchthat y\in\C} 
              \qquad
              \nullspace{(t-2)^2}=\set{\colvec{-y-(1/2)z \\ y \\ z}
                                          \suchthat y,z\in\C}
            \end{equation*}
            dan
            \begin{equation*} 
              \nullspace{(t-2)^3}=\C^3
            \end{equation*}
            sehingga
            tindakan $t-2$ pada
            basis rantai terkait adalah
            $\vec{\beta}_1\mapsto\vec{\beta}_2\mapsto
                   \vec{\beta}_3\mapsto\zero$.
            Bentuk Jordannya adalah berikut,
            \begin{equation*}
                  \begin{mat}[r]
                    2  &0  &0  \\
                    1  &2  &0  \\
                    0  &1  &2
                  \end{mat}
            \end{equation*}
            dan berikut salah satu pilihan basisnya.
            \begin{equation*}
              B=\sequence{\colvec[r]{0 \\ 1 \\ 0},
                          \colvec[r]{7 \\ -9 \\ 4},
                          \colvec[r]{-2 \\ 2 \\ 0}}
            \end{equation*}
        \partsitem Polinomial karakteristik
            \( c(x)=(1-x)^3=-(x-1)^3 \) hanya memiliki satu akar,
            sehingga matriks tersebut hanya memiliki satu nilai eigen $\lambda=1$.
            Mencari pangkat-pangkat $T-1I$
            dan menghitung ruang nol
            \begin{equation*}
               \nullspace{t-1}=\set{\colvec{-2y+z \\ y \\ z}
                                      \suchthat y,z\in\C} 
               \qquad
               \nullspace{(t-1)^2}=\C^3 
            \end{equation*}
            menunjukkan bahwa tindakan pemetaan nilpoten $t-1$ pada basis
            rantai adalah
            $\vec{\beta}_1\mapsto\vec{\beta}_2\mapsto\zero$ dan
            $\vec{\beta}_3\mapsto\zero$.
            Dengan demikian, bentuk Jordannya adalah
            \begin{equation*}
                  J=
                  \begin{mat}[r]
                    1  &0  &0  \\
                    1  &1  &0  \\
                    0  &0  &1
                  \end{mat}
            \end{equation*}
            dan basis yang sesuai (basis rantai yang terkait dengan
            $t-1$) adalah berikut.
            \begin{equation*}
              B=\sequence{\colvec[r]{0 \\ 1 \\ 0},
                          \colvec[r]{2 \\ -2 \\ -2},
                          \colvec[r]{1 \\ 0 \\ 1}}
            \end{equation*}
        \partsitem Polinomial karakteristiknya agak besar untuk dihitung
            dengan tangan, tetapi masih dapat ditangani,
            \( c(x)=x^4-24x^3+216x^2-864x+1296=(x-6)^4 \).
            Ini merupakan pemetaan dengan satu nilai eigen, sehingga
            transformasi $t-6$ nilpoten.
            Ruang-ruang nol
            \begin{equation*}
               \nullspace{t-6}=\set{\colvec{-z-w \\ -z-w \\ z \\ w}
                                      \suchthat z,w\in\C} 
               \qquad
               \nullspace{(t-6)^2}=\set{\colvec{x \\ -z-w \\ z \\ w}
                                      \suchthat x,z,w\in\C} 
            \end{equation*}
            dan
            \begin{equation*}
               \nullspace{(t-6)^3}=\C^4 
            \end{equation*}
            beserta nulitasnya
            menunjukkan bahwa tindakan $t-6$ pada basis rantai adalah
            $\vec{\beta}_1\mapsto\vec{\beta}_2\mapsto
                   \vec{\beta}_3\mapsto\zero$ dan
            $\vec{\beta}_4\mapsto\zero$.
            Bentuk Jordannya adalah
            \begin{equation*}
              \begin{mat}[r]
                6  &0  &0  &0  \\
                1  &6  &0  &0  \\
                0  &1  &6  &0  \\
                0  &0  &0  &6  \\
              \end{mat}
            \end{equation*}
            dan pencarian basis rantai yang sesuai bersifat rutin.
            \begin{equation*}
              B=\sequence{\colvec[r]{0 \\ 0 \\ 0 \\ 1},
                          \colvec[r]{2 \\ -1 \\ -1 \\ 2},
                          \colvec[r]{3 \\ 3 \\ -6 \\ 3},
                          \colvec[r]{-1 \\ -1 \\ 1 \\ 0}}
            \end{equation*}
      \end{exparts}  
    \end{answer}
  \recommended \item
    Carilah semua bentuk Jordan yang mungkin dari transformasi dengan polinomial
    karakteristik \( (x-1)^2(x+2)^2 \).
    \begin{answer}
      Terdapat dua nilai eigen, $\lambda_1=-2$ dan $\lambda_2=1$.
      Restriksi $t+2$ pada
      $\gennullspace{t+2}$ dapat memiliki salah satu dari tindakan berikut
      pada basis rantai terkait.
      \begin{equation*}
        \vec{\beta}_1\mapsto\vec{\beta}_2\mapsto\zero
        \qquad
        \begin{array}[t]{l}
          \vec{\beta}_1\mapsto\zero  \\
          \vec{\beta}_2\mapsto\zero 
        \end{array}
      \end{equation*}
      Restriksi $t-1$ pada
      $\gennullspace{t-1}$ dapat memiliki salah satu dari tindakan berikut
      pada basis rantai terkait.
      \begin{equation*}
        \vec{\beta}_3\mapsto\vec{\beta}_4\mapsto\zero
        \qquad
        \begin{array}[t]{l}
          \vec{\beta}_3\mapsto\zero  \\
          \vec{\beta}_4\mapsto\zero 
        \end{array}
      \end{equation*}
      Jika digabungkan, ini menghasilkan empat bentuk Jordan yang mungkin:
      kedua tindakan pertama, tindakan kedua dan pertama, tindakan pertama dan kedua,
      serta kedua tindakan kedua.
      \begin{equation*}
        \begin{mat}[r]
          -2  &0  &0  &0  \\
           1  &-2 &0  &0  \\
           0  &0  &1  &0  \\
           0  &0  &1  &1
        \end{mat}
        \quad
        \begin{mat}[r]
          -2  &0  &0  &0  \\
           0  &-2 &0  &0  \\
           0  &0  &1  &0  \\
           0  &0  &1  &1
        \end{mat}
        \quad
        \begin{mat}[r]
          -2  &0  &0  &0  \\
           1  &-2 &0  &0  \\
           0  &0  &1  &0  \\
           0  &0  &0  &1
        \end{mat}
        \quad
        \begin{mat}[r]
          -2  &0  &0  &0  \\
           0  &-2 &0  &0  \\
           0  &0  &1  &0  \\
           0  &0  &0  &1
        \end{mat}
     \end{equation*}  
    \end{answer}
  \item 
    Carilah semua bentuk Jordan yang mungkin dari transformasi dengan polinomial
    karakteristik \( (x-1)^3(x+2) \).
    \begin{answer}
     Restriksi $t+2$ pada
     $\gennullspace{t+2}$ hanya dapat memiliki tindakan
     $\vec{\beta}_1\mapsto\zero$.
     Restriksi $t-1$ pada $\gennullspace{t-1}$ dapat memiliki salah satu
     dari ketiga tindakan berikut pada basis rantai terkait.
     \begin{equation*}
        \vec{\beta}_2\mapsto\vec{\beta}_3\mapsto\vec{\beta}_4\mapsto\zero
        \qquad
        \begin{array}[t]{l}
          \vec{\beta}_2\mapsto\vec{\beta}_3\mapsto\zero  \\
          \vec{\beta}_4\mapsto\zero 
        \end{array}
        \qquad
        \begin{array}[t]{l}
          \vec{\beta}_2\mapsto\zero  \\
          \vec{\beta}_3\mapsto\zero  \\
          \vec{\beta}_4\mapsto\zero 
        \end{array}
     \end{equation*}
     Jika digabungkan, terdapat tiga bentuk Jordan yang mungkin:
     bentuk yang berasal dari tindakan pertama $t-1$ (bersama satu-satunya
     tindakan $t+2$), bentuk yang berasal dari tindakan kedua, dan
     bentuk yang berasal dari tindakan ketiga.
     \begin{equation*}
       \begin{mat}[r]
         -2  &0  &0  &0  \\
          0  &1  &0  &0  \\
          0  &1  &1  &0  \\
          0  &0  &1  &1
       \end{mat}
       \quad
       \begin{mat}[r]
         -2  &0  &0  &0  \\
          0  &1  &0  &0  \\
          0  &1  &1  &0  \\
          0  &0  &0  &1
       \end{mat}
       \quad
       \begin{mat}[r]
         -2  &0  &0  &0  \\
          0  &1  &0  &0  \\
          0  &0  &1  &0  \\
          0  &0  &0  &1
       \end{mat}
     \end{equation*}
    \end{answer}
  \recommended \item
    Carilah semua bentuk Jordan yang mungkin dari transformasi dengan polinomial
    karakteristik \( (x-2)^3(x+1) \) dan polinomial minimal \( (x-2)^2(x+1) \).
    \begin{answer}
      Tindakan $t+1$ pada basis rantai untuk $\gennullspace{t+1}$
      harus berupa $\vec{\beta}_1\mapsto\zero$.
      Karena pangkat \( x-2 \) dalam polinomial minimal, suatu
      rantai dalam basis rantai untuk $t-2$ memiliki panjang dua, sehingga
      tindakan \( t-2 \) pada \( \gennullspace{t-2} \)
      harus berbentuk berikut.
      \begin{equation*}
        \begin{array}[t]{l}
          \vec{\beta}_2\mapsto\vec{\beta}_3\mapsto\zero  \\
          \vec{\beta}_4\mapsto\zero 
        \end{array}        
      \end{equation*}
      Jadi, hanya ada satu bentuk Jordan yang mungkin.
      \begin{equation*}
          \begin{mat}[r]
            -1  &0  &0  &0  \\
             0  &2  &0  &0  \\
             0  &1  &2  &0  \\
             0  &0  &0  &2
          \end{mat}
       \end{equation*}
     \end{answer}
  \item 
    Carilah semua bentuk Jordan yang mungkin dari transformasi dengan polinomial
    karakteristik \( (x-2)^4(x+1) \) dan polinomial minimal \( (x-2)^2(x+1) \).
    \begin{answer}
      Terdapat dua bentuk Jordan yang mungkin.
      Tindakan $t+1$ pada basis rantai untuk $\gennullspace{t+1}$
      harus berupa $\vec{\beta}_1\mapsto\zero$.
      Terdapat dua tindakan $t-2$ pada basis rantai untuk
      $\gennullspace{t-2}$ yang dimungkinkan oleh polinomial karakteristik
      dan polinomial minimal ini.
      \begin{equation*}
        \begin{array}[t]{l}
          \vec{\beta}_2\mapsto\vec{\beta}_3\mapsto\zero  \\
          \vec{\beta}_4\mapsto\vec{\beta}_5\mapsto\zero  
        \end{array}        
        \qquad
        \begin{array}[t]{l}
          \vec{\beta}_2\mapsto\vec{\beta}_3\mapsto\zero  \\
          \vec{\beta}_4\mapsto\zero                      \\
          \vec{\beta}_5\mapsto\zero                      
        \end{array}        
      \end{equation*}
      Matriks bentuk Jordan yang dihasilkan adalah berikut.
      \begin{equation*}
        \begin{mat}[r]
          -1  &0  &0  &0  &0  \\
           0  &2  &0  &0  &0  \\
           0  &1  &2  &0  &0  \\
           0  &0  &0  &2  &0  \\
           0  &0  &0  &1  &2
        \end{mat}
        \qquad
        \begin{mat}[r]
          -1  &0  &0  &0  &0  \\
           0  &2  &0  &0  &0  \\
           0  &1  &2  &0  &0  \\
           0  &0  &0  &2  &0  \\
           0  &0  &0  &0  &2
        \end{mat}
     \end{equation*}  
    \end{answer}
  \recommended \item Diagonalkan matriks-matriks berikut.
    \begin{exparts*}
       \partsitem \( \begin{mat}[r]
                  1  &1  \\
                  0  &0
                \end{mat}  \)
       \partsitem \( \begin{mat}[r]
                  0  &1  \\
                  1  &0
                \end{mat}  \)
    \end{exparts*}
    \begin{answer}
      \begin{exparts}
        \partsitem Polinomial karakteristiknya adalah \( c(x)=x(x-1) \).
          Untuk $\lambda_1=0$, kita mempunyai
          \begin{equation*}
            \nullspace{t-0}=\set{\colvec{-y \\ y}
                                 \suchthat y\in\C }  
          \end{equation*} 
          (tentu saja, ruang nol $t^2$ sama).
          Untuk $\lambda_2=1$,
          \begin{equation*}
            \nullspace{t-1}=\set{\colvec{x \\ 0}
                                     \suchthat x\in\C }  
          \end{equation*}
          (dan ruang nol $(t-1)^2$ sama).
          Kita dapat mengambil basis berikut
          \begin{equation*}
            B=\sequence{\colvec[r]{1 \\ -1},\colvec[r]{1 \\ 0}}
          \end{equation*}
          untuk memperoleh diagonalisasi.
          \begin{equation*}
            \begin{mat}[r]
              1  &1  \\
             -1  &0
            \end{mat}^{-1}
            \begin{mat}[r]
              1  &1  \\
              0  &0
            \end{mat}
            \begin{mat}[r]
              1  &1  \\
             -1  &0
            \end{mat}
            =
            \begin{mat}[r]
              0  &0  \\
              0  &1
            \end{mat}
          \end{equation*}
        \partsitem Polinomial karakteristiknya adalah
          \( c(x)=x^2-1=(x+1)(x-1) \).
          Untuk $\lambda_1=-1$,
          \begin{equation*}
            \nullspace{t+1}=\set{\colvec{-y \\ y}
                                    \suchthat y\in\C } 
          \end{equation*}
          dan ruang nol $(t+1)^2$ sama.
          Untuk $\lambda_2=1$,
          \begin{equation*}
            \nullspace{t-1}=\set{\colvec{y \\ y}
                                    \suchthat y\in\C } 
          \end{equation*}
          dan ruang nol $(t-1)^2$ sama.
          Kita dapat mengambil basis berikut
          \begin{equation*}
            B=\sequence{\colvec[r]{1 \\ -1},\colvec[r]{1 \\ 1}}
          \end{equation*}
          untuk memperoleh suatu diagonalisasi.
          \begin{equation*}
            \begin{mat}[r]
              1  &1  \\
             -1  &1
            \end{mat}^{-1}
            \begin{mat}[r]
              0  &1  \\
              1  &0
            \end{mat}
            \begin{mat}[r]
              1   &1  \\
              -1  &1
            \end{mat}
            =
            \begin{mat}[r]
              -1  &0  \\
              0   &1
            \end{mat}
          \end{equation*}
      \end{exparts}  
     \end{answer}
  \recommended \item 
    Carilah matriks Jordan yang merepresentasikan operator diferensiasi
    pada \( \polyspace_3 \).
    \begin{answer}
      Transformasi $\map{d/dx}{\polyspace_3}{\polyspace_3}$
      bersifat nilpoten.
      Tindakannya pada basis \( B=\sequence{x^3,3x^2,6x,6} \)
      adalah $x^3\mapsto 3x^2\mapsto 6x\mapsto 6\mapsto 0$.
      Bentuk Jordannya sama dengan bentuk kanoniknya sebagai matriks nilpoten.
      \begin{equation*}
         J=
         \begin{mat}[r]
           0  &0  &0  &0  \\
           1  &0  &0  &0  \\
           0  &1  &0  &0  \\
           0  &0  &1  &0
         \end{mat}
      \end{equation*}
    \end{answer}
   \recommended \item 
      Tentukan apakah kedua matriks berikut serupa.
      \begin{equation*}
         \begin{mat}[r]
            1  &-1 \\
            4  &-3 \\
         \end{mat}
         \qquad
         \begin{mat}[r]
           -1  &0  \\
            1  &-1 \\
         \end{mat}
      \end{equation*}
      \begin{answer}
        Ya.
        Masing-masing memiliki polinomial karakteristik $(x+1)^2$.
        Perhitungan pangkat-pangkat $T_1+1\cdot I$ dan
        $T_2+1\cdot I$ memberikan kedua ruang berikut.
        \begin{equation*}
          \nullspace{t_1+1}=\set{\colvec{y/2 \\ y} \suchthat y\in\C}
          \qquad
          \nullspace{t_2+1}=\set{\colvec{0 \\ y} \suchthat y\in\C}
        \end{equation*}
        (Tentu saja, untuk masing-masing, ruang nol kuadratnya
        adalah seluruh ruang.)
        Pola pertumbuhan nulitas menunjukkan bahwa masing-masing
        serupa dengan matriks berbentuk Jordan berikut,
        \begin{equation*}
           \begin{mat}[r]
             -1  &0  \\
              1  &-1 \\
           \end{mat}
        \end{equation*}
        sehingga keduanya serupa satu sama lain.
      \end{answer}
  \item 
     Carilah bentuk Jordan matriks berikut.
     \begin{equation*}
        \begin{mat}[r]
           0  &-1  \\
           1  &0
        \end{mat}
     \end{equation*}
     Berikan pula suatu basis Jordan.
     \begin{answer}
       Polinomial karakteristiknya adalah
       \( c(x)=x^2+1 \), yang memiliki akar kompleks
       \( x^2+1=(x+i)(x-i) \).
       Karena akar-akarnya berbeda,
       matriks ini dapat didiagonalkan dan bentuk Jordannya adalah
       matriks diagonal berikut.
       \begin{equation*}
         \begin{mat}[r]
           -i  &0  \\
            0  &i
         \end{mat}
       \end{equation*}  
       Untuk mencari basis terkait, kita menghitung ruang nol.
       \begin{equation*}
         \nullspace{t+i}=\set{\colvec{-iy \\ y}
                                        \suchthat y\in\C} 
         \qquad
         \nullspace{t-i}=\set{\colvec{iy \\ y}
                                        \suchthat y\in\C} 
       \end{equation*}
       Sebagai contoh,
       \begin{equation*}
         T+i\cdot I=
         \begin{mat}[r]
           i  &-1  \\
           1  &i
         \end{mat}
       \end{equation*}
       sehingga kita memperoleh deskripsi ruang nol $t+i$ dengan menyelesaikan
       sistem linear berikut.
       \begin{equation*}
         \begin{linsys}{2}
           ix  &-  &y  &=  &0  \\
            x  &+  &iy &=  &0
         \end{linsys}
         \grstep{i\rho_1+\rho_2}
         \begin{linsys}{2}
           ix  &-  &y  &=  &0  \\
               &   &0  &=  &0
         \end{linsys}
       \end{equation*}
       (Untuk mengubah hubungan $ix=y$ agar variabel utama $x$
       dinyatakan dalam variabel bebas $y$, kita dapat mengalikan kedua
       ruas dengan $-i$.)

       Hasilnya, salah satu basis semacam itu adalah berikut.
       \begin{equation*}
         B=\sequence{\colvec[r]{-i \\ 1},
                     \colvec[r]{i \\ 1}}
       \end{equation*}  
     \end{answer}
  \item 
    Berapa banyak kelas keserupaan untuk matriks \( \nbyn{3} \)
    yang nilai eigennya hanya \( -3 \) dan \( 4 \)?
    \begin{answer}
     Kita dapat menghitung kelas-kelas yang mungkin dengan menghitung
     wakil kanonik yang mungkin, yakni matriks berbentuk Jordan yang mungkin.
     Polinomial karakteristiknya harus berupa $c_1(x)=(x+3)^2(x-4)$
     atau $c_2(x)=(x+3)(x-4)^2$.
     Dalam kasus $c_1$, terdapat dua kemungkinan tindakan
     $t+3$ pada basis rantai untuk $\gennullspace{t+3}$.
     \begin{equation*}
       \vec{\beta}_1\mapsto\vec{\beta}_2\mapsto \zero
       \qquad
       \begin{array}[t]{l}
         \vec{\beta}_1\mapsto\zero \\
         \vec{\beta}_2\mapsto\zero 
       \end{array}
     \end{equation*}
     Terdapat dua matriks berbentuk Jordan yang terkait.
     \begin{equation*}
        \begin{mat}[r]
          -3  &0  &0  \\
           1  &-3 &0  \\
           0  &0  &4
        \end{mat}
        \qquad
        \begin{mat}[r]
          -3  &0  &0  \\
           0  &-3 &0  \\
           0  &0  &4
        \end{mat}
      \end{equation*}
      Demikian pula, terdapat dua matriks berbentuk Jordan yang mungkin
      timbul dari $c_2$.
      \begin{equation*}
        \begin{mat}[r]
          -3  &0  &0  \\
           0  &4  &0  \\
           0  &1  &4
        \end{mat}
        \qquad
        \begin{mat}[r]
          -3  &0  &0  \\
           0  &4  &0  \\
           0  &0  &4
        \end{mat}
     \end{equation*}  
     Jadi, seluruhnya terdapat empat bentuk Jordan yang mungkin.
    \end{answer}
  \recommended \item
    Buktikan bahwa suatu matriks dapat didiagonalkan jika dan hanya jika
    polinomial minimalnya hanya memiliki faktor linear tanpa pengulangan.
    \begin{answer}
       Bentuk Jordan bersifat unik.
       Suatu matriks diagonal sudah berbentuk Jordan.
       Jadi, bentuk Jordan matriks yang dapat didiagonalkan adalah diagonalisasinya.
       Jika polinomial minimal memiliki faktor berpangkat lebih tinggi daripada satu,
       maka bentuk Jordannya memiliki entri subdiagonal \( 1 \), sehingga tidak
       diagonal.
     \end{answer}
  \item 
    Berikan contoh transformasi linear pada ruang vektor
    yang tidak memiliki subruang invarian nontrivial.
    \begin{answer}
      Salah satu contohnya adalah transformasi pada \( \C \) yang
       mengirim \( x \) ke \( -x \).
     \end{answer}
  \item 
    Tunjukkan bahwa suatu subruang invarian terhadap \( t-\lambda_1 \)
    jika dan hanya jika invarian terhadap \( t-\lambda_2 \).
    \begin{answer}
      Terapkan \nearbylemma{le:tInvIfftMinLambdaInv} dua kali;
      subruang tersebut invarian terhadap $t-\lambda_1$ jika dan hanya jika
      invarian terhadap $t$, yang selanjutnya berlaku jika dan hanya jika
      subruang itu invarian terhadap $t-\lambda_2$.
    \end{answer}
  \item 
     Buktikan atau sangkal: dua matriks \( \nbyn{n} \)
     serupa jika dan hanya jika memiliki polinomial karakteristik dan
     polinomial minimal yang sama.
     \begin{answer}
       Salah; kedua matriks $\nbyn{4}$ berikut masing-masing memiliki $c(x)=(x-3)^4$
       dan $m(x)=(x-3)^2$.
       \begin{equation*}
          \begin{mat}[r]
             3  &0  &0  &0  \\
             1  &3  &0  &0  \\
             0  &0  &3  &0  \\
             0  &0  &1  &3
          \end{mat}
          \quad
          \begin{mat}[r]
             3  &0  &0  &0  \\
             1  &3  &0  &0  \\
             0  &0  &3  &0  \\
             0  &0  &0  &3
          \end{mat}
       \end{equation*} 
     \end{answer}
  \item 
    \definend{jejak}\index{jejak}\index{matriks!jejak}
    suatu matriks persegi adalah jumlah entri diagonalnya.
    \begin{exparts}
       \partsitem Carilah rumus polinomial karakteristik suatu matriks
         $\nbyn{2}$.
       \partsitem Tunjukkan bahwa
         jejak invarian terhadap keserupaan, sehingga kita dapat secara bermakna
         berbicara tentang `jejak suatu pemetaan'.\index{pemetaan linear!jejak}
         (\textit{Petunjuk:} lihat butir sebelumnya.)
       \partsitem Apakah jejak invarian terhadap ekuivalensi matriks?
       \partsitem Tunjukkan bahwa jejak suatu pemetaan adalah jumlah nilai eigennya
         (dengan memperhitungkan multiplisitas).
       \partsitem Tunjukkan bahwa jejak suatu pemetaan nilpoten adalah nol.
         Apakah konversnya berlaku?
    \end{exparts}
    \begin{answer}
      \begin{exparts}
         \partsitem Berikut polinomial karakteristiknya.
           \begin{equation*}
             \begin{vmat}
               a-x  &b  \\
               c  &d-x
             \end{vmat}
             =(a-x)(d-x)-bc=ad-(a+d)x+x^2-bc
             =x^2-(a+d)x+(ad-bc)
           \end{equation*}
           Perhatikan bahwa determinan muncul sebagai suku konstan.
         \partsitem Ingatlah bahwa polinomial karakteristik
            \( \deter{T-xI} \) invarian terhadap keserupaan.
            Gunakan rumus ekspansi permutasi untuk menunjukkan bahwa setelah
            polinomial itu dinormalkan menjadi monik, jejak adalah negatif
            dari koefisien \( x^{n-1} \).
         \partsitem Tidak. Terdapat matriks $T$ dan $S$ yang
            ekuivalen, $S=PTQ$ (untuk suatu $P$ dan $Q$ nonsingular),
            tetapi memiliki jejak berbeda.
            Berikut contoh yang mudah.
            \begin{equation*}
               PTQ=
               \begin{mat}[r]
                  2  &0  \\
                  0  &1
               \end{mat}
               \begin{mat}[r]
                  1  &0  \\
                  0  &1
               \end{mat}
               \begin{mat}[r]
                  1  &0  \\
                  0  &1
               \end{mat}
               =
               \begin{mat}[r]
                  2  &0  \\
                  0  &1
               \end{mat}
            \end{equation*}
            Contoh yang lebih mudah lagi dapat dibuat dengan matriks $\nbyn{1}$.
         \partsitem Ubah matriks ke bentuk Jordan.
            Berdasarkan butir kedua, jejak tidak berubah.
         \partsitem Bagian pertama mudah; gunakan butir keempat.
            Konversnya tidak berlaku:~matriks berikut
            \begin{equation*}
               \begin{mat}[r]
                  1  &0  \\
                  0  &-1
               \end{mat}
            \end{equation*}
            memiliki jejak nol, tetapi tidak nilpoten.
       \end{exparts}  
     \end{answer}
  \item 
    Untuk menggunakan \nearbydefinition{def:invariant} guna memeriksa apakah suatu subruang
    invarian terhadap $t$, tampaknya kita harus memeriksa tak terhingga banyaknya
    vektor dalam subruang (nontrivial) untuk melihat apakah semuanya memenuhi syarat.
    Buktikan bahwa suatu subruang invarian terhadap \( t \) jika dan hanya jika basisnya
    memiliki sifat bahwa untuk setiap elemennya, $t(\vec{\beta})$ berada dalam
    subruang tersebut.
    \begin{answer}
      Misalkan \( B_M \) adalah basis untuk suatu subruang \( M \) dari sebuah
      ruang vektor.
      Salah satu implikasinya jelas; jika \( M \) invarian terhadap \( t \), maka
      secara khusus, jika \( \vec{m}\in B_M \), berlaku \( t(\vec{m})\in M \).
      Untuk implikasi lainnya, misalkan
      \( B_M=\sequence{\vec{\beta}_1,\dots,\vec{\beta}_q} \), dan perhatikan bahwa
      \( t(\vec{m})=t(m_1\vec{\beta}_1+\dots+m_q\vec{\beta}_q)
                   =m_1t(\vec{\beta}_1)+\dots+m_qt(\vec{\beta}_q) \)
      berada dalam \( M \) karena setiap subruang tertutup terhadap kombinasi
      linear.
    \end{answer}
  \recommended \item 
    Apakah keinvarianan terhadap \( t \) terpelihara dalam irisan?
    Dalam gabungan?
    Komplemen?
    Jumlah subruang?
    \begin{answer}
      Ya, irisan subruang-subruang yang invarian terhadap $t$
      juga invarian terhadap $t$.
      Andaikan \( M \) dan \( N \) invarian terhadap \( t \).
      Jika \( \vec{v}\in M\intersection N \), maka \( t(\vec{v})\in M \)
      berdasarkan keinvarianan \( M \), dan \( t(\vec{v})\in N \) berdasarkan
      keinvarianan \( N \).

      Tentu saja, gabungan dua subruang belum tentu merupakan subruang
      (ingat bahwa sumbu-$x$\hbox{} dan sumbu-$y$ merupakan subruang bidang
      $\Re^2$, tetapi gabungan kedua sumbu tidak tertutup terhadap
      penjumlahan vektor; misalnya, gabungan itu tidak memuat
      $\vec{e}_1+\vec{e}_2$.)
      Namun, gabungan himpunan-himpunan invarian merupakan himpunan invarian; jika
      \( \vec{v}\in M\union N \), maka \( \vec{v}\in M \) atau \( \vec{v}\in N \),
      sehingga \( t(\vec{v})\in M \) atau \( t(\vec{v})\in N \).

      Tidak, komplemen suatu subruang invarian belum tentu invarian.
      Perhatikan subruang
      \begin{equation*}
        \set{\colvec{x \\ 0}\suchthat x\in\C}
      \end{equation*}
      dari \( \C^2 \) di bawah transformasi nol.

      Ya, jumlah dua subruang invarian juga invarian.
      Pemeriksaannya mudah.
    \end{answer}
  \item  \label{exer:OrderC}
     Berikan cara mengurutkan blok-blok Jordan jika sebagian nilai eigennya
     merupakan bilangan kompleks.
     Artinya, sarankan pengurutan yang wajar untuk bilangan kompleks.
     \begin{answer}
       Salah satu pengurutan semacam itu adalah \definend{urutan kamus}.
       Urutkan terlebih dahulu berdasarkan komponen real, kemudian berdasarkan koefisien \( i \).
       Sebagai contoh, \( 3+2i<4+1i \), tetapi \( 4+1i<4+2i \).
     \end{answer}
  \item
    Misalkan \( \polyspace_j(\Re) \) adalah ruang vektor atas bilangan
    real yang terdiri atas polinomial berderajat paling tinggi \( j \).
    Tunjukkan bahwa jika \( j\le k \), maka \( \polyspace_j(\Re) \) merupakan
    subruang invarian dari \( \polyspace_k(\Re) \) di bawah operator diferensiasi.
    Dalam \( \polyspace_7(\Re) \), adakah salah satu dari \( \polyspace_0(\Re) \),
    \ldots, \( \polyspace_6(\Re) \) yang memiliki komplemen invarian?
    \begin{answer}
     Bagian pertama mudah\Dash turunan sebarang polinomial real adalah
      polinomial real berderajat lebih rendah.
      Jawaban untuk bagian kedua adalah `tidak'; setiap komplemen dari
      \( \polyspace_j(\Re) \) harus memuat polinomial berderajat \( j+1 \),
      dan turunan polinomial tersebut berada dalam \( \polyspace_j(\Re)\).
     \end{answer}
  \item
    Dalam \( \polyspace_n(\Re) \), yaitu ruang vektor (atas bilangan
    real) dari polinomial berderajat paling tinggi \( n \),
    \begin{equation*}
      \mathcal{E}=
      \set{p(x)\in\polyspace_n(\Re)\suchthat p(-x)=p(x) \text{\ untuk semua\ }x}
    \end{equation*}
    dan
    \begin{equation*}
      \mathcal{O}=
      \set{p(x)\in\polyspace_n(\Re)\suchthat p(-x)=-p(x) \text{\ untuk semua\ }x}
    \end{equation*}
    masing-masing merupakan polinomial \definend{genap}\index{polinomial genap}\index{polinomial!genap}
    dan \definend{ganjil}\index{polinomial ganjil}\index{polinomial!ganjil};
    \( p(x)=x^2 \) genap, sedangkan \( p(x)=x^3 \) ganjil.
    Tunjukkan bahwa keduanya merupakan subruang.
    Apakah keduanya saling melengkapi?
    Apakah keduanya invarian di bawah transformasi diferensiasi?
    \begin{answer}
      Untuk bagian pertama, tunjukkan bahwa masing-masing merupakan subruang,
      lalu amati bahwa setiap polinomial dapat dituliskan secara unik
       sebagai jumlah suku berpangkat genap dan suku berpangkat ganjil (polinomial
       nol termasuk keduanya).
       Jawaban untuk bagian kedua adalah `tidak': \( x^2 \) genap, sedangkan
       \( 2x \) ganjil.
     \end{answer}
  % Lemma gone when moved to new JF presentation
  % \item 
  %   \nearbylemma{le:InvCompSubspSplitTrans} says that if \( M \) and
  %   \( N \) are
  %   invariant complements then \( t \) has a representation in the given
  %   block form (with respect to the same ending as starting basis, of course).
  %   Does the implication reverse?
  %   \begin{answer}
  %     Yes.
  %     If \( \rep{t}{B,B} \) has the given block form, take \( B_M \) to
  %     be the first \( j \) vectors of \( B \), where \( J \) is the
  %     \( \nbyn{j} \) upper left submatrix.
  %     Take \( B_N \) to be the remaining \( k \) vectors in \( B \).
  %     Let \( M \) and \( N \) be the spans of \( B_M \) and \( B_N \).
  %     Clearly \( M \) and \( N \) are complementary.
  %     To see \( M \) is invariant (\( N \) works the same way), represent
  %     any \( \vec{m}\in M \) with respect to \( B \), note the last
  %     \( k \) components are zeroes, and multiply by the given block
  %     matrix.
  %     The final \( k \) components of the result are zeroes, so that
  %     result is again in \( M \). 
  %    \end{answer}
   \item 
     Matriks \( S \) merupakan \definend{akar kuadrat}\index{akar kuadrat}
     dari matriks lain \( T \) jika \( S^2=T \).
     Tunjukkan bahwa setiap matriks nonsingular memiliki akar kuadrat.
     \begin{answer}
         Ubah matriks ke bentuk Jordan.
         Karena nonsingular, tidak ada nilai eigen nol pada diagonal.
         Tiru contoh berikut:
          \begin{equation*}
             \begin{mat}[r]
                9  &0  &0 \\
                1  &9  &0 \\
                0  &0  &4
             \end{mat}
             =
             \begin{mat}[r]
                3  &0  &0 \\
               1/6 &3  &0 \\
                0  &0  &2
             \end{mat}^2
          \end{equation*}
         untuk membangun akar kuadrat.
         Tunjukkan bahwa hal ini terpelihara di bawah keserupaan: jika \( S^2=T \), maka
         \( (PSP^{-1})(PSP^{-1})=PTP^{-1} \). 
    \end{answer}
\index{bentuk Jordan|)}
\end{exercises}
